\documentclass[11pt, ukrainian]{article}
\setlength\textwidth{180mm} \setlength\textheight{277mm}
\setlength\voffset{-38mm} \setlength\hoffset{-28mm}
%\setlength\parindent{0mm}
\pagestyle{empty}
\usepackage[T2A]{fontenc}
\usepackage[utf8]{inputenc}
\usepackage[ukrainian]{babel}
\begin{document}
{\Large \center  Варіанти завдань до лабораторної роботи №4\par 2020/21 н.р.\par \bigskip\bigskip}

Порядок полів у записах та інформацію додаткового файлу фіксовано у варіанті.

Для порівняння числових значень використовується традиційний порядок.

Для порівняння рядкових значень використовується лексикографічний порядок, якщо в умові варіанту явно не вказано інше. 

У випадку сортування за кількома полями результуюче сортування будується 
за порядками окремих полів за допомогою конструкції лексикографічного порядку. 
Послідовність полів, за якими має відбуватися сортування, задається в умові. 
Якщо не вказано інше, то поле сортується за зростанням.

Якщо умова вимагає знайти $k$ най… і $k$-е місце займають кілька, то слід виводити всіх їх.

\newpage
\par \bigskip

%begin
{\center Варіант  \No \textbf { 1 }\par}
\bigskip
%type = meteo
%variant = 114
Обробляються дані спостережень за погодою. 

\underline{Записи основного файлу містять поля}: день, місяць, код метеостанції, максимальна денна температура, рік, мінімальна денна температура, середня денна температура, кількість опадів, сила вітру, відносна вологість.

\underline{У допоміжному файлі наявні ключі}: середнє значення дня, округлене до 1 знаку; загальна кількість записів.

\underline{Знайти} метеостанції, по яких жодний з показників не відхилявся від середніх по всіх спостереженнях більше ніж на 50\%. Вивести по кожному з них інформацію:\par
{\noindent}--- на першому рядку:\par
метеостанція, середні показники по метеостанції: кількість опадів, сила вітру, середня денна температура;

{\noindent}--- на наступних рядках, починаючи з табуляції, вивести для них дані по місяцях (по одному на рядок):\par
рік, місяць, середня середньоденна температура, середня вологість

{\noindent}у такому сортуванні: рік, місяць.

%Метеостанції сортувати: кількість спостережень (за спаданням), середня сила вітру, метеостанція.

\bigskip\textit{Обмеження}

Код метеостанції є рядком довжини від 4 до 9 включно.
Може містити літери алфавіту, десяткові цифри, підкреслення. Решта полів числові.

Рік є чотиризначним цілим.
Кількість опадів є дійсним значенням, подається у форматі з 2 знаками після крапки.
Сила вітру є дійсним значенням, подається у форматі з 1 знаком після крапки.
Температури є дійсними значеннями, подаються у форматі з 3 знаками після крапки.
Цілі значення подаються без десяткової крапки.
Для дійсних використовується зображення з фіксованою крапкою.

\textit{Вихідне форматування дійсних}

Усі дійсні значення виводяться у форматі з фіксованою крапкою з тією ж кількістю знаків після крапки, що і у вхідному зображенні.

Якщо у варіанті раніше не було зазначено, то \underline{обчислювані програмою} середні значення виводяться у форматі з 2 знаками після крапки.
%\bigskip%\textit{Для деяких варіантів може знадобитись}




\par
%end
\newpage
%begin
{\center Варіант  \No \textbf { 2 }\par}
\bigskip
%type = meteo
%variant = 111
Обробляються дані спостережень за погодою. 

\underline{Записи основного файлу містять поля}: середня денна температура, кількість опадів, відносна вологість, сила вітру, мінімальна денна температура, максимальна денна температура, код метеостанції, місяць, день, рік.

\underline{У допоміжному файлі наявні ключі}: перший рік, в якому були спостереження; кількість спостережень, в яких середня денна температура була більша 10.

\underline{Знайти} 3 метеостанції, на яких була зафіксована найбільша кількість опадів за день. Вивести по кожному з них інформацію:\par
{\noindent}--- на першому рядку:\par
кількість спостережень, середня кількість опадів, найбільша кількість опадів за день, метеостанція;

{\noindent}--- на наступних рядках, починаючи з табуляції, вивести для них дані по місяцях (по одному на рядок):\par
рік, місяць, середня сила вітру за місяць, кількість опадів за місяць

{\noindent}у такому сортуванні: місяць, рік.

%Метеостанції сортувати: кількість спостережень (за спаданням), метеостанція.

\bigskip\textit{Обмеження}

Код метеостанції є рядком довжини від 5 до 8 включно.
Може містити літери алфавіту, десяткові цифри, підкреслення. Решта полів числові.

Рік є чотиризначним цілим.
Кількість опадів є дійсним значенням, подається у форматі з 1 знаком після крапки.
Сила вітру є цілим значенням.
Температури є дійсними значеннями, подаються у форматі з 2 знаками після крапки.
Цілі значення подаються без десяткової крапки.
Для дійсних використовується зображення з фіксованою крапкою.

\textit{Вихідне форматування дійсних}

Усі дійсні значення виводяться у форматі з фіксованою крапкою з тією ж кількістю знаків після крапки, що і у вхідному зображенні.

Якщо у варіанті раніше не було зазначено, то \underline{обчислювані програмою} середні значення виводяться у форматі з 1 знаком після крапки.
%\bigskip%\textit{Для деяких варіантів може знадобитись}




\par
%end
\newpage
%begin
{\center Варіант  \No \textbf { 3 }\par}
\bigskip
%type = absense
%variant = -1
Обробляються дані семестрового журналу перевірок відвідування занять студентами деканатом (фіксуються відсутні). 

\underline{Записи основного файлу містять поля}: курс, прізвище, ім’я, по-батькові, день тижня, пара, вид занять, аудиторія, код групи, навчальний тиждень, предмет.

\underline{У допоміжному файлі наявні ключі}: сума днів тижня, сума курсів.

\underline{Знайти} top-10 прогульщиків. Вивести по кожному з них інформацію:\par
{\noindent}--- на першому рядку:\par
прізвище, ім’я, по-батькові, кількість пропущених пар, кількість пропущених лекцій, кількість пропущених лабораторних, кількість пропущених практичних;

{\noindent}--- на наступних рядках, починаючи з табуляції, вивести пропуски студента (по одному на рядок):\par
предмет, вид занять, тиждень, день, пара, аудиторія

{\noindent}у такому сортуванні: тиждень, день, пара, предмет, вид занять.

%Студентів сортувати: кількість пропусків (за спаданням), прізвище, ім’я, по-батькові.

\bigskip\textit{Обмеження}

Навчальний тиждень може приймати цілі значення від 1 до 14.
День тижня є цілим від 1 до 5. Пара є цілим від 1 до 4.
Курс є цілим від 1 до 4. Аудиторія є додатнім цілим.
Вид занять може приймати значення:
L --- лекція, Practice --- практичне, Сем. --- семінар, Lab --- лабораторне.

Решта полів є непорожніми рядками.
Рядки починатися та закінчуватися білими символами не можуть.

Групи кодуються в межах курсу.
Код групи складається не більш ніж з 5 символів.
Може містити літери алфавіту, десяткові цифри та дефіс.

Предмет є рядком, може мати від 3 до 21 символів включно.
Може містити літери алфавіту, десяткові цифри, апостроф, дефіс, пробіл.
Предмети різних курсів вважати різними навіть за однакових назв.

Прізвище, ім’я та по-батькові (за наявності у варіанті)
можуть мати від 3 до 26, 20, 25 відповідно символів включно кожне.
Можуть містити літери алфавіту, апостроф, пробіл та дефіс.
Студенти однозначно ідентифікуються прізвищем, ім’ям та по-батькові (за наявності у варіанті).
Жоден студент не вчиться одночасно на різних курсах чи в різних групах.

Види занять сортуються так: лабораторне < практичне < лекція < семінар.

%\bigskip%\textit{Для деяких варіантів може знадобитись}




\par
%end
\newpage
%begin
{\center Варіант  \No \textbf { 4 }\par}
\bigskip
%type = session
%variant = 74
Обробляються результати складання студентами екзаменів зимової сесії.

\underline{Записи основного файлу містять поля}: оцінка за державною шкалою, набрані на екзамені бали, номер залікової книжки, прізвище, набрані впродовж семестру бали, сумарна оцінка в балах за 100-бальною системою, ім’я, предмет, код групи.

\underline{У допоміжному файлі наявні ключі}: найменша довжина назви предмета, загальна кількість записів.

\underline{Знайти} предмети, середній бал за які, вищий середнього за всю сесію. Вивести по кожному з них інформацію:\par
{\noindent}--- на першому рядку:\par
предмет, середній бал за предметом (округлений до 1 знаку), кількість студентів, що склали на «відмінно»;

{\noindent}--- на наступних рядках, починаючи з табуляції, вивести для предмета студентів, що його не склали (по одному на рядок):\par
прізвище, ім’я, по-батькові, номер залікової книжки, група, оцінка за державною шкалою прописом

{\noindent}у такому сортуванні: номер залікової книжки.

%Предмети сортувати: середній бал (округлений), кількість студентів, що склали на «відмінно», предмет.

\bigskip\textit{Обмеження}

Студент не може складати предмет більше одного разу.

Бали та оцінка є значеннями цілих числових типів. Решта полів є непорожніми рядками.
Рядки розпочинатись та закінчуватись білими символами не можуть.

Оцінка за державною шкалою може приймати значення: 2-5~--- як зазвичай, 0~--- не з’явився,$-1$~--- не допущений. 
Оцінки 3-5 вважаються задовільними, решта~--- незадовільними. 
Бали є цілими та невід’ємними.
Набрані впродовж семестру бали не можуть перевищувати 60. 
Набрані на екзамені бали не можуть перевищувати 40. 
Набрані на екзамені бали повинні бути не менше 24 або дорівнювати 0(в останньому випадку оцінка не може бути задовільною). 
Сумарна оцінка в балах є сумою балів, набраних у семестрі та на екзамені.
Сумарна оцінка в балах та за державною шкалою мають бути узгоджені між собою.

Прізвище, ім’я та по-батькові (якщо наявне у варіанті) можуть мати до 25, 26, 30 відповідно символів включно кожне. 
Крім літер алфавіту можуть містити апостроф, дефіс та пробіл.

Предмет може мати від 4 до 58 символів включно. 
Крім літер алфавіту може містити подвійні лапки, пробіл та дефіс.

Код групи є непорожнім та складається не більш ніж з 4 літер. 
Крім літер алфавіту може містити десяткові цифри та дефіс.


Номер залікової книжки однозначно ідентифікує студента та складається з 6 десяткових цифр. 

\bigskip\textit{Для деяких варіантів може знадобитись}

Середній бал/оцінка за предметом визначається як середнє арифметичнесумарних балів/державних оцінок за цим предметом. 
Середній бал/оцінка студента визначається як середнє арифметичне сумарних балів/державних оцінокцього студента. 
За обчислення середньої оцінки недопущені та ті, що не з’явилися, рахуються як 2.
У решті випадків недопуски підсумовуються як $-1$, а «не з’явився» як 0.
Якщо студент не гірше ніж задовільно склав усі предмети, 
то його рейтинг за балами/оцінкою визначається як середній бал/оцінка; інакше дорівнює 0.

Стабільність студента визначається як модуль різниці між його кращим та найгіршим сумарним балом.
Стабільність предмета визначається як модуль різниці між найкращим та найгіршимсумарним балом за цим предметом.






\par
%end
\newpage
%begin
{\center Варіант  \No \textbf { 5 }\par}
\bigskip
%type = meteo
%variant = 101
Обробляються дані спостережень за погодою. 

\underline{Записи основного файлу містять поля}: мінімальна денна температура, максимальна денна температура, день, місяць, сила вітру, код метеостанції, кількість опадів, середня денна температура, рік, відносна вологість.

\underline{У допоміжному файлі наявні ключі}: найменше значення відносної вологості, Failed successfullY.

\underline{Знайти} дати, по яких середня максимальна денна температура була більша ніж в середньому по всіх спостереженнях. Вивести по кожному з них інформацію:\par
{\noindent}--- на першому рядку:\par
рік, місяць, день, кількість опадів за цей рік, середня середньоденна температура (беремо середні по датах і потім їх середнє);

{\noindent}--- на наступних рядках, починаючи з табуляції, вивести для дати спостереження метеостанції, по яких різниця максимальної та мінімальної температур менша 5.13 (по одному на рядок):\par
середня денна температура, максимальна денна температура, мінімальна денна температура, метеостанція

{\noindent}у такому сортуванні: різниця максимальної та мінімальної температур, метеостанція.

%Дати сортувати: рік, місяць, день.

\bigskip\textit{Обмеження}

Код метеостанції є рядком довжини від 3 до 11 включно.
Може містити літери алфавіту, десяткові цифри, підкреслення. Решта полів числові.

Рік є чотиризначним цілим.
Кількість опадів є цілим значенням.
Сила вітру є дійсним значенням, подається у форматі з 2 знаками після крапки.
Температури є дійсними значеннями, подаються у форматі з 1 знаком після крапки.
Цілі значення подаються без десяткової крапки.
Для дійсних використовується зображення з фіксованою крапкою.

\textit{Вихідне форматування дійсних}

Усі дійсні значення виводяться у форматі з фіксованою крапкою з тією ж кількістю знаків після крапки, що і у вхідному зображенні.

Якщо у варіанті раніше не було зазначено, то \underline{обчислювані програмою} середні значення виводяться у форматі з 3 знаками після крапки.
%\bigskip%\textit{Для деяких варіантів може знадобитись}




\par
%end
\newpage
%begin
{\center Варіант  \No \textbf { 6 }\par}
\bigskip
%type = session
%variant = 53
Обробляються результати складання студентами екзаменів зимової сесії.

\underline{Записи основного файлу містять поля}: код групи, набрані на екзамені бали, набрані впродовж семестру бали, сумарна оцінка в балах за 100-бальною системою, предмет, ім’я, оцінка за державною шкалою, прізвище, номер залікової книжки.

\underline{У допоміжному файлі наявні ключі}: найбільша довжина предмету, загальна кількість записів.

\underline{Знайти} 7 найкращих студентів (за рейтингом за національною шкалою). Вивести по кожному з них інформацію:\par
{\noindent}--- на першому рядку:\par
прізвище, ім’я, номер залікової книжки, рейтинг за балами (округлений до цілого), рейтинг за оцінкою (округлений до 2 знаків);

{\noindent}--- на наступних рядках, починаючи з табуляції, вивести для студента всі предмети, де було отримано не «відмінно» (по одному на рядок):\par
сумарна оцінка, бали в семестрі, державна оцінка, предмет

{\noindent}у такому сортуванні: предмет.

%Студентів сортувати: прізвище, ім’я, сума балів за сесію, номер залікової книжки.

\bigskip\textit{Обмеження}

Студент не може складати предмет більше одного разу.

Бали та оцінка є значеннями цілих числових типів. Решта полів є непорожніми рядками.
Рядки розпочинатись та закінчуватись білими символами не можуть.

Оцінка за державною шкалою може приймати значення: 2-5~--- як зазвичай, 0~--- не з’явився,$-1$~--- не допущений. 
Оцінки 3-5 вважаються задовільними, решта~--- незадовільними. 
Бали є цілими та невід’ємними.
Набрані впродовж семестру бали не можуть перевищувати 60. 
Набрані на екзамені бали не можуть перевищувати 40. 
Набрані на екзамені бали повинні бути не менше 24 або дорівнювати 0(в останньому випадку оцінка не може бути задовільною). 
Сумарна оцінка в балах є сумою балів, набраних у семестрі та на екзамені.
Сумарна оцінка в балах та за державною шкалою мають бути узгоджені між собою.

Прізвище, ім’я та по-батькові (якщо наявне у варіанті) можуть мати до 20, 22, 21 відповідно символів включно кожне. 
Крім літер алфавіту можуть містити апостроф, дефіс та пробіл.

Предмет може мати від 3 до 67 символів включно. 
Крім літер алфавіту може містити подвійні лапки, пробіл та дефіс.

Код групи є непорожнім та складається не більш ніж з 6 літер. 
Крім літер алфавіту може містити десяткові цифри та дефіс.


Номер залікової книжки однозначно ідентифікує студента та складається з 7 десяткових цифр. 

\bigskip\textit{Для деяких варіантів може знадобитись}

Середній бал/оцінка за предметом визначається як середнє арифметичнесумарних балів/державних оцінок за цим предметом. 
Середній бал/оцінка студента визначається як середнє арифметичне сумарних балів/державних оцінокцього студента. 
За обчислення середньої оцінки недопущені та ті, що не з’явилися, рахуються як 2.
У решті випадків недопуски підсумовуються як $-1$, а «не з’явився» як 0.
Якщо студент не гірше ніж задовільно склав усі предмети, 
то його рейтинг за балами/оцінкою визначається як середній бал/оцінка; інакше дорівнює 0.

Стабільність студента визначається як модуль різниці між його кращим та найгіршим сумарним балом.
Стабільність предмета визначається як модуль різниці між найкращим та найгіршимсумарним балом за цим предметом.






\par
%end
\newpage
%begin
{\center Варіант  \No \textbf { 7 }\par}
\bigskip
%type = session
%variant = 70
Обробляються результати складання студентами екзаменів зимової сесії.

\underline{Записи основного файлу містять поля}: прізвище, код групи, ім’я, набрані на екзамені бали, предмет, сумарна оцінка в балах за 100-бальною системою, номер залікової книжки, оцінка за державною шкалою, набрані впродовж семестру бали.

\underline{У допоміжному файлі наявні ключі}: загальна кількість предметів, загальна кількість записів.

\underline{Знайти} предмети, за якими кількість незадовільних оцінок найменша. Вивести по кожному з них інформацію:\par
{\noindent}--- на першому рядку:\par
предмет, середній бал (округлений до цілого), кількість студентів, що складали;

{\noindent}--- на наступних рядках, починаючи з табуляції, вивести для предмета всіх студентів, що його складали (по одному на рядок):\par
прізвище, ім’я, по-батькові, номер залікової книжки, група, оцінка в балах, оцінка за державною шкалою, бали за екзамен, бали за семестр

{\noindent}у такому сортуванні: оцінка в балах (за спаданням), номер залікової книжки.

%Предмети сортувати: середній бал (округлений), предмет.

\bigskip\textit{Обмеження}

Студент не може складати предмет більше одного разу.

Бали та оцінка є значеннями цілих числових типів. Решта полів є непорожніми рядками.
Рядки розпочинатись та закінчуватись білими символами не можуть.

Оцінка за державною шкалою може приймати значення: 2-5~--- як зазвичай, 0~--- не з’явився,$-1$~--- не допущений. 
Оцінки 3-5 вважаються задовільними, решта~--- незадовільними. 
Бали є цілими та невід’ємними.
Набрані впродовж семестру бали не можуть перевищувати 60. 
Набрані на екзамені бали не можуть перевищувати 40. 
Набрані на екзамені бали повинні бути не менше 24 або дорівнювати 0(в останньому випадку оцінка не може бути задовільною). 
Сумарна оцінка в балах є сумою балів, набраних у семестрі та на екзамені.
Сумарна оцінка в балах та за державною шкалою мають бути узгоджені між собою.

Прізвище, ім’я та по-батькові (якщо наявне у варіанті) можуть мати до 23, 29, 29 відповідно символів включно кожне. 
Крім літер алфавіту можуть містити апостроф, дефіс та пробіл.

Предмет може мати від 5 до 52 символів включно. 
Крім літер алфавіту може містити подвійні лапки, пробіл та дефіс.

Код групи є непорожнім та складається не більш ніж з 5 літер. 
Крім літер алфавіту може містити десяткові цифри та дефіс.


Номер залікової книжки однозначно ідентифікує студента та складається з 8 десяткових цифр. 

\bigskip\textit{Для деяких варіантів може знадобитись}

Середній бал/оцінка за предметом визначається як середнє арифметичнесумарних балів/державних оцінок за цим предметом. 
Середній бал/оцінка студента визначається як середнє арифметичне сумарних балів/державних оцінокцього студента. 
За обчислення середньої оцінки недопущені та ті, що не з’явилися, рахуються як 2.
У решті випадків недопуски підсумовуються як $-1$, а «не з’явився» як 0.
Якщо студент не гірше ніж задовільно склав усі предмети, 
то його рейтинг за балами/оцінкою визначається як середній бал/оцінка; інакше дорівнює 0.

Стабільність студента визначається як модуль різниці між його кращим та найгіршим сумарним балом.
Стабільність предмета визначається як модуль різниці між найкращим та найгіршимсумарним балом за цим предметом.






\par
%end
\newpage
%begin
{\center Варіант  \No \textbf { 8 }\par}
\bigskip
%type = meteo
%variant = 103
Обробляються дані спостережень за погодою. 

\underline{Записи основного файлу містять поля}: максимальна денна температура, код метеостанції, кількість опадів, мінімальна денна температура, середня денна температура, відносна вологість, сила вітру, рік, день, місяць.

\underline{У допоміжному файлі наявні ключі}: середня сила вітру, середня відносна вологість.

\underline{Знайти} 5 метеостанцій, по яких у середньому випала найбільша кількість опадів. Вивести по кожному з них інформацію:\par
{\noindent}--- на першому рядку:\par
метеостанція, кількість спостережень, середня кількість опадів, найбільша сила вітру;

{\noindent}--- на наступних рядках, починаючи з табуляції, вивести для них дані по датах (по одному на рядок):\par
рік, місяць, день, середня денна температура, максимальна денна температура, мінімальна денна температура, сила вітру, вологість

{\noindent}у такому сортуванні: місяць, день, рік.

%Метеостанції сортувати: кількість спостережень (за спаданням), найбільша сила вітру, метеостанція.

\bigskip\textit{Обмеження}

Код метеостанції є рядком довжини від 5 до 7 включно.
Може містити літери алфавіту, десяткові цифри, підкреслення. Решта полів числові.

Рік є чотиризначним цілим.
Кількість опадів є дійсним значенням, подається у форматі з 1 знаком після крапки.
Сила вітру є дійсним значенням, подається у форматі з 2 знаками після крапки.
Температури є дійсними значеннями, подаються у форматі з 1 знаком після крапки.
Цілі значення подаються без десяткової крапки.
Для дійсних використовується зображення з фіксованою крапкою.

\textit{Вихідне форматування дійсних}

Усі дійсні значення виводяться у форматі з фіксованою крапкою з тією ж кількістю знаків після крапки, що і у вхідному зображенні.

Якщо у варіанті раніше не було зазначено, то \underline{обчислювані програмою} середні значення виводяться у форматі з 2 знаками після крапки.
%\bigskip%\textit{Для деяких варіантів може знадобитись}




\par
%end
\newpage
%begin
{\center Варіант  \No \textbf { 9 }\par}
\bigskip
%type = absense
%variant = 24
Обробляються дані семестрового журналу перевірок відвідування занять студентами деканатом (фіксуються відсутні). 

\underline{Записи основного файлу містять поля}: курс, код групи, навчальний тиждень, предмет, ім’я, прізвище, вид занять, день тижня, пара, аудиторія.

\underline{У допоміжному файлі наявні ключі}: найменша довжина назви предмета, найбільша довжина назви предмета, загальна кількість студентів.

\underline{Знайти} всіх студентів, що пропускали пари тільки з одного предмета. Вивести по кожному з них інформацію:\par
{\noindent}--- на першому рядку:\par
прізвище, ім’я, кількість пропущених пар, кількість пропущених лекцій, кількість пропущених лабораторних, предмет, що пропускався;

{\noindent}--- на наступних рядках, починаючи з табуляції, вивести пропуски студента (по одному на рядок):\par
тиждень, день, пара, аудиторія, вид занять

{\noindent}у такому сортуванні: пара, тиждень, день, вид занять.

%Студентів сортувати: прізвище, ім’я.

\bigskip\textit{Обмеження}

Навчальний тиждень може приймати цілі значення від 1 до 15.
День тижня є цілим від 1 до 5. Пара є цілим від 1 до 4.
Курс є цілим від 1 до 4. Аудиторія є додатнім цілим.
Вид занять може приймати значення:
Lect --- лекція, ПР --- практичне, Sem --- семінар, ЛАБ --- лабораторне.

Решта полів є непорожніми рядками.
Рядки починатися та закінчуватися білими символами не можуть.

Групи кодуються в межах курсу.
Код групи складається не більш ніж з 4 символів.
Може містити літери алфавіту, десяткові цифри та дефіс.

Предмет є рядком, може мати від 6 до 31 символів включно.
Може містити літери алфавіту, десяткові цифри, апостроф, дефіс, пробіл.
Предмети різних курсів вважати різними навіть за однакових назв.

Прізвище, ім’я та по-батькові (за наявності у варіанті)
можуть мати від 6 до 23, 25, 28 відповідно символів включно кожне.
Можуть містити літери алфавіту, апостроф, пробіл та дефіс.
Студенти однозначно ідентифікуються прізвищем, ім’ям та по-батькові (за наявності у варіанті).
Жоден студент не вчиться одночасно на різних курсах чи в різних групах.

Види занять сортуються так: лекція < практичне < лабораторне < семінар.

%\bigskip%\textit{Для деяких варіантів може знадобитись}




\par
%end
\newpage
%begin
{\center Варіант  \No \textbf { 10 }\par}
\bigskip
%type = tasks
%variant = 5
Обробляються результати розв’язування студентами задач на контрольних.
Одиничний факт роз\-в’я\-зу\-вання студентом задачі надалі розуміється як спроба.

\underline{Записи основного файлу містять поля}: відсоток набраних за розв’язання задачі балів, рік, по-бать\-ко\-ві, умова задачі, ім’я, прізвище.

\underline{У допоміжному файлі наявні ключі}: сумарна довжина прізвищ по всіх спробах, кількість записів.

\underline{Знайти} всі спроби, коли задача була розв’язана краще, ніж в середньому по цій задачі. Вивести по кожному з них інформацію:\par
{\noindent}--- на першому рядку:\par
умова, середня розв’язуваність задачі (округлена до 1 знаку);

{\noindent}--- на наступних рядках, починаючи з табуляції, вивести спроби по цих задачах (по одному на рядок):\par
рік, прізвище, ім’я, відсоток набраних балів

{\noindent}у такому сортуванні: відсоток набраних балів, рік, прізвище, ім’я.

%Задачі сортувати: середня розв’язуваність, умова.

\bigskip\textit{Обмеження}

Студент міг розв’язувати задачу кілька разів.

Відсоток набраних за розв’язання задачі балів є цілим від 0 до 100.
День, місяць, рік є числовими полями, що зображують дату.
Раніше 2001 року статистика не велася. Рік є 4-значним.
Решта полів є непорожніми рядками.
Рядки починатися та закінчуватися білими символами не можуть.

Прізвище, ім’я та по-батькові (за наявності у варіанті) можуть мати до 28, 30, 22 відповідно символів включно кожне. 
Крім літер алфавіту можуть містити апостроф, дефіс та пробіл.

Умова задачі може мати від 5 до 57 символів включно.
Крім літер алфавіту може містити подвійні лапки, десяткові цифри, пробіл, дефіс,
а також символи . , \_ \{ \} \$ $+$ $*$ {\textbackslash} ( ) \% $[$ $]$ .

Студенти однозначно ідентифікуються прізвищем, ім’ям та по-батькові (за наявності у варіанті).
Задачі однозначно ідентифікуються умовою.

\bigskip\textit{Для деяких варіантів може знадобитись}

Розв'язуваність задачі --- середній відсоток набраних за задачу балів.

Розв'язуваність студента --- середній відсоток набраних студентом балів по всіх його спробах.

Розв'язуваність --- середній відсоток набраних студентами за задачі балів (рахується по всім даним).




\par
%end
\newpage
%begin
{\center Варіант  \No \textbf { 11 }\par}
\bigskip
%type = booking
%variant = 128
Обробляється інформація щодо відвантажених накладних.

\underline{Записи основного файлу містять поля}: номер товарної позиції у межах накладної, номер накладної, назва товару, кількість, вартість, ціна.

\underline{У допоміжному файлі наявні ключі}: сума вартостей, сума цін.

\underline{Знайти} 5 накладних, в яких було відвантажено товарів на найменшу суму. Вивести по кожному з них інформацію:\par
{\noindent}--- на першому рядку:\par
середня ціна товарної позиції (з 2 знаками після крапки), загальна сума, номер накладної, кількість товарних позицій;

{\noindent}--- на наступних рядках, починаючи з табуляції, вивести для накладної всі її товарні позиції з найбільшою вартістю (по одному на рядок):\par
назва товару, ціна, кількість, вартість

{\noindent}у такому сортуванні: вартість, назва товару.

%Накладні сортувати: загальна сума, номер накладної.

\bigskip\textit{Обмеження}

Кількість, ціна та вартість є значеннями дійсного типу.
Решта полів є непорожніми рядками.
Рядки починатися та закінчуватися білими символами не можуть.

Номер накладної складається з 13 символів (цифри, латинські літери).
Назва товару може мати від 3 до 22 символів включно.
Може містити тільки цифри, літери, пробіли, апостроф, дефіс, підкреслення.

Кількість, ціна мають бути додатними.
Кількість вказується з рівно 4 знаками після крапки.
Ціна вказується з рівно 4 знаками після крапки.
Вартість має бути узгоджена з кількістю та ціною та записується рівно з двома десятковими знаками.
(Застосовується округлення.)

Товари однозначно ідентифікуються назвами та одиницями виміру (за наявності).

Одиниця виміру (за наявності у варіанті) може приймати значення: шт., кг, м , кор.

Кожна накладна має свій власний унікальний номер.
У накладній можуть відвантажуватися кілька товарів (не менше 1).
Товарна позиція --- рядок накладної: товар, кількість, ціна, вартість та, за наявності, одиниця виміру.
Товарні позиції кожної накладної нумеруються послідовними натуральними числами, починаючи з 1.
В одній накладній товари можуть повторюватися, 
тобто товарні позиції однієї накладної можуть містити однакові товари 
(при цьому ціна може бути і різною, і однаковою).

\bigskip\textit{Для деяких варіантів може знадобитись}

Недоотриману вигоду розраховувати як суму виразів
{\center (найбільша ціна відвантаження товару – ціна відвантаження товару)*кількість\par}

\bigskip
{\noindent} округлених до 2 знаків після крапки.




\par
%end
\newpage

%begin
{\center Варіант  \No \textbf { 13 }\par}
\bigskip
%type = absense
%variant = 42
Обробляються дані семестрового журналу перевірок відвідування занять студентами деканатом (фіксуються відсутні). 

\underline{Записи основного файлу містять поля}: курс, код групи, предмет, вид занять, прізвище, ім’я, навчальний тиждень, день тижня, пара, аудиторія.

\underline{У допоміжному файлі наявні ключі}: кількість різних предметів, кількість записів.

\underline{Знайти} всі предмети, заняття з яких пропускалися тільки в аудиторіях 18 та 42. Вивести по кожному з них інформацію:\par
{\noindent}--- на першому рядку:\par
курс, предмет, кількість пропусків, кількість пропусків лабораторних;

{\noindent}--- на наступних рядках, починаючи з табуляції, вивести пропуски предмета (по одному на рядок):\par
тиждень, день, пара, аудиторія, вид занять, кількість пропусків

{\noindent}у такому сортуванні: тиждень, день, пара, аудиторія, вид занять.

%Предмети сортувати: курс, предмет.

\bigskip\textit{Обмеження}

Навчальний тиждень може приймати цілі значення від 1 до 20.
День тижня є цілим від 1 до 5. Пара є цілим від 1 до 4.
Курс є цілим від 1 до 4. Аудиторія є додатнім цілим.
Вид занять може приймати значення:
lect --- лекція, P --- практичне, с. --- семінар, lab --- лабораторне.

Решта полів є непорожніми рядками.
Рядки починатися та закінчуватися білими символами не можуть.

Групи кодуються в межах курсу.
Код групи складається не більш ніж з 6 символів.
Може містити літери алфавіту, десяткові цифри та дефіс.

Предмет є рядком, може мати від 5 до 39 символів включно.
Може містити літери алфавіту, десяткові цифри, апостроф, дефіс, пробіл.
Предмети різних курсів вважати різними навіть за однакових назв.

Прізвище, ім’я та по-батькові (за наявності у варіанті)
можуть мати від 5 до 22, 30, 22 відповідно символів включно кожне.
Можуть містити літери алфавіту, апостроф, пробіл та дефіс.
Студенти однозначно ідентифікуються прізвищем, ім’ям та по-батькові (за наявності у варіанті).
Жоден студент не вчиться одночасно на різних курсах чи в різних групах.

Види занять сортуються так: лабораторне < практичне < лекція < семінар.

%\bigskip%\textit{Для деяких варіантів може знадобитись}




\par
%end
\newpage
%begin
{\center Варіант  \No \textbf { 14 }\par}
\bigskip
%type = meteo
%variant = 110
Обробляються дані спостережень за погодою. 

\underline{Записи основного файлу містять поля}: мінімальна денна температура, день, місяць, максимальна денна температура, відносна вологість, середня денна температура, код метеостанції, кількість опадів, сила вітру, рік.

\underline{У допоміжному файлі наявні ключі}: середня сила вітру, загальна кількість записів.

\underline{Знайти} 6 метеостанцій, на яких була зафіксована найбільша сила вітру. Вивести по кожному з них інформацію:\par
{\noindent}--- на першому рядку:\par
найбільша сила вітру по метеостанції, найменша вологість, середня кількість опадів, метеостанція;

{\noindent}--- на наступних рядках, починаючи з табуляції, вивести для них дані по датах, коли сила вітру була хоча б 50\% від максимальної зафіксованої (по одному на рядок):\par
місяць, день, вологість, середня денна температура, максимальна денна температура, мінімальна денна температура, сила вітру, рік

{\noindent}у такому сортуванні: рік, місяць, день.

%Метеостанції сортувати: сила вітру (за спаданням), метеостанція.

\bigskip\textit{Обмеження}

Код метеостанції є рядком довжини від 4 до 12 включно.
Може містити літери алфавіту, десяткові цифри, підкреслення. Решта полів числові.

Рік є чотиризначним цілим.
Кількість опадів є дійсним значенням, подається у форматі з 2 знаками після крапки.
Сила вітру є цілим значенням.
Температури є дійсними значеннями, подаються у форматі з 3 знаками після крапки.
Цілі значення подаються без десяткової крапки.
Для дійсних використовується зображення з фіксованою крапкою.

\textit{Вихідне форматування дійсних}

Усі дійсні значення виводяться у форматі з фіксованою крапкою з тією ж кількістю знаків після крапки, що і у вхідному зображенні.

Якщо у варіанті раніше не було зазначено, то \underline{обчислювані програмою} середні значення виводяться у форматі з 1 знаком після крапки.
%\bigskip%\textit{Для деяких варіантів може знадобитись}




\par
%end
\newpage
%begin
{\center Варіант  \No \textbf { 15 }\par}
\bigskip
%type = meteo
%variant = 90
Обробляються дані спостережень за погодою. 

\underline{Записи основного файлу містять поля}: рік, максимальна денна температура, день, місяць, кількість опадів, код метеостанції, мінімальна денна температура, відносна вологість, сила вітру, середня денна температура.

\underline{У допоміжному файлі наявні ключі}: кількість спостережень, коли сила вітру перевищувала 15; загальна кількість записів.

\underline{Знайти} 7 дат, по яких фіксувалася найбільша мінімальна денна температура. Вивести по кожному з них інформацію:\par
{\noindent}--- на першому рядку:\par
рік, місяць, день, мінімальна температура, середня температура, максимальна температура, вологість;

{\noindent}--- на наступних рядках, починаючи з табуляції, вивести для дати спостереження метеостанції, де сила вітру була меша середньої за дату (по одному на рядок):\par
сила вітру, метеостанція, вологість, різниця максимальної та мінімальної денних температур

{\noindent}у такому сортуванні: різниця максимальної та мінімальної температур, метеостанція.

%Дати сортувати: вологість, місяць, день, рік.

\bigskip\textit{Обмеження}

Код метеостанції є рядком довжини від 3 до 6 включно.
Може містити літери алфавіту, десяткові цифри, підкреслення. Решта полів числові.

Рік є чотиризначним цілим.
Кількість опадів є цілим значенням.
Сила вітру є дійсним значенням, подається у форматі з 1 знаком після крапки.
Температури є дійсними значеннями, подаються у форматі з 2 знаками після крапки.
Цілі значення подаються без десяткової крапки.
Для дійсних використовується зображення з фіксованою крапкою.

\textit{Вихідне форматування дійсних}

Усі дійсні значення виводяться у форматі з фіксованою крапкою з тією ж кількістю знаків після крапки, що і у вхідному зображенні.

Якщо у варіанті раніше не було зазначено, то \underline{обчислювані програмою} середні значення виводяться у форматі з 3 знаками після крапки.
%\bigskip%\textit{Для деяких варіантів може знадобитись}




\par
%end
\newpage

%begin
{\center Варіант  \No \textbf { 17 }\par}
\bigskip
%type = session
%variant = 71
Обробляються результати складання студентами екзаменів зимової сесії.

\underline{Записи основного файлу містять поля}: номер залікової книжки, предмет, код групи, сумарна оцінка в балах за 100-бальною системою, набрані на екзамені бали, набрані впродовж семестру бали, прізвище, ім’я, оцінка за державною шкалою, по-батькові.

\underline{У допоміжному файлі наявні ключі}: довжина найдовшого прізвища, загальна кількість «хвостів».

\underline{Знайти} предмети з найкращим середнім балом. Вивести по кожному з них інформацію:\par
{\noindent}--- на першому рядку:\par
предмет, середній бал (округлений до цілого), середня оцінка (округлена до 1 знаку);

{\noindent}--- на наступних рядках, починаючи з табуляції, вивести для предмета всіх студентів, що його складали (по одному на рядок):\par
оцінка в балах, бали за екзамен, прізвище, ім’я, група, номер залікової книжки

{\noindent}у такому сортуванні: прізвище, ім’я, номер залікової книжки.

%Предмети сортувати: середній бал (округлений), предмет.

\bigskip\textit{Обмеження}

Студент не може складати предмет більше одного разу.

Бали та оцінка є значеннями цілих числових типів. Решта полів є непорожніми рядками.
Рядки розпочинатись та закінчуватись білими символами не можуть.

Оцінка за державною шкалою може приймати значення: 2-5~--- як зазвичай, 0~--- не з’явився,$-1$~--- не допущений. 
Оцінки 3-5 вважаються задовільними, решта~--- незадовільними. 
Бали є цілими та невід’ємними.
Набрані впродовж семестру бали не можуть перевищувати 60. 
Набрані на екзамені бали не можуть перевищувати 40. 
Набрані на екзамені бали повинні бути не менше 24 або дорівнювати 0(в останньому випадку оцінка не може бути задовільною). 
Сумарна оцінка в балах є сумою балів, набраних у семестрі та на екзамені.
Сумарна оцінка в балах та за державною шкалою мають бути узгоджені між собою.

Прізвище, ім’я та по-батькові (якщо наявне у варіанті) можуть мати до 21, 21, 28 відповідно символів включно кожне. 
Крім літер алфавіту можуть містити апостроф, дефіс та пробіл.

Предмет може мати від 3 до 54 символів включно. 
Крім літер алфавіту може містити подвійні лапки, пробіл та дефіс.

Код групи є непорожнім та складається не більш ніж з 7 літер. 
Крім літер алфавіту може містити десяткові цифри та дефіс.


Номер залікової книжки однозначно ідентифікує студента та складається з 7 десяткових цифр. 

\bigskip\textit{Для деяких варіантів може знадобитись}

Середній бал/оцінка за предметом визначається як середнє арифметичнесумарних балів/державних оцінок за цим предметом. 
Середній бал/оцінка студента визначається як середнє арифметичне сумарних балів/державних оцінокцього студента. 
За обчислення середньої оцінки недопущені та ті, що не з’явилися, рахуються як 2.
У решті випадків недопуски підсумовуються як $-1$, а «не з’явився» як 0.
Якщо студент не гірше ніж задовільно склав усі предмети, 
то його рейтинг за балами/оцінкою визначається як середній бал/оцінка; інакше дорівнює 0.

Стабільність студента визначається як модуль різниці між його кращим та найгіршим сумарним балом.
Стабільність предмета визначається як модуль різниці між найкращим та найгіршимсумарним балом за цим предметом.






\par
%end
\newpage
%begin
{\center Варіант  \No \textbf { 18 }\par}
\bigskip
%type = meteo
%variant = 106
Обробляються дані спостережень за погодою. 

\underline{Записи основного файлу містять поля}: середня денна температура, максимальна денна температура, відносна вологість, рік, мінімальна денна температура, кількість опадів, день, місяць, код метеостанції, сила вітру.

\underline{У допоміжному файлі наявні ключі}: найбільша вологість, найменша вологість, кількість різних метеостанцій.

\underline{Знайти} 7 метеостанцій, по яких у середньому була найвища мінімальна денна температура. Вивести по кожному з них інформацію:\par
{\noindent}--- на першому рядку:\par
метеостанція, середня мінімальна температура, найбільша вологість;

{\noindent}--- на наступних рядках, починаючи з табуляції, вивести для них дані по датах (по одному на рядок):\par
місяць, день, середня денна температура, максимальна денна температура, мінімальна денна температура, сила вітру, вологість, рік

{\noindent}у такому сортуванні: сила вітру, місяць, день, рік.

%Метеостанції сортувати: кількість спостережень (за спаданням), найбільша вологість (за спаданням), метеостанція.

\bigskip\textit{Обмеження}

Код метеостанції є рядком довжини від 3 до 10 включно.
Може містити літери алфавіту, десяткові цифри, підкреслення. Решта полів числові.

Рік є чотиризначним цілим.
Кількість опадів є дійсним значенням, подається у форматі з 2 знаками після крапки.
Сила вітру є цілим значенням.
Температури є дійсними значеннями, подаються у форматі з 1 знаком після крапки.
Цілі значення подаються без десяткової крапки.
Для дійсних використовується зображення з фіксованою крапкою.

\textit{Вихідне форматування дійсних}

Усі дійсні значення виводяться у форматі з фіксованою крапкою з тією ж кількістю знаків після крапки, що і у вхідному зображенні.

Якщо у варіанті раніше не було зазначено, то \underline{обчислювані програмою} середні значення виводяться у форматі з 2 знаками після крапки.
%\bigskip%\textit{Для деяких варіантів може знадобитись}




\par
%end
\newpage
%begin
{\center Варіант  \No \textbf { 19 }\par}
\bigskip
%type = tasks
%variant = 4
Обробляються результати розв’язування студентами задач на контрольних.
Одиничний факт роз\-в’я\-зу\-вання студентом задачі надалі розуміється як спроба.

\underline{Записи основного файлу містять поля}: відсоток набраних за розв’язання задачі балів, умова задачі, прізвище, місяць, рік, день, ім’я.

\underline{У допоміжному файлі наявні ключі}: сума років по всіх спробах, кількість різних студентів.

\underline{Знайти} всіх студентів, хто в середньому (округлити до цілого) розв’язував задачі найкраще. Вивести по кожному з них інформацію:\par
{\noindent}--- на першому рядку:\par
прізвище, ім’я, середня розв’язуваність студента, кількість спроб;

{\noindent}--- на наступних рядках, починаючи з табуляції, вивести їх спроби (по одному на рядок):\par
умова, рік, відсоток набраних за задачу балів

{\noindent}у такому сортуванні: відсоток набраних балів (за спаданням), рік (за спаданням), умова.

%Студентів сортувати: ім’я, прізвище.

\bigskip\textit{Обмеження}

Студент міг розв’язувати задачу кілька разів.

Відсоток набраних за розв’язання задачі балів є цілим від 0 до 100.
День, місяць, рік є числовими полями, що зображують дату.
Раніше 2001 року статистика не велася. Рік є 4-значним.
Решта полів є непорожніми рядками.
Рядки починатися та закінчуватися білими символами не можуть.

Прізвище, ім’я та по-батькові (за наявності у варіанті) можуть мати до 30, 25, 20 відповідно символів включно кожне. 
Крім літер алфавіту можуть містити апостроф, дефіс та пробіл.

Умова задачі може мати від 4 до 52 символів включно.
Крім літер алфавіту може містити подвійні лапки, десяткові цифри, пробіл, дефіс,
а також символи . , \_ \{ \} \$ $+$ $*$ {\textbackslash} ( ) \% $[$ $]$ .

Студенти однозначно ідентифікуються прізвищем, ім’ям та по-батькові (за наявності у варіанті).
Задачі однозначно ідентифікуються умовою.

\bigskip\textit{Для деяких варіантів може знадобитись}

Розв'язуваність задачі --- середній відсоток набраних за задачу балів.

Розв'язуваність студента --- середній відсоток набраних студентом балів по всіх його спробах.

Розв'язуваність --- середній відсоток набраних студентами за задачі балів (рахується по всім даним).




\par
%end
\newpage
%begin
{\center Варіант  \No \textbf { 20 }\par}
\bigskip
%type = tasks
%variant = 18
Обробляються результати розв’язування студентами задач на контрольних.
Одиничний факт роз\-в’я\-зу\-вання студентом задачі надалі розуміється як спроба.

\underline{Записи основного файлу містять поля}: по-батькові, прізвище, день, відсоток набраних за розв’язан\-ня задачі балів, місяць, умова задачі, ім’я, рік.

\underline{У допоміжному файлі наявні ключі}: найкоротша умова, кількість записів.

\underline{Знайти} всіх студентів, що зробили рівно одну спробу з відсотком 99\%. Вивести по кожному з них інформацію:\par
{\noindent}--- на першому рядку:\par
прізвище, ім’я, по-батькові, найкраща спроба, найгірша спроба, кількість спроб з відсотком 99\%;

{\noindent}--- на наступних рядках, починаючи з табуляції, вивести їх спроби (по одному на рядок):\par
рік, місяць, день, відсоток набраних за задачу балів, умова

{\noindent}у такому сортуванні: рік(за спаданням), місяць (за спаданням), день (за спаданням), відсоток (за спаданням), умова.

%Студентів сортувати: найкраща спроба (за спаданням), найгірша спроба, прізвище, ім’я, по-батькові.

\bigskip\textit{Обмеження}

Студент міг розв’язувати задачу кілька разів.

Відсоток набраних за розв’язання задачі балів є цілим від 0 до 100.
День, місяць, рік є числовими полями, що зображують дату.
Раніше 2001 року статистика не велася. Рік є 4-значним.
Решта полів є непорожніми рядками.
Рядки починатися та закінчуватися білими символами не можуть.

Прізвище, ім’я та по-батькові (за наявності у варіанті) можуть мати до 22, 27, 26 відповідно символів включно кожне. 
Крім літер алфавіту можуть містити апостроф, дефіс та пробіл.

Умова задачі може мати від 5 до 63 символів включно.
Крім літер алфавіту може містити подвійні лапки, десяткові цифри, пробіл, дефіс,
а також символи . , \_ \{ \} \$ $+$ $*$ {\textbackslash} ( ) \% $[$ $]$ .

Студенти однозначно ідентифікуються прізвищем, ім’ям та по-батькові (за наявності у варіанті).
Задачі однозначно ідентифікуються умовою.

\bigskip\textit{Для деяких варіантів може знадобитись}

Розв'язуваність задачі --- середній відсоток набраних за задачу балів.

Розв'язуваність студента --- середній відсоток набраних студентом балів по всіх його спробах.

Розв'язуваність --- середній відсоток набраних студентами за задачі балів (рахується по всім даним).




\par
%end
\newpage
%begin
{\center Варіант  \No \textbf { 21 }\par}
\bigskip
%type = meteo
%variant = 95
Обробляються дані спостережень за погодою. 

\underline{Записи основного файлу містять поля}: відносна вологість, код метеостанції, рік, кількість опадів, середня денна температура, максимальна денна температура, місяць, день, мінімальна денна температура, сила вітру.

\underline{У допоміжному файлі наявні ключі}: максимальна сила вітру, що спостерігалася; загальна кількість записів.

\underline{Знайти} дати, по яких фіксувалася середньоденна температура більша ніж в середньому по всіх спостереженнях. Вивести по кожному з них інформацію:\par
{\noindent}--- на першому рядку:\par
рік, місяць, день, кількість опадів за місяць дати;

{\noindent}--- на наступних рядках, починаючи з табуляції, вивести для дати спостереження метеостанції, по яких спостерігалися опади більше 5 (по одному на рядок):\par
метеостанція, кількість опадів, сила вітру

{\noindent}у такому сортуванні: кількість опадів, метеостанція.

%Дати сортувати: рік, місяць, день.

\bigskip\textit{Обмеження}

Код метеостанції є рядком довжини від 5 до 12 включно.
Може містити літери алфавіту, десяткові цифри, підкреслення. Решта полів числові.

Рік є чотиризначним цілим.
Кількість опадів є дійсним значенням, подається у форматі з 1 знаком після крапки.
Сила вітру є дійсним значенням, подається у форматі з 1 знаком після крапки.
Температури є дійсними значеннями, подаються у форматі з 2 знаками після крапки.
Цілі значення подаються без десяткової крапки.
Для дійсних використовується зображення з фіксованою крапкою.

\textit{Вихідне форматування дійсних}

Усі дійсні значення виводяться у форматі з фіксованою крапкою з тією ж кількістю знаків після крапки, що і у вхідному зображенні.

Якщо у варіанті раніше не було зазначено, то \underline{обчислювані програмою} середні значення виводяться у форматі з 1 знаком після крапки.
%\bigskip%\textit{Для деяких варіантів може знадобитись}




\par
%end
\newpage
%begin
{\center Варіант  \No \textbf { 22 }\par}
\bigskip
%type =     tasks
%variant = 1
Обробляються результати розв’язування студентами задач на контрольних.
Одиничний факт роз\-в’я\-зу\-вання студентом задачі надалі розуміється як спроба.

\underline{Записи основного файлу містять поля}: по-батькові, рік, відсоток набраних за розв’язання задачі балів, місяць, ім’я, прізвище, день, умова задачі.

\underline{У допоміжному файлі наявні ключі}: кількість різних задач, кількість записів.

\underline{Знайти} всі задачі, які були розв’язані на 100\% у найбільшій кількості спроб. Вивести по кожному з них інформацію:\par
{\noindent}--- на першому рядку:\par
умова, кількість спроб на 100\%, кількість спроб;

{\noindent}--- на наступних рядках, починаючи з табуляції, вивести спроби, що набрали якнайменш 75\% за цю задачу (по одному на рядок):\par
рік, місяць, відсоток набраних за задачу балів, прізвище, ім’я, по-батькові

{\noindent}у такому сортуванні: відсоток набраних за задачу балів (за спаданням), рік, місяць, прізвище, ім’я, по-батькові.

%Задачі сортувати: умова.

\bigskip\textit{Обмеження}

Студент міг розв’язувати задачу кілька разів.

Відсоток набраних за розв’язання задачі балів є цілим від 0 до 100.
День, місяць, рік є числовими полями, що зображують дату.
Раніше 2001 року статистика не велася. Рік є 4-значним.
Решта полів є непорожніми рядками.
Рядки починатися та закінчуватися білими символами не можуть.

Прізвище, ім’я та по-батькові (за наявності у варіанті) можуть мати до 26, 28, 27 відповідно символів включно кожне. 
Крім літер алфавіту можуть містити апостроф, дефіс та пробіл.

Умова задачі може мати від 4 до 68 символів включно.
Крім літер алфавіту може містити подвійні лапки, десяткові цифри, пробіл, дефіс,
а також символи . , \_ \{ \} \$ $+$ $*$ {\textbackslash} ( ) \% $[$ $]$ .

Студенти однозначно ідентифікуються прізвищем, ім’ям та по-батькові (за наявності у варіанті).
Задачі однозначно ідентифікуються умовою.

\bigskip\textit{Для деяких варіантів може знадобитись}

Розв'язуваність задачі --- середній відсоток набраних за задачу балів.

Розв'язуваність студента --- середній відсоток набраних студентом балів по всіх його спробах.

Розв'язуваність --- середній відсоток набраних студентами за задачі балів (рахується по всім даним).




\par
%end
\newpage
%begin
{\center Варіант  \No \textbf { 23 }\par}
\bigskip
%type = meteo
%variant = 108
Обробляються дані спостережень за погодою. 

\underline{Записи основного файлу містять поля}: сила вітру, відносна вологість, середня денна температура, кількість опадів, мінімальна денна температура, максимальна денна температура, день, місяць, рік, код метеостанції.

\underline{У допоміжному файлі наявні ключі}: найбільша середня температура, найменша вологість.

\underline{Знайти} 8 метеостанцій, на яких була зафіксована найменша максимальна денна температура. Вивести по кожному з них інформацію:\par
{\noindent}--- на першому рядку:\par
середні значення показників: середня, максимальна та мінімальна денна  температура, вологість, сила вітру; та код метеостанції;

{\noindent}--- на наступних рядках, починаючи з табуляції, вивести для них дані по датах (по одному на рядок):\par
день, рік, місяць, різниця максимальної та мінімальної температур, вологість

{\noindent}у такому сортуванні: вологість, місяць, день, рік.

%Метеостанція сортувати: середня вологість, середня сила вітру, метеостанція.

\bigskip\textit{Обмеження}

Код метеостанції є рядком довжини від 4 до 6 включно.
Може містити літери алфавіту, десяткові цифри, підкреслення. Решта полів числові.

Рік є чотиризначним цілим.
Кількість опадів є цілим значенням.
Сила вітру є дійсним значенням, подається у форматі з 2 знаками після крапки.
Температури є дійсними значеннями, подаються у форматі з 3 знаками після крапки.
Цілі значення подаються без десяткової крапки.
Для дійсних використовується зображення з фіксованою крапкою.

\textit{Вихідне форматування дійсних}

Усі дійсні значення виводяться у форматі з фіксованою крапкою з тією ж кількістю знаків після крапки, що і у вхідному зображенні.

Якщо у варіанті раніше не було зазначено, то \underline{обчислювані програмою} середні значення виводяться у форматі з 3 знаками після крапки.
%\bigskip%\textit{Для деяких варіантів може знадобитись}




\par
%end
\newpage
%begin
{\center Варіант  \No \textbf { 24 }\par}
\bigskip
%type = booking
%variant = 130
Обробляється інформація щодо відвантажених накладних.

\underline{Записи основного файлу містять поля}: назва товару, вартість, кількість, ціна, номер товарної позиції у межах накладної, одиниця виміру, номер накладної.

\underline{У допоміжному файлі наявні ключі}: найбільша ціна товарної позиції, кількість накладних.

\underline{Знайти} 10 накладних, в яких було відвантажено найменшу кількість різних товарів. Вивести по кожному з них інформацію:\par
{\noindent}--- на першому рядку:\par
номер накладної, кількість різних товарів у накладній, найбільша вартість товарної позиції накладної;

{\noindent}--- на наступних рядках, починаючи з табуляції, вивести для накладної всі її товарні позиції (по одному на рядок):\par
номер товарної позиції, назва товару, ціна, кількість, вартість

{\noindent}у такому сортуванні: назва товару, ціна (за спаданням).

%Накладні сортувати: кількість різних товарів, номер накладної.

\bigskip\textit{Обмеження}

Кількість, ціна та вартість є значеннями дійсного типу.
Решта полів є непорожніми рядками.
Рядки починатися та закінчуватися білими символами не можуть.

Номер накладної складається з 12 символів (цифри, латинські літери).
Назва товару може мати від 1 до 21 символів включно.
Може містити тільки цифри, літери, пробіли, апостроф, дефіс, підкреслення.

Кількість, ціна мають бути додатними.
Кількість вказується з рівно 2 знаками після крапки.
Ціна вказується з рівно 3 знаками після крапки.
Вартість має бути узгоджена з кількістю та ціною та записується рівно з двома десятковими знаками.
(Застосовується округлення.)

Товари однозначно ідентифікуються назвами та одиницями виміру (за наявності).

Одиниця виміру (за наявності у варіанті) може приймати значення: шт., кг, м , кор.

Кожна накладна має свій власний унікальний номер.
У накладній можуть відвантажуватися кілька товарів (не менше 1).
Товарна позиція --- рядок накладної: товар, кількість, ціна, вартість та, за наявності, одиниця виміру.
Товарні позиції кожної накладної нумеруються послідовними натуральними числами, починаючи з 1.
В одній накладній товари можуть повторюватися, 
тобто товарні позиції однієї накладної можуть містити однакові товари 
(при цьому ціна може бути і різною, і однаковою).

\bigskip\textit{Для деяких варіантів може знадобитись}

Недоотриману вигоду розраховувати як суму виразів
{\center (найбільша ціна відвантаження товару – ціна відвантаження товару)*кількість\par}

\bigskip
{\noindent} округлених до 2 знаків після крапки.




\par
%end
\newpage
%begin
{\center Варіант  \No \textbf { 25 }\par}
\bigskip
%type = tasks
%variant = 14
Обробляються результати розв’язування студентами задач на контрольних.
Одиничний факт роз\-в’я\-зу\-вання студентом задачі надалі розуміється як спроба.

\underline{Записи основного файлу містять поля}: день, прізвище, місяць, умова задачі, відсоток набраних за розв’язання задачі балів, ім’я, рік.

\underline{У допоміжному файлі наявні ключі}: загальна довжина прізвищ, найбільший відсоток набраних балів.

\underline{Знайти} всі задачі, за якими більше половини спроб отримало більше 60\%. Вивести по кожному з них інформацію:\par
{\noindent}--- на першому рядку:\par
умова, кількість спроб, кількість спроб кращих за 60\%, відсоток найкращої спроби;

{\noindent}--- на наступних рядках, починаючи з табуляції, вивести спроби по цих задачах (по одному на рядок):\par
прізвище, ім’я, відсоток набраних балів, рік, місяць, день

{\noindent}у такому сортуванні: рік, місяць, день, відсоток набраних балів, прізвище, ім’я.

%Задачі сортувати: відсоток найкращої спроби, кількість спроб, умова.

\bigskip\textit{Обмеження}

Студент міг розв’язувати задачу кілька разів.

Відсоток набраних за розв’язання задачі балів є цілим від 0 до 100.
День, місяць, рік є числовими полями, що зображують дату.
Раніше 2001 року статистика не велася. Рік є 4-значним.
Решта полів є непорожніми рядками.
Рядки починатися та закінчуватися білими символами не можуть.

Прізвище, ім’я та по-батькові (за наявності у варіанті) можуть мати до 23, 21, 21 відповідно символів включно кожне. 
Крім літер алфавіту можуть містити апостроф, дефіс та пробіл.

Умова задачі може мати від 3 до 56 символів включно.
Крім літер алфавіту може містити подвійні лапки, десяткові цифри, пробіл, дефіс,
а також символи . , \_ \{ \} \$ $+$ $*$ {\textbackslash} ( ) \% $[$ $]$ .

Студенти однозначно ідентифікуються прізвищем, ім’ям та по-батькові (за наявності у варіанті).
Задачі однозначно ідентифікуються умовою.

\bigskip\textit{Для деяких варіантів може знадобитись}

Розв'язуваність задачі --- середній відсоток набраних за задачу балів.

Розв'язуваність студента --- середній відсоток набраних студентом балів по всіх його спробах.

Розв'язуваність --- середній відсоток набраних студентами за задачі балів (рахується по всім даним).




\par
%end
\newpage
%begin
{\center Варіант  \No \textbf { 26 }\par}
\bigskip
%type = session
%variant = 79
Обробляються результати складання студентами екзаменів зимової сесії.

\underline{Записи основного файлу містять поля}: прізвище, набрані на екзамені бали, предмет, номер залікової книжки, оцінка за державною шкалою, ім’я, код групи, сумарна оцінка в балах за 100-бальною системою, набрані впродовж семестру бали.

\underline{У допоміжному файлі наявні ключі}: FAILED SUCCESSFULLY, загальна кількість записів, найменший номер залікової книжки.

\underline{Знайти} предмети, за якими оцінок «відмінно» було більше, ніж оцінок «задовільно». Вивести по кожному з них інформацію:\par
{\noindent}--- на першому рядку:\par
кількість «відмінно» -(мінус) кількість «задовільно», предмет;

{\noindent}--- на наступних рядках, починаючи з табуляції, вивести для предмета тільки тих, хто склав на «відмінно» або «задовільно» (по одному на рядок):\par
ім’я, прізвище, номер залікової книжки, група, оцінка в балах

{\noindent}у такому сортуванні: прізвище, ім’я, номер залікової книжки.

%Предмети сортувати: за першим полем, що виводиться, предмет.

\bigskip\textit{Обмеження}

Студент не може складати предмет більше одного разу.

Бали та оцінка є значеннями цілих числових типів. Решта полів є непорожніми рядками.
Рядки розпочинатись та закінчуватись білими символами не можуть.

Оцінка за державною шкалою може приймати значення: 2-5~--- як зазвичай, 0~--- не з’явився,$-1$~--- не допущений. 
Оцінки 3-5 вважаються задовільними, решта~--- незадовільними. 
Бали є цілими та невід’ємними.
Набрані впродовж семестру бали не можуть перевищувати 60. 
Набрані на екзамені бали не можуть перевищувати 40. 
Набрані на екзамені бали повинні бути не менше 24 або дорівнювати 0(в останньому випадку оцінка не може бути задовільною). 
Сумарна оцінка в балах є сумою балів, набраних у семестрі та на екзамені.
Сумарна оцінка в балах та за державною шкалою мають бути узгоджені між собою.

Прізвище, ім’я та по-батькові (якщо наявне у варіанті) можуть мати до 29, 25, 24 відповідно символів включно кожне. 
Крім літер алфавіту можуть містити апостроф, дефіс та пробіл.

Предмет може мати від 5 до 68 символів включно. 
Крім літер алфавіту може містити подвійні лапки, пробіл та дефіс.

Код групи є непорожнім та складається не більш ніж з 3 літер. 
Крім літер алфавіту може містити десяткові цифри та дефіс.


Номер залікової книжки однозначно ідентифікує студента та складається з 8 десяткових цифр. 

\bigskip\textit{Для деяких варіантів може знадобитись}

Середній бал/оцінка за предметом визначається як середнє арифметичнесумарних балів/державних оцінок за цим предметом. 
Середній бал/оцінка студента визначається як середнє арифметичне сумарних балів/державних оцінокцього студента. 
За обчислення середньої оцінки недопущені та ті, що не з’явилися, рахуються як 2.
У решті випадків недопуски підсумовуються як $-1$, а «не з’явився» як 0.
Якщо студент не гірше ніж задовільно склав усі предмети, 
то його рейтинг за балами/оцінкою визначається як середній бал/оцінка; інакше дорівнює 0.

Стабільність студента визначається як модуль різниці між його кращим та найгіршим сумарним балом.
Стабільність предмета визначається як модуль різниці між найкращим та найгіршимсумарним балом за цим предметом.






\par
%end
\newpage
%begin
{\center Варіант  \No \textbf { 27 }\par}
\bigskip
%type = session
%variant = 65
Обробляються результати складання студентами екзаменів зимової сесії.

\underline{Записи основного файлу містять поля}: оцінка за державною шкалою, набрані впродовж семестру бали, набрані на екзамені бали, номер залікової книжки, код групи, прізвище, по-батькові, сумарна оцінка в балах за 100-бальною системою, предмет, ім’я.

\underline{У допоміжному файлі наявні ключі}: кількість записів, скільки разів на екзамені набирали 24 бали.

\underline{Знайти} всіх студентів, чий рейтинг за оцінкою вище середнього рейтингу. Вивести по кожному з них інформацію:\par
{\noindent}--- на першому рядку:\par
ім’я, по-батькові, прізвище, кількість екзаменів, сума балів, номер залікової книжки, рейтинг за оцінкою (округлений до цілого);

{\noindent}--- на наступних рядках, починаючи з табуляції, вивести для студента всі його результати (по одному на рядок):\par
державна оцінка, сумарна оцінка, предмет

{\noindent}у такому сортуванні: сумарний бал, предмет.

%Студентів сортувати: ім’я, прізвище, рейтинг за оцінкою( після округлення), номер залікової книжки.

\bigskip\textit{Обмеження}

Студент не може складати предмет більше одного разу.

Бали та оцінка є значеннями цілих числових типів. Решта полів є непорожніми рядками.
Рядки розпочинатись та закінчуватись білими символами не можуть.

Оцінка за державною шкалою може приймати значення: 2-5~--- як зазвичай, 0~--- не з’явився,$-1$~--- не допущений. 
Оцінки 3-5 вважаються задовільними, решта~--- незадовільними. 
Бали є цілими та невід’ємними.
Набрані впродовж семестру бали не можуть перевищувати 60. 
Набрані на екзамені бали не можуть перевищувати 40. 
Набрані на екзамені бали повинні бути не менше 24 або дорівнювати 0(в останньому випадку оцінка не може бути задовільною). 
Сумарна оцінка в балах є сумою балів, набраних у семестрі та на екзамені.
Сумарна оцінка в балах та за державною шкалою мають бути узгоджені між собою.

Прізвище, ім’я та по-батькові (якщо наявне у варіанті) можуть мати до 27, 24, 22 відповідно символів включно кожне. 
Крім літер алфавіту можуть містити апостроф, дефіс та пробіл.

Предмет може мати від 4 до 55 символів включно. 
Крім літер алфавіту може містити подвійні лапки, пробіл та дефіс.

Код групи є непорожнім та складається не більш ніж з 3 літер. 
Крім літер алфавіту може містити десяткові цифри та дефіс.


Номер залікової книжки однозначно ідентифікує студента та складається з 6 десяткових цифр. 

\bigskip\textit{Для деяких варіантів може знадобитись}

Середній бал/оцінка за предметом визначається як середнє арифметичнесумарних балів/державних оцінок за цим предметом. 
Середній бал/оцінка студента визначається як середнє арифметичне сумарних балів/державних оцінокцього студента. 
За обчислення середньої оцінки недопущені та ті, що не з’явилися, рахуються як 2.
У решті випадків недопуски підсумовуються як $-1$, а «не з’явився» як 0.
Якщо студент не гірше ніж задовільно склав усі предмети, 
то його рейтинг за балами/оцінкою визначається як середній бал/оцінка; інакше дорівнює 0.

Стабільність студента визначається як модуль різниці між його кращим та найгіршим сумарним балом.
Стабільність предмета визначається як модуль різниці між найкращим та найгіршимсумарним балом за цим предметом.






\par
%end
\newpage
%begin
{\center Варіант  \No \textbf { 28 }\par}
\bigskip
%type = session
%variant = 73
Обробляються результати складання студентами екзаменів зимової сесії.

\underline{Записи основного файлу містять поля}: код групи, оцінка за державною шкалою, предмет, сумарна оцінка в балах за 100-бальною системою, набрані впродовж семестру бали, набрані на екзамені бали, ім’я, прізвище, по-батькові, номер залікової книжки.

\underline{У допоміжному файлі наявні ключі}: кількість предметів, сума балів.

\underline{Знайти} найбільш стабільні предмети. Вивести по кожному з них інформацію:\par
{\noindent}--- на першому рядку:\par
предмет, найменший бал, найбільший бал, стабільність, кількість студентів, що успішно склали;

{\noindent}--- на наступних рядках, починаючи з табуляції, вивести для предмета всіх студентів, що його складали (по одному на рядок):\par
група, середня оцінка по групі в балах (округлена до 2 знаків), кількість студентів групи

{\noindent}у такому сортуванні: група.

%Предмети сортувати: середній бал (округлений)(за спаданням), предмет.

\bigskip\textit{Обмеження}

Студент не може складати предмет більше одного разу.

Бали та оцінка є значеннями цілих числових типів. Решта полів є непорожніми рядками.
Рядки розпочинатись та закінчуватись білими символами не можуть.

Оцінка за державною шкалою може приймати значення: 2-5~--- як зазвичай, 0~--- не з’явився,$-1$~--- не допущений. 
Оцінки 3-5 вважаються задовільними, решта~--- незадовільними. 
Бали є цілими та невід’ємними.
Набрані впродовж семестру бали не можуть перевищувати 60. 
Набрані на екзамені бали не можуть перевищувати 40. 
Набрані на екзамені бали повинні бути не менше 24 або дорівнювати 0(в останньому випадку оцінка не може бути задовільною). 
Сумарна оцінка в балах є сумою балів, набраних у семестрі та на екзамені.
Сумарна оцінка в балах та за державною шкалою мають бути узгоджені між собою.

Прізвище, ім’я та по-батькові (якщо наявне у варіанті) можуть мати до 22, 20, 23 відповідно символів включно кожне. 
Крім літер алфавіту можуть містити апостроф, дефіс та пробіл.

Предмет може мати від 3 до 66 символів включно. 
Крім літер алфавіту може містити подвійні лапки, пробіл та дефіс.

Код групи є непорожнім та складається не більш ніж з 5 літер. 
Крім літер алфавіту може містити десяткові цифри та дефіс.


Номер залікової книжки однозначно ідентифікує студента та складається з 8 десяткових цифр. 

\bigskip\textit{Для деяких варіантів може знадобитись}

Середній бал/оцінка за предметом визначається як середнє арифметичнесумарних балів/державних оцінок за цим предметом. 
Середній бал/оцінка студента визначається як середнє арифметичне сумарних балів/державних оцінокцього студента. 
За обчислення середньої оцінки недопущені та ті, що не з’явилися, рахуються як 2.
У решті випадків недопуски підсумовуються як $-1$, а «не з’явився» як 0.
Якщо студент не гірше ніж задовільно склав усі предмети, 
то його рейтинг за балами/оцінкою визначається як середній бал/оцінка; інакше дорівнює 0.

Стабільність студента визначається як модуль різниці між його кращим та найгіршим сумарним балом.
Стабільність предмета визначається як модуль різниці між найкращим та найгіршимсумарним балом за цим предметом.






\par
%end
\newpage
%begin
{\center Варіант  \No \textbf { 29 }\par}
\bigskip
%type = absense
%variant = 28
Обробляються дані семестрового журналу перевірок відвідування занять студентами деканатом (фіксуються відсутні). 

\underline{Записи основного файлу містять поля}: ім’я, по-батькові, прізвище, курс, код групи, аудиторія, вид занять, день тижня, пара, навчальний тиждень, предмет.

\underline{У допоміжному файлі наявні ключі}: кількість студентів, найбільша довжина прізвища.

\underline{Знайти} всіх студентів, що пропускали лекції, але не пропускали практичні та лабораторні заняття. Вивести по кожному з них інформацію:\par
{\noindent}--- на першому рядку:\par
кількість пропусків, прізвище, ім’я, по-батькові;

{\noindent}--- на наступних рядках, починаючи з табуляції, вивести пропуски студента (по одному на рядок):\par
предмет, тиждень, кількість пропусків лекцій, кількість пропусків практичних, кількість пропусків лабораторних, кількість пропусків семінарів, загальна кількість пропусків

{\noindent}у такому сортуванні: предмет, тиждень.

%Студентів сортувати: кількість пропусків(за спаданням), прізвище, ім’я, по-батькові.

\bigskip\textit{Обмеження}

Навчальний тиждень може приймати цілі значення від 1 до 18.
День тижня є цілим від 1 до 5. Пара є цілим від 1 до 4.
Курс є цілим від 1 до 4. Аудиторія є додатнім цілим.
Вид занять може приймати значення:
ЛЕ --- лекція, практ. --- практичне, S --- семінар, Лаб. --- лабораторне.

Решта полів є непорожніми рядками.
Рядки починатися та закінчуватися білими символами не можуть.

Групи кодуються в межах курсу.
Код групи складається не більш ніж з 5 символів.
Може містити літери алфавіту, десяткові цифри та дефіс.

Предмет є рядком, може мати від 4 до 30 символів включно.
Може містити літери алфавіту, десяткові цифри, апостроф, дефіс, пробіл.
Предмети різних курсів вважати різними навіть за однакових назв.

Прізвище, ім’я та по-батькові (за наявності у варіанті)
можуть мати від 4 до 25, 22, 23 відповідно символів включно кожне.
Можуть містити літери алфавіту, апостроф, пробіл та дефіс.
Студенти однозначно ідентифікуються прізвищем, ім’ям та по-батькові (за наявності у варіанті).
Жоден студент не вчиться одночасно на різних курсах чи в різних групах.

Види занять сортуються так: практичне < семінар < лабораторне < лекція.

%\bigskip%\textit{Для деяких варіантів може знадобитись}




\par
%end
\newpage
%begin
{\center Варіант  \No \textbf { 30 }\par}
\bigskip
%type = tasks
%variant = 20
Обробляються результати розв’язування студентами задач на контрольних.
Одиничний факт роз\-в’я\-зу\-вання студентом задачі надалі розуміється як спроба.

\underline{Записи основного файлу містять поля}: по-батькові, умова задачі, відсоток набраних за розв’язання задачі балів, день, ім’я, прізвище, місяць, рік.

\underline{У допоміжному файлі наявні ключі}: середній відсоток з округленням до цілого, довжина найкоротшої умови.

\underline{Знайти} всіх студентів, що не мають спроб з відсотком менше 60\%. Вивести по кожному з них інформацію:\par
{\noindent}--- на першому рядку:\par
найкраща спроба, найгірша спроба, прізвище, ім’я, по-батькові, кількість спроб з відсотком 75\% та більше;

{\noindent}--- на наступних рядках, починаючи з табуляції, вивести їх спроби (по одному на рядок):\par
відсоток набраних за задачу балів, умова, рік

{\noindent}у такому сортуванні: відсоток (за спаданням), рік (за спаданням), місяць, умова.

%Студентів сортувати: найгірша спроба, кількість спроб з відсотком не менше 75\%, прізвище, ім’я, по-батькові.

\bigskip\textit{Обмеження}

Студент міг розв’язувати задачу кілька разів.

Відсоток набраних за розв’язання задачі балів є цілим від 0 до 100.
День, місяць, рік є числовими полями, що зображують дату.
Раніше 2001 року статистика не велася. Рік є 4-значним.
Решта полів є непорожніми рядками.
Рядки починатися та закінчуватися білими символами не можуть.

Прізвище, ім’я та по-батькові (за наявності у варіанті) можуть мати до 24, 29, 24 відповідно символів включно кожне. 
Крім літер алфавіту можуть містити апостроф, дефіс та пробіл.

Умова задачі може мати від 4 до 51 символів включно.
Крім літер алфавіту може містити подвійні лапки, десяткові цифри, пробіл, дефіс,
а також символи . , \_ \{ \} \$ $+$ $*$ {\textbackslash} ( ) \% $[$ $]$ .

Студенти однозначно ідентифікуються прізвищем, ім’ям та по-батькові (за наявності у варіанті).
Задачі однозначно ідентифікуються умовою.

\bigskip\textit{Для деяких варіантів може знадобитись}

Розв'язуваність задачі --- середній відсоток набраних за задачу балів.

Розв'язуваність студента --- середній відсоток набраних студентом балів по всіх його спробах.

Розв'язуваність --- середній відсоток набраних студентами за задачі балів (рахується по всім даним).




\par
%end
\newpage

%begin
{\center Варіант  \No \textbf { 32 }\par}
\bigskip
%type = booking
%variant = 123
Обробляється інформація щодо відвантажених накладних.

\underline{Записи основного файлу містять поля}: назва товару, кількість, ціна, номер товарної позиції у межах накладної, вартість, номер накладної.

\underline{У допоміжному файлі наявні ключі}: середня вартість, найбільша ціна товарної позиції.

\underline{Знайти} 7 найдешевших товарів (з найменшою ціною відвантаження). Вивести по кожному з них інформацію:\par
{\noindent}--- на першому рядку:\par
назва товару, загальна відвантажена кількість, середня ціна відвантаження (з 3 знаками після крапки), найменша ціна відвантаження;

{\noindent}--- на наступних рядках, починаючи з табуляції, вивести для товару всі накладні, де він відвантажувався за ціною, що менша середньої (по одному на рядок):\par
номер накладної, кількість, ціна, вартість

{\noindent}у такому сортуванні: ціна (за спаданням), кількість, номер накладної.

%Товари сортувати: середня ціна відвантаження, назва товару.

\bigskip\textit{Обмеження}

Кількість, ціна та вартість є значеннями дійсного типу.
Решта полів є непорожніми рядками.
Рядки починатися та закінчуватися білими символами не можуть.

Номер накладної складається з 15 символів (цифри, латинські літери).
Назва товару може мати від 2 до 29 символів включно.
Може містити тільки цифри, літери, пробіли, апостроф, дефіс, підкреслення.

Кількість, ціна мають бути додатними.
Кількість вказується з рівно 3 знаками після крапки.
Ціна вказується з рівно 2 знаками після крапки.
Вартість має бути узгоджена з кількістю та ціною та записується рівно з двома десятковими знаками.
(Застосовується округлення.)

Товари однозначно ідентифікуються назвами та одиницями виміру (за наявності).

Одиниця виміру (за наявності у варіанті) може приймати значення: шт., кг, м , кор.

Кожна накладна має свій власний унікальний номер.
У накладній можуть відвантажуватися кілька товарів (не менше 1).
Товарна позиція --- рядок накладної: товар, кількість, ціна, вартість та, за наявності, одиниця виміру.
Товарні позиції кожної накладної нумеруються послідовними натуральними числами, починаючи з 1.
В одній накладній товари можуть повторюватися, 
тобто товарні позиції однієї накладної можуть містити однакові товари 
(при цьому ціна може бути і різною, і однаковою).

\bigskip\textit{Для деяких варіантів може знадобитись}

Недоотриману вигоду розраховувати як суму виразів
{\center (найбільша ціна відвантаження товару – ціна відвантаження товару)*кількість\par}

\bigskip
{\noindent} округлених до 2 знаків після крапки.




\par
%end
\newpage
%begin
{\center Варіант  \No \textbf { 33 }\par}
\bigskip
%type = meteo
%variant = 91
Обробляються дані спостережень за погодою. 

\underline{Записи основного файлу містять поля}: середня денна температура, рік, сила вітру, код метеостанції, відносна вологість, максимальна денна температура, мінімальна денна температура, кількість опадів, місяць, день.

\underline{У допоміжному файлі наявні ключі}: середня силу вітру, загальна кількість спостережень.

\underline{Знайти} 4 дати, по яких фіксувалася найменша мінімальна  денна температура. Вивести по кожному з них інформацію:\par
{\noindent}--- на першому рядку:\par
найменша мінімальна температура, найбільша мінімальна температура, середня вологість, середня сила вітру, середня кількість опадів, місяць, день, рік;

{\noindent}--- на наступних рядках, починаючи з табуляції, вивести для дати спостереження метеостанції, де сила вітру не перевищувала 12 (по одному на рядок):\par
метеостанція, максимальна денна температура, мінімальна денна температура, сила вітру, вологість

{\noindent}у такому сортуванні: сила вітру (за спаданням), метеостанція.

%Дати сортувати: середня вологість, середня кількість опадів, рік, місяць, день.

\bigskip\textit{Обмеження}

Код метеостанції є рядком довжини від 5 до 10 включно.
Може містити літери алфавіту, десяткові цифри, підкреслення. Решта полів числові.

Рік є чотиризначним цілим.
Кількість опадів є дійсним значенням, подається у форматі з 2 знаками після крапки.
Сила вітру є дійсним значенням, подається у форматі з 1 знаком після крапки.
Температури є дійсними значеннями, подаються у форматі з 2 знаками після крапки.
Цілі значення подаються без десяткової крапки.
Для дійсних використовується зображення з фіксованою крапкою.

\textit{Вихідне форматування дійсних}

Усі дійсні значення виводяться у форматі з фіксованою крапкою з тією ж кількістю знаків після крапки, що і у вхідному зображенні.

Якщо у варіанті раніше не було зазначено, то \underline{обчислювані програмою} середні значення виводяться у форматі з 3 знаками після крапки.
%\bigskip%\textit{Для деяких варіантів може знадобитись}




\par
%end
\newpage
%begin
{\center Варіант  \No \textbf { 34 }\par}
\bigskip
%type = meteo
%variant = 112
Обробляються дані спостережень за погодою. 

\underline{Записи основного файлу містять поля}: середня денна температура, кількість опадів, код метеостанції, сила вітру, рік, день, місяць, відносна вологість, максимальна денна температура, мінімальна денна температура.

\underline{У допоміжному файлі наявні ключі}: кількість метеостанцій, середнє значення середньоденної температури.

\underline{Знайти} метеостанції, по яких фіксувалася кількість опадів у два рази менша за середню по всіх даних. Вивести по кожному з них інформацію:\par
{\noindent}--- на першому рядку:\par
найбільша кількість опадів, кількість спостережень, середня мінімальна температура, метеостанція;

{\noindent}--- на наступних рядках, починаючи з табуляції, вивести для них агреговані дані (по одному на рядок):\par
день, середня вологість (за цей день року по всіх місяцях по цій метеостанції), рік

{\noindent}у такому сортуванні: день, рік.

%Метеостанції сортувати: середня мінімальна температура, кількість спостережень (за спаданням), метеостанція.

\bigskip\textit{Обмеження}

Код метеостанції є рядком довжини від 3 до 9 включно.
Може містити літери алфавіту, десяткові цифри, підкреслення. Решта полів числові.

Рік є чотиризначним цілим.
Кількість опадів є дійсним значенням, подається у форматі з 1 знаком після крапки.
Сила вітру є цілим значенням.
Температури є дійсними значеннями, подаються у форматі з 1 знаком після крапки.
Цілі значення подаються без десяткової крапки.
Для дійсних використовується зображення з фіксованою крапкою.

\textit{Вихідне форматування дійсних}

Усі дійсні значення виводяться у форматі з фіксованою крапкою з тією ж кількістю знаків після крапки, що і у вхідному зображенні.

Якщо у варіанті раніше не було зазначено, то \underline{обчислювані програмою} середні значення виводяться у форматі з 1 знаком після крапки.
%\bigskip%\textit{Для деяких варіантів може знадобитись}




\par
%end
\newpage
%begin
{\center Варіант  \No \textbf { 35 }\par}
\bigskip
%type = tasks
%variant = 9
Обробляються результати розв’язування студентами задач на контрольних.
Одиничний факт роз\-в’я\-зу\-вання студентом задачі надалі розуміється як спроба.

\underline{Записи основного файлу містять поля}: день, рік, прізвище, умова задачі, ім’я, відсоток набраних за розв’язання задачі балів, місяць.

\underline{У допоміжному файлі наявні ключі}: кількість задач, кількість записів.

\underline{Знайти} всі задачі, які в кожній спробі були розв’язані не більш ніж на 90\%. Вивести по кожному з них інформацію:\par
{\noindent}--- на першому рядку:\par
умова, кількість спроб, середня розв’язуваність (округлена до цілого);

{\noindent}--- на наступних рядках, починаючи з табуляції, вивести спроби по цих задачах (по одному на рядок):\par
рік, прізвище, ім’я, відсоток набраних балів

{\noindent}у такому сортуванні: прізвище, ім’я, відсоток набраних балів (за спаданням), рік.

%Задачі сортувати: кількість спроб, умова.

\bigskip\textit{Обмеження}

Студент міг розв’язувати задачу кілька разів.

Відсоток набраних за розв’язання задачі балів є цілим від 0 до 100.
День, місяць, рік є числовими полями, що зображують дату.
Раніше 2001 року статистика не велася. Рік є 4-значним.
Решта полів є непорожніми рядками.
Рядки починатися та закінчуватися білими символами не можуть.

Прізвище, ім’я та по-батькові (за наявності у варіанті) можуть мати до 27, 23, 29 відповідно символів включно кожне. 
Крім літер алфавіту можуть містити апостроф, дефіс та пробіл.

Умова задачі може мати від 3 до 61 символів включно.
Крім літер алфавіту може містити подвійні лапки, десяткові цифри, пробіл, дефіс,
а також символи . , \_ \{ \} \$ $+$ $*$ {\textbackslash} ( ) \% $[$ $]$ .

Студенти однозначно ідентифікуються прізвищем, ім’ям та по-батькові (за наявності у варіанті).
Задачі однозначно ідентифікуються умовою.

\bigskip\textit{Для деяких варіантів може знадобитись}

Розв'язуваність задачі --- середній відсоток набраних за задачу балів.

Розв'язуваність студента --- середній відсоток набраних студентом балів по всіх його спробах.

Розв'язуваність --- середній відсоток набраних студентами за задачі балів (рахується по всім даним).




\par
%end
\newpage
%begin
{\center Варіант  \No \textbf { 36 }\par}
\bigskip
%type = absense
%variant = 34
Обробляються дані семестрового журналу перевірок відвідування занять студентами деканатом (фіксуються відсутні). 

\underline{Записи основного файлу містять поля}: предмет, вид занять, аудиторія, прізвище, навчальний тиждень, день тижня, пара, ім’я, курс.

\underline{У допоміжному файлі наявні ключі}: сума пар, найбільша довжина назви предмета.

\underline{Знайти} всіх студентів, що пропустили найбільше практичних занять. Вивести по кожному з них інформацію:\par
{\noindent}--- на першому рядку:\par
кількість пропусків практичних, прізвище, ім’я, кількість пропущених пар;

{\noindent}--- на наступних рядках, починаючи з табуляції, вивести пропуски студента (по одному на рядок):\par
предмет, тиждень, кількість пропусків, вид занять

{\noindent}у такому сортуванні: предмет, тиждень, вид занять, кількість пропусків.

%Студентів сортувати: кількість пропусків практичних (за спаданням), прізвище, ім’я.

\bigskip\textit{Обмеження}

Навчальний тиждень може приймати цілі значення від 1 до 17.
День тижня є цілим від 1 до 5. Пара є цілим від 1 до 4.
Курс є цілим від 1 до 4. Аудиторія є додатнім цілим.
Вид занять може приймати значення:
Лекц. --- лекція, p --- практичне, sem --- семінар, Lab --- лабораторне.

Решта полів є непорожніми рядками.
Рядки починатися та закінчуватися білими символами не можуть.

Групи кодуються в межах курсу.
Код групи складається не більш ніж з 4 символів.
Може містити літери алфавіту, десяткові цифри та дефіс.

Предмет є рядком, може мати від 6 до 35 символів включно.
Може містити літери алфавіту, десяткові цифри, апостроф, дефіс, пробіл.
Предмети різних курсів вважати різними навіть за однакових назв.

Прізвище, ім’я та по-батькові (за наявності у варіанті)
можуть мати від 6 до 28, 21, 26 відповідно символів включно кожне.
Можуть містити літери алфавіту, апостроф, пробіл та дефіс.
Студенти однозначно ідентифікуються прізвищем, ім’ям та по-батькові (за наявності у варіанті).
Жоден студент не вчиться одночасно на різних курсах чи в різних групах.

Види занять сортуються так: практичне < лекція < семінар < лабораторне.

%\bigskip%\textit{Для деяких варіантів може знадобитись}




\par
%end
\newpage
%begin
{\center Варіант  \No \textbf { 37 }\par}
\bigskip
%type = tasks
%variant = 19
Обробляються результати розв’язування студентами задач на контрольних.
Одиничний факт роз\-в’я\-зу\-вання студентом задачі надалі розуміється як спроба.

\underline{Записи основного файлу містять поля}: відсоток набраних за розв’язання задачі балів, умова задачі, ім’я, по-батькові, рік, прізвище.

\underline{У допоміжному файлі наявні ключі}: найбільший відсоток, найменший відсоток, загальна кількість записів.

\underline{Знайти} всіх студентів, у яких більше половини спроб отримали не менше 75\%. Вивести по кожному з них інформацію:\par
{\noindent}--- на першому рядку:\par
прізвище, ім’я, по-батькові, найкраща спроба, середня розв’язуваність (округлена до цілого), кількість спроб з відсотком не нижче 75\%;

{\noindent}--- на наступних рядках, починаючи з табуляції, вивести їх спроби (по одному на рядок):\par
відсоток набраних за задачу балів, умова, рік

{\noindent}у такому сортуванні: відсоток (за спаданням), рік (за спаданням), умова.

%Студентів сортувати: найкраща спроба (за зростанням), середня розв’язуваність, прізвище, ім’я, по-батькові.

\bigskip\textit{Обмеження}

Студент міг розв’язувати задачу кілька разів.

Відсоток набраних за розв’язання задачі балів є цілим від 0 до 100.
День, місяць, рік є числовими полями, що зображують дату.
Раніше 2001 року статистика не велася. Рік є 4-значним.
Решта полів є непорожніми рядками.
Рядки починатися та закінчуватися білими символами не можуть.

Прізвище, ім’я та по-батькові (за наявності у варіанті) можуть мати до 20, 24, 25 відповідно символів включно кожне. 
Крім літер алфавіту можуть містити апостроф, дефіс та пробіл.

Умова задачі може мати від 5 до 64 символів включно.
Крім літер алфавіту може містити подвійні лапки, десяткові цифри, пробіл, дефіс,
а також символи . , \_ \{ \} \$ $+$ $*$ {\textbackslash} ( ) \% $[$ $]$ .

Студенти однозначно ідентифікуються прізвищем, ім’ям та по-батькові (за наявності у варіанті).
Задачі однозначно ідентифікуються умовою.

\bigskip\textit{Для деяких варіантів може знадобитись}

Розв'язуваність задачі --- середній відсоток набраних за задачу балів.

Розв'язуваність студента --- середній відсоток набраних студентом балів по всіх його спробах.

Розв'язуваність --- середній відсоток набраних студентами за задачі балів (рахується по всім даним).




\par
%end
\newpage
%begin
{\center Варіант  \No \textbf { 38 }\par}
\bigskip
%type =     booking
%variant = 116
Обробляється інформація щодо відвантажених накладних.

\underline{Записи основного файлу містять поля}: номер накладної, ціна, вартість, назва товару, кількість, номер товарної позиції у межах накладної.

\underline{У допоміжному файлі наявні ключі}: середня вартість, кількість записів.

\underline{Знайти} 4 найменш популярні товари (ті, що сумарно відвантажувалися в найменшій кількості). Вивести по кожному з них інформацію:\par
{\noindent}--- на першому рядку:\par
назва товару, загальна відвантажена кількість, найбільша вартість відвантаження (з 2 знаками після крапки);

{\noindent}--- на наступних рядках, починаючи з табуляції, вивести для товару всі накладні, де його було відвантажено більше 2 одиниць (по одному на рядок):\par
номер накладної, кількість, ціна, вартість

{\noindent}у такому сортуванні: номер накладної, кількість (за спаданням).

%Накладні сортувати: загальна відвантажена кількість (за спаданням), назва.

\bigskip\textit{Обмеження}

Кількість, ціна та вартість є значеннями дійсного типу.
Решта полів є непорожніми рядками.
Рядки починатися та закінчуватися білими символами не можуть.

Номер накладної складається з 14 символів (цифри, латинські літери).
Назва товару може мати від 4 до 24 символів включно.
Може містити тільки цифри, літери, пробіли, апостроф, дефіс, підкреслення.

Кількість, ціна мають бути додатними.
Кількість вказується з рівно 1 знаками після крапки.
Ціна вказується з рівно 3 знаками після крапки.
Вартість має бути узгоджена з кількістю та ціною та записується рівно з двома десятковими знаками.
(Застосовується округлення.)

Товари однозначно ідентифікуються назвами та одиницями виміру (за наявності).

Одиниця виміру (за наявності у варіанті) може приймати значення: шт., кг, м , кор.

Кожна накладна має свій власний унікальний номер.
У накладній можуть відвантажуватися кілька товарів (не менше 1).
Товарна позиція --- рядок накладної: товар, кількість, ціна, вартість та, за наявності, одиниця виміру.
Товарні позиції кожної накладної нумеруються послідовними натуральними числами, починаючи з 1.
В одній накладній товари можуть повторюватися, 
тобто товарні позиції однієї накладної можуть містити однакові товари 
(при цьому ціна може бути і різною, і однаковою).

\bigskip\textit{Для деяких варіантів може знадобитись}

Недоотриману вигоду розраховувати як суму виразів
{\center (найбільша ціна відвантаження товару – ціна відвантаження товару)*кількість\par}

\bigskip
{\noindent} округлених до 2 знаків після крапки.




\par
%end
\newpage
%begin
{\center Варіант  \No \textbf { 39 }\par}
\bigskip
%type = booking
%variant = 119
Обробляється інформація щодо відвантажених накладних.

\underline{Записи основного файлу містять поля}: ціна, кількість, номер товарної позиції у межах накладної, назва товару, номер накладної, вартість.

\underline{У допоміжному файлі наявні ключі}: кількість різних товарів, кількість накладних.

\underline{Знайти} 4 найбільш популярні товари (ті, що сумарно відвантажувалися на найбільшу суму). Вивести по кожному з них інформацію:\par
{\noindent}--- на першому рядку:\par
назва товару, загальна сума відвантаження (з 2 знаками після крапки), середня ціна відвантаження (з 2 знаками після крапки);

{\noindent}--- на наступних рядках, починаючи з табуляції, вивести для товару всі накладні, де він відвантажувався в найбільшій кількості (по одному на рядок):\par
номер накладної, вартість, ціна, кількість

{\noindent}у такому сортуванні: номер накладної, ціна (за спаданням).

%Товари сортувати: назва товару.

\bigskip\textit{Обмеження}

Кількість, ціна та вартість є значеннями дійсного типу.
Решта полів є непорожніми рядками.
Рядки починатися та закінчуватися білими символами не можуть.

Номер накладної складається з 10 символів (цифри, латинські літери).
Назва товару може мати від 1 до 28 символів включно.
Може містити тільки цифри, літери, пробіли, апостроф, дефіс, підкреслення.

Кількість, ціна мають бути додатними.
Кількість вказується з рівно 3 знаками після крапки.
Ціна вказується з рівно 4 знаками після крапки.
Вартість має бути узгоджена з кількістю та ціною та записується рівно з двома десятковими знаками.
(Застосовується округлення.)

Товари однозначно ідентифікуються назвами та одиницями виміру (за наявності).

Одиниця виміру (за наявності у варіанті) може приймати значення: шт., кг, м , кор.

Кожна накладна має свій власний унікальний номер.
У накладній можуть відвантажуватися кілька товарів (не менше 1).
Товарна позиція --- рядок накладної: товар, кількість, ціна, вартість та, за наявності, одиниця виміру.
Товарні позиції кожної накладної нумеруються послідовними натуральними числами, починаючи з 1.
В одній накладній товари можуть повторюватися, 
тобто товарні позиції однієї накладної можуть містити однакові товари 
(при цьому ціна може бути і різною, і однаковою).

\bigskip\textit{Для деяких варіантів може знадобитись}

Недоотриману вигоду розраховувати як суму виразів
{\center (найбільша ціна відвантаження товару – ціна відвантаження товару)*кількість\par}

\bigskip
{\noindent} округлених до 2 знаків після крапки.




\par
%end
\newpage
%begin
{\center Варіант  \No \textbf { 40 }\par}
\bigskip
%type = booking
%variant = 121
Обробляється інформація щодо відвантажених накладних.

\underline{Записи основного файлу містять поля}: вартість, номер накладної, одиниця виміру, назва товару, кількість, номер товарної позиції у межах накладної, ціна.

\underline{У допоміжному файлі наявні ключі}: сума цін, найменша ціна товарної позиції.

\underline{Знайти} 5 найбільш популярних товарів (ті, що відвантажувалися в найбільшій кількості накладних). Вивести по кожному з них інформацію:\par
{\noindent}--- на першому рядку:\par
кількість накладних, в яких відвантажувалися товари, назва товару, одиниця виміру, сума відвантаження (з 2 знаками після крапки), найменша відпускна ціна товару (з 4 знаками після крапки);

{\noindent}--- на наступних рядках, починаючи з табуляції, вивести для товару всі накладні, де він відвантажувався менше ніж в кількості 2 одиниць (по одному на рядок):\par
номер накладної, кількість, ціна

{\noindent}у такому сортуванні: кількість (за спаданням), номер накладної.

%Товари сортувати: кількість накладних з товаром (за спаданням), назва товару, одиниця виміру.

\bigskip\textit{Обмеження}

Кількість, ціна та вартість є значеннями дійсного типу.
Решта полів є непорожніми рядками.
Рядки починатися та закінчуватися білими символами не можуть.

Номер накладної складається з 11 символів (цифри, латинські літери).
Назва товару може мати від 3 до 30 символів включно.
Може містити тільки цифри, літери, пробіли, апостроф, дефіс, підкреслення.

Кількість, ціна мають бути додатними.
Кількість вказується з рівно 4 знаками після крапки.
Ціна вказується з рівно 2 знаками після крапки.
Вартість має бути узгоджена з кількістю та ціною та записується рівно з двома десятковими знаками.
(Застосовується округлення.)

Товари однозначно ідентифікуються назвами та одиницями виміру (за наявності).

Одиниця виміру (за наявності у варіанті) може приймати значення: шт., кг, м , кор.

Кожна накладна має свій власний унікальний номер.
У накладній можуть відвантажуватися кілька товарів (не менше 1).
Товарна позиція --- рядок накладної: товар, кількість, ціна, вартість та, за наявності, одиниця виміру.
Товарні позиції кожної накладної нумеруються послідовними натуральними числами, починаючи з 1.
В одній накладній товари можуть повторюватися, 
тобто товарні позиції однієї накладної можуть містити однакові товари 
(при цьому ціна може бути і різною, і однаковою).

\bigskip\textit{Для деяких варіантів може знадобитись}

Недоотриману вигоду розраховувати як суму виразів
{\center (найбільша ціна відвантаження товару – ціна відвантаження товару)*кількість\par}

\bigskip
{\noindent} округлених до 2 знаків після крапки.




\par
%end
\newpage
%begin
{\center Варіант  \No \textbf { 41 }\par}
\bigskip
%type = meteo
%variant = 85
Обробляються дані спостережень за погодою. 

\underline{Записи основного файлу містять поля}: код метеостанції, середня денна температура, рік, відносна вологість, мінімальна денна температура, сила вітру, місяць, день, кількість опадів, максимальна денна температура.

\underline{У допоміжному файлі наявні ключі}: рік початку спостережень, кількість записів по 13-их числах.

\underline{Знайти} 7 дат, по яких фіксувалася найбільша сила вітру. Вивести по кожному з них інформацію:\par
{\noindent}--- на першому рядку:\par
найбільша сила вітру, найменша сила вітру, сумарна кількість опадів, місяць, рік, день, середня вологість;

{\noindent}--- на наступних рядках, починаючи з табуляції, вивести для дати спостереження метеостанції, де вологість була менша за середню вологість за дату (по одному на рядок):\par
вологість, метеостанція, сила вітру

{\noindent}у такому сортуванні: вологість, сила вітру (за спаданням), метеостанція.

%Дати сортувати: рік, місяць, день, сумарна кількість опадів.

\bigskip\textit{Обмеження}

Код метеостанції є рядком довжини від 4 до 8 включно.
Може містити літери алфавіту, десяткові цифри, підкреслення. Решта полів числові.

Рік є чотиризначним цілим.
Кількість опадів є цілим значенням.
Сила вітру є дійсним значенням, подається у форматі з 2 знаками після крапки.
Температури є дійсними значеннями, подаються у форматі з 3 знаками після крапки.
Цілі значення подаються без десяткової крапки.
Для дійсних використовується зображення з фіксованою крапкою.

\textit{Вихідне форматування дійсних}

Усі дійсні значення виводяться у форматі з фіксованою крапкою з тією ж кількістю знаків після крапки, що і у вхідному зображенні.

Якщо у варіанті раніше не було зазначено, то \underline{обчислювані програмою} середні значення виводяться у форматі з 2 знаками після крапки.
%\bigskip%\textit{Для деяких варіантів може знадобитись}




\par
%end
\newpage
%begin
{\center Варіант  \No \textbf { 42 }\par}
\bigskip
%type =     meteo
%variant = 82
Обробляються дані спостережень за погодою. 

\underline{Записи основного файлу містять поля}: максимальна денна температура, середня денна температура, день, місяць, мінімальна денна температура, кількість опадів, сила вітру, відносна вологість, рік, код метеостанції.

\underline{У допоміжному файлі наявні ключі}: кількість записів, найбільша сила вітру.

\underline{Знайти} 5 дат, за які випала найбільша кількість опадів (сумарно по всіх метеостанціях). Вивести по кожному з них інформацію:\par
{\noindent}--- на першому рядку:\par
рік, місяць, день, кількість опадів, середня середньоденна температура;

{\noindent}--- на наступних рядках, починаючи з табуляції, вивести для дати спостереження метеостанцій, де максимальна денна температура була більша за середню (по всіх метеостанціях) середньоденну за цю дату (по одному на рядок):\par
метеостанція, середня денна температура, максимальна денна температура, мінімальна денна температура, сила вітру

{\noindent}у такому сортуванні: сила вітру, метеостанція.

%Дати сортувати: кількість опадів, середня середньоденна температура, рік, місяць, день.

\bigskip\textit{Обмеження}

Код метеостанції є рядком довжини від 3 до 7 включно.
Може містити літери алфавіту, десяткові цифри, підкреслення. Решта полів числові.

Рік є чотиризначним цілим.
Кількість опадів є дійсним значенням, подається у форматі з 2 знаками після крапки.
Сила вітру є дійсним значенням, подається у форматі з 2 знаками після крапки.
Температури є дійсними значеннями, подаються у форматі з 2 знаками після крапки.
Цілі значення подаються без десяткової крапки.
Для дійсних використовується зображення з фіксованою крапкою.

\textit{Вихідне форматування дійсних}

Усі дійсні значення виводяться у форматі з фіксованою крапкою з тією ж кількістю знаків після крапки, що і у вхідному зображенні.

Якщо у варіанті раніше не було зазначено, то \underline{обчислювані програмою} середні значення виводяться у форматі з 2 знаками після крапки.
%\bigskip%\textit{Для деяких варіантів може знадобитись}




\par
%end
\newpage
%begin
{\center Варіант  \No \textbf { 43 }\par}
\bigskip
%type = booking
%variant = 118
Обробляється інформація щодо відвантажених накладних.

\underline{Записи основного файлу містять поля}: вартість, номер товарної позиції у межах накладної, ціна, номер накладної, одиниця виміру, кількість, назва товару.

\underline{У допоміжному файлі наявні ключі}: найбільша ціна товарної позиції, сума вартостей.

\underline{Знайти} 3 найменш популярні товари (ті, що сумарно відвантажувалися на найменшу суму). Вивести по кожному з них інформацію:\par
{\noindent}--- на першому рядку:\par
найбільша відвантажена кількість в одній накладній (з 1 знаком після крапки),  назва товару, одиниця виміру;

{\noindent}--- на наступних рядках, починаючи з табуляції, вивести для товару всі накладні, де він відвантажувався в найбільшій кількості (по одному на рядок):\par
вартість, ціна, кількість, номер накладної

{\noindent}у такому сортуванні: ціна (за спаданням), номер накладної.

%Товари сортувати: загальна відвантажена кількість (за спаданням), назва товару, одиниця виміру.

\bigskip\textit{Обмеження}

Кількість, ціна та вартість є значеннями дійсного типу.
Решта полів є непорожніми рядками.
Рядки починатися та закінчуватися білими символами не можуть.

Номер накладної складається з 12 символів (цифри, латинські літери).
Назва товару може мати від 4 до 27 символів включно.
Може містити тільки цифри, літери, пробіли, апостроф, дефіс, підкреслення.

Кількість, ціна мають бути додатними.
Кількість вказується з рівно 1 знаками після крапки.
Ціна вказується з рівно 4 знаками після крапки.
Вартість має бути узгоджена з кількістю та ціною та записується рівно з двома десятковими знаками.
(Застосовується округлення.)

Товари однозначно ідентифікуються назвами та одиницями виміру (за наявності).

Одиниця виміру (за наявності у варіанті) може приймати значення: шт., кг, м , кор.

Кожна накладна має свій власний унікальний номер.
У накладній можуть відвантажуватися кілька товарів (не менше 1).
Товарна позиція --- рядок накладної: товар, кількість, ціна, вартість та, за наявності, одиниця виміру.
Товарні позиції кожної накладної нумеруються послідовними натуральними числами, починаючи з 1.
В одній накладній товари можуть повторюватися, 
тобто товарні позиції однієї накладної можуть містити однакові товари 
(при цьому ціна може бути і різною, і однаковою).

\bigskip\textit{Для деяких варіантів може знадобитись}

Недоотриману вигоду розраховувати як суму виразів
{\center (найбільша ціна відвантаження товару – ціна відвантаження товару)*кількість\par}

\bigskip
{\noindent} округлених до 2 знаків після крапки.




\par
%end
\newpage
%begin
{\center Варіант  \No \textbf { 44 }\par}
\bigskip
%type = booking
%variant = 122
Обробляється інформація щодо відвантажених накладних.

\underline{Записи основного файлу містять поля}: назва товару, кількість, ціна, номер товарної позиції у межах накладної, номер накладної, вартість, одиниця виміру.

\underline{У допоміжному файлі наявні ключі}: кількість накладних, кількість записів.

\underline{Знайти} 5 найдорожчих товарів (з найбільшою ціною відвантаження). Вивести по кожному з них інформацію:\par
{\noindent}--- на першому рядку:\par
найбільша ціна відвантаження, найменша ціна відвантаження, назва товару, одиниця виміру, загальна відвантажена кількість;

{\noindent}--- на наступних рядках, починаючи з табуляції, вивести для товару всі накладні, де він відвантажувався за ціною, що відрізняється від найбільшої (по одному на рядок):\par
номер накладної, кількість, ціна, недоотримана вигода (з 1 знаком після крапки)

{\noindent}у такому сортуванні: номер накладної.

%Товари сортувати: різниця максимальної та мінімальної відпускних цін за спаданням, назва товару, одиниця виміру.

\bigskip\textit{Обмеження}

Кількість, ціна та вартість є значеннями дійсного типу.
Решта полів є непорожніми рядками.
Рядки починатися та закінчуватися білими символами не можуть.

Номер накладної складається з 13 символів (цифри, латинські літери).
Назва товару може мати від 2 до 23 символів включно.
Може містити тільки цифри, літери, пробіли, апостроф, дефіс, підкреслення.

Кількість, ціна мають бути додатними.
Кількість вказується з рівно 2 знаками після крапки.
Ціна вказується з рівно 3 знаками після крапки.
Вартість має бути узгоджена з кількістю та ціною та записується рівно з двома десятковими знаками.
(Застосовується округлення.)

Товари однозначно ідентифікуються назвами та одиницями виміру (за наявності).

Одиниця виміру (за наявності у варіанті) може приймати значення: шт., кг, м , кор.

Кожна накладна має свій власний унікальний номер.
У накладній можуть відвантажуватися кілька товарів (не менше 1).
Товарна позиція --- рядок накладної: товар, кількість, ціна, вартість та, за наявності, одиниця виміру.
Товарні позиції кожної накладної нумеруються послідовними натуральними числами, починаючи з 1.
В одній накладній товари можуть повторюватися, 
тобто товарні позиції однієї накладної можуть містити однакові товари 
(при цьому ціна може бути і різною, і однаковою).

\bigskip\textit{Для деяких варіантів може знадобитись}

Недоотриману вигоду розраховувати як суму виразів
{\center (найбільша ціна відвантаження товару – ціна відвантаження товару)*кількість\par}

\bigskip
{\noindent} округлених до 2 знаків після крапки.




\par
%end
\newpage
%begin
{\center Варіант  \No \textbf { 45 }\par}
\bigskip
%type = tasks
%variant = 12
Обробляються результати розв’язування студентами задач на контрольних.
Одиничний факт роз\-в’я\-зу\-вання студентом задачі надалі розуміється як спроба.

\underline{Записи основного файлу містять поля}: місяць, день, прізвище, умова задачі, ім’я, по-батькові, відсоток набраних за розв’язання задачі балів, рік.

\underline{У допоміжному файлі наявні ключі}: найменший відсоток, найбільший відсоток, кількість різних задач.

\underline{Знайти} всі задачі, за якими не було жодної спроби краще 50\%. Вивести по кожному з них інформацію:\par
{\noindent}--- на першому рядку:\par
умова, кількість спроб кращих за 30\%, кількість спроб, відсоток найкращої спроби;

{\noindent}--- на наступних рядках, починаючи з табуляції, вивести спроби по цих задачах (по одному на рядок):\par
прізвище, ім’я, по-батькові, відсоток набраних балів, рік

{\noindent}у такому сортуванні: відсоток набраних балів (спочатку нулеві, далі інші за зростанням), рік, прізвище, ім’я, по-батькові.

%Задачі сортувати: відсоток найкращої спроби, кількість спроб, умова.

\bigskip\textit{Обмеження}

Студент міг розв’язувати задачу кілька разів.

Відсоток набраних за розв’язання задачі балів є цілим від 0 до 100.
День, місяць, рік є числовими полями, що зображують дату.
Раніше 2001 року статистика не велася. Рік є 4-значним.
Решта полів є непорожніми рядками.
Рядки починатися та закінчуватися білими символами не можуть.

Прізвище, ім’я та по-батькові (за наявності у варіанті) можуть мати до 21, 22, 28 відповідно символів включно кожне. 
Крім літер алфавіту можуть містити апостроф, дефіс та пробіл.

Умова задачі може мати від 3 до 62 символів включно.
Крім літер алфавіту може містити подвійні лапки, десяткові цифри, пробіл, дефіс,
а також символи . , \_ \{ \} \$ $+$ $*$ {\textbackslash} ( ) \% $[$ $]$ .

Студенти однозначно ідентифікуються прізвищем, ім’ям та по-батькові (за наявності у варіанті).
Задачі однозначно ідентифікуються умовою.

\bigskip\textit{Для деяких варіантів може знадобитись}

Розв'язуваність задачі --- середній відсоток набраних за задачу балів.

Розв'язуваність студента --- середній відсоток набраних студентом балів по всіх його спробах.

Розв'язуваність --- середній відсоток набраних студентами за задачі балів (рахується по всім даним).




\par
%end
\newpage
%begin
{\center Варіант  \No \textbf { 46 }\par}
\bigskip
%type = booking
%variant = -5
Обробляється інформація щодо відвантажених накладних.

\underline{Записи основного файлу містять поля}: номер накладної, номер товарної позиції у межах накладної, назва товару, ціна, кількість, вартість.

\underline{У допоміжному файлі наявні ключі}: загальна кількість записів, сума кількостей.

\underline{Знайти} всі накладні, що містять повторювані товари. Вивести по кожному з них інформацію:\par
{\noindent}--- на першому рядку:\par
номер накладної, кількість товару, кількість товарних позицій, загальна сума;

{\noindent}--- на наступних рядках, починаючи з табуляції, вивести для накладної всі її товарні позиції, що містять повторювані товари (по одному на рядок):\par
назва товару, ціна, кількість, вартість, номер товарної позиції

{\noindent}у такому сортуванні: назва товару, кількість, ціна.

%Накладні сортувати: назва товару, ціна (за спаданням), кількість, номер накладної.

\bigskip\textit{Обмеження}

Кількість, ціна та вартість є значеннями дійсного типу.
Решта полів є непорожніми рядками.
Рядки починатися та закінчуватися білими символами не можуть.

Номер накладної складається з 11 символів (цифри, латинські літери).
Назва товару може мати від 2 до 26 символів включно.
Може містити тільки цифри, літери, пробіли, апостроф, дефіс, підкреслення.

Кількість, ціна мають бути додатними.
Кількість вказується з рівно 3 знаками після крапки.
Ціна вказується з рівно 2 знаками після крапки.
Вартість має бути узгоджена з кількістю та ціною та записується рівно з двома десятковими знаками.
(Застосовується округлення.)

Товари однозначно ідентифікуються назвами та одиницями виміру (за наявності).

Одиниця виміру (за наявності у варіанті) може приймати значення: шт., кг, м , кор.

Кожна накладна має свій власний унікальний номер.
У накладній можуть відвантажуватися кілька товарів (не менше 1).
Товарна позиція --- рядок накладної: товар, кількість, ціна, вартість та, за наявності, одиниця виміру.
Товарні позиції кожної накладної нумеруються послідовними натуральними числами, починаючи з 1.
В одній накладній товари можуть повторюватися, 
тобто товарні позиції однієї накладної можуть містити однакові товари 
(при цьому ціна може бути і різною, і однаковою).

\bigskip\textit{Для деяких варіантів може знадобитись}

Недоотриману вигоду розраховувати як суму виразів
{\center (найбільша ціна відвантаження товару – ціна відвантаження товару)*кількість\par}

\bigskip
{\noindent} округлених до 2 знаків після крапки.




\par
%end
\newpage
%begin
{\center Варіант  \No \textbf { 47 }\par}
\bigskip
%type = session
%variant = 60
Обробляються результати складання студентами екзаменів зимової сесії.

\underline{Записи основного файлу містять поля}: номер залікової книжки, набрані впродовж семестру бали, сумарна оцінка в балах за 100-бальною системою, набрані на екзамені бали, прізвище, ім’я, код групи, предмет, оцінка за державною шкалою.

\underline{У допоміжному файлі наявні ключі}: довжина найкоротшого прізвища, загальна кількість записів.

\underline{Знайти} всіх студентів, що по кожному предмету отримали сумарний бал рівно 60. Вивести по кожному з них інформацію:\par
{\noindent}--- на першому рядку:\par
номер залікової книжки, прізвище, ім’я, найменший бал на екзамені, середній бал в семестрі (округлений до цілого);

{\noindent}--- на наступних рядках, починаючи з табуляції, вивести для студента всі його результати (по одному на рядок):\par
предмет, бали на екзамені, бали в семестрі

{\noindent}у такому сортуванні: бали в семестрі, предмет.

%Студентів сортувати: середній бал в семестрі (після округлення), номер залікової книжки.

\bigskip\textit{Обмеження}

Студент не може складати предмет більше одного разу.

Бали та оцінка є значеннями цілих числових типів. Решта полів є непорожніми рядками.
Рядки розпочинатись та закінчуватись білими символами не можуть.

Оцінка за державною шкалою може приймати значення: 2-5~--- як зазвичай, 0~--- не з’явився,$-1$~--- не допущений. 
Оцінки 3-5 вважаються задовільними, решта~--- незадовільними. 
Бали є цілими та невід’ємними.
Набрані впродовж семестру бали не можуть перевищувати 60. 
Набрані на екзамені бали не можуть перевищувати 40. 
Набрані на екзамені бали повинні бути не менше 24 або дорівнювати 0(в останньому випадку оцінка не може бути задовільною). 
Сумарна оцінка в балах є сумою балів, набраних у семестрі та на екзамені.
Сумарна оцінка в балах та за державною шкалою мають бути узгоджені між собою.

Прізвище, ім’я та по-батькові (якщо наявне у варіанті) можуть мати до 26, 30, 20 відповідно символів включно кожне. 
Крім літер алфавіту можуть містити апостроф, дефіс та пробіл.

Предмет може мати від 5 до 50 символів включно. 
Крім літер алфавіту може містити подвійні лапки, пробіл та дефіс.

Код групи є непорожнім та складається не більш ніж з 6 літер. 
Крім літер алфавіту може містити десяткові цифри та дефіс.


Номер залікової книжки однозначно ідентифікує студента та складається з 6 десяткових цифр. 

\bigskip\textit{Для деяких варіантів може знадобитись}

Середній бал/оцінка за предметом визначається як середнє арифметичнесумарних балів/державних оцінок за цим предметом. 
Середній бал/оцінка студента визначається як середнє арифметичне сумарних балів/державних оцінокцього студента. 
За обчислення середньої оцінки недопущені та ті, що не з’явилися, рахуються як 2.
У решті випадків недопуски підсумовуються як $-1$, а «не з’явився» як 0.
Якщо студент не гірше ніж задовільно склав усі предмети, 
то його рейтинг за балами/оцінкою визначається як середній бал/оцінка; інакше дорівнює 0.

Стабільність студента визначається як модуль різниці між його кращим та найгіршим сумарним балом.
Стабільність предмета визначається як модуль різниці між найкращим та найгіршимсумарним балом за цим предметом.






\par
%end
\newpage
%begin
{\center Варіант  \No \textbf { 48 }\par}
\bigskip
%type = absense
%variant = 48
Обробляються дані семестрового журналу перевірок відвідування занять студентами деканатом (фіксуються відсутні). 

\underline{Записи основного файлу містять поля}: прізвище, навчальний тиждень, аудиторія, предмет, вид занять, ім’я, курс, код групи, день тижня, пара.

\underline{У допоміжному файлі наявні ключі}: кількість пропусків перших пар, загальна кількість записів.

\underline{Знайти} всі предмети, з яких було пропущено найбільше семінарів. Вивести по кожному з них інформацію:\par
{\noindent}--- на першому рядку:\par
курс, предмет, кількість пропусків семінарів, кількість пропусків лекцій, кількість пропусків практичних, кількість пропусків лабораторних занять, кількість пропусків;

{\noindent}--- на наступних рядках, починаючи з табуляції, вивести пропуски семінарів з предмету (по одному на рядок):\par
прізвище, ім’я, по-батькові, тиждень, день, аудиторія, пара

{\noindent}у такому сортуванні: прізвище, ім’я, по-батькові, тиждень, день, пара, аудиторія.

%Предмети сортувати: кількість пропусків (за спаданням), курс, предмет.

\bigskip\textit{Обмеження}

Навчальний тиждень може приймати цілі значення від 1 до 16.
День тижня є цілим від 1 до 5. Пара є цілим від 1 до 4.
Курс є цілим від 1 до 4. Аудиторія є додатнім цілим.
Вид занять може приймати значення:
Le --- лекція, pr --- практичне, s --- семінар, лаб. --- лабораторне.

Решта полів є непорожніми рядками.
Рядки починатися та закінчуватися білими символами не можуть.

Групи кодуються в межах курсу.
Код групи складається не більш ніж з 5 символів.
Може містити літери алфавіту, десяткові цифри та дефіс.

Предмет є рядком, може мати від 3 до 33 символів включно.
Може містити літери алфавіту, десяткові цифри, апостроф, дефіс, пробіл.
Предмети різних курсів вважати різними навіть за однакових назв.

Прізвище, ім’я та по-батькові (за наявності у варіанті)
можуть мати від 3 до 24, 23, 20 відповідно символів включно кожне.
Можуть містити літери алфавіту, апостроф, пробіл та дефіс.
Студенти однозначно ідентифікуються прізвищем, ім’ям та по-батькові (за наявності у варіанті).
Жоден студент не вчиться одночасно на різних курсах чи в різних групах.

Види занять сортуються так: лабораторне < практичне < лекція < семінар.

%\bigskip%\textit{Для деяких варіантів може знадобитись}




\par
%end
\newpage
%begin
{\center Варіант  \No \textbf { 49 }\par}
\bigskip
%type = booking
%variant = 124
Обробляється інформація щодо відвантажених накладних.

\underline{Записи основного файлу містять поля}: номер накладної, кількість, назва товару, вартість, ціна, номер товарної позиції у межах накладної.

\underline{У допоміжному файлі наявні ключі}: найбільша ціна товарної позиції, сума вартостей.

\underline{Знайти} 5 найдорожчих товарів (з найбільшою середньою ціною відвантаження). Вивести по кожному з них інформацію:\par
{\noindent}--- на першому рядку:\par
назва товару, загальна відвантажена кількість, середня ціна відвантаження (з 2 знаками після крапки), найменша ціна відвантаження,  найбільша ціна відвантаження;

{\noindent}--- на наступних рядках, починаючи з табуляції, вивести для товару всі накладні, де він відвантажувався за ненайбільшою своєю ціною (по одному на рядок):\par
ціна, кількість, ціна, вартість, номер накладної

{\noindent}у такому сортуванні: вартість (за спаданням), ціна(за спаданням), номер накладної.

%Товари сортувати: середня ціна відвантаження (за спаданням), назва товару.

\bigskip\textit{Обмеження}

Кількість, ціна та вартість є значеннями дійсного типу.
Решта полів є непорожніми рядками.
Рядки починатися та закінчуватися білими символами не можуть.

Номер накладної складається з 14 символів (цифри, латинські літери).
Назва товару може мати від 1 до 20 символів включно.
Може містити тільки цифри, літери, пробіли, апостроф, дефіс, підкреслення.

Кількість, ціна мають бути додатними.
Кількість вказується з рівно 2 знаками після крапки.
Ціна вказується з рівно 4 знаками після крапки.
Вартість має бути узгоджена з кількістю та ціною та записується рівно з двома десятковими знаками.
(Застосовується округлення.)

Товари однозначно ідентифікуються назвами та одиницями виміру (за наявності).

Одиниця виміру (за наявності у варіанті) може приймати значення: шт., кг, м , кор.

Кожна накладна має свій власний унікальний номер.
У накладній можуть відвантажуватися кілька товарів (не менше 1).
Товарна позиція --- рядок накладної: товар, кількість, ціна, вартість та, за наявності, одиниця виміру.
Товарні позиції кожної накладної нумеруються послідовними натуральними числами, починаючи з 1.
В одній накладній товари можуть повторюватися, 
тобто товарні позиції однієї накладної можуть містити однакові товари 
(при цьому ціна може бути і різною, і однаковою).

\bigskip\textit{Для деяких варіантів може знадобитись}

Недоотриману вигоду розраховувати як суму виразів
{\center (найбільша ціна відвантаження товару – ціна відвантаження товару)*кількість\par}

\bigskip
{\noindent} округлених до 2 знаків після крапки.




\par
%end
\newpage

%begin
{\center Варіант  \No \textbf { 51 }\par}
\bigskip
%type = absense
%variant = -2
Обробляються дані семестрового журналу перевірок відвідування занять студентами деканатом (фіксуються відсутні). 

\underline{Записи основного файлу містять поля}: курс, день тижня, пара, вид занять, прізвище, ім’я, по-батькові, аудиторія, код групи, навчальний тиждень, предмет.

\underline{У допоміжному файлі наявні ключі}: сума днів тижня, сума курсів.

\underline{Знайти} що пропустили найменше занять. Вивести по кожному з них інформацію:\par
{\noindent}--- на першому рядку:\par
прізвище, ім’я, по-батькові, кількість пропущених пар, кількість пропущених лекцій, кількість пропущених лабораторних, кількість пропущених практичних;

{\noindent}--- на наступних рядках, починаючи з табуляції, вивести пропуски студента (по одному на рядок):\par
аудиторія, тиждень, день, пара, предмет, вид занять

{\noindent}у такому сортуванні: вид занять, тиждень, день, пара, предмет.

%Студентів сортувати: кількість пропусків, прізвище, ім’я, по-батькові.

\bigskip\textit{Обмеження}

Навчальний тиждень може приймати цілі значення від 1 до 19.
День тижня є цілим від 1 до 5. Пара є цілим від 1 до 4.
Курс є цілим від 1 до 4. Аудиторія є додатнім цілим.
Вид занять може приймати значення:
Lecture --- лекція, Pr --- практичне, сем. --- семінар, ЛАБ --- лабораторне.

Решта полів є непорожніми рядками.
Рядки починатися та закінчуватися білими символами не можуть.

Групи кодуються в межах курсу.
Код групи складається не більш ніж з 6 символів.
Може містити літери алфавіту, десяткові цифри та дефіс.

Предмет є рядком, може мати від 5 до 27 символів включно.
Може містити літери алфавіту, десяткові цифри, апостроф, дефіс, пробіл.
Предмети різних курсів вважати різними навіть за однакових назв.

Прізвище, ім’я та по-батькові (за наявності у варіанті)
можуть мати від 5 до 30, 26, 24 відповідно символів включно кожне.
Можуть містити літери алфавіту, апостроф, пробіл та дефіс.
Студенти однозначно ідентифікуються прізвищем, ім’ям та по-батькові (за наявності у варіанті).
Жоден студент не вчиться одночасно на різних курсах чи в різних групах.

Види занять сортуються так: практичне < семінар < лабораторне < лекція.

%\bigskip%\textit{Для деяких варіантів може знадобитись}




\par
%end
\newpage
%begin
{\center Варіант  \No \textbf { 52 }\par}
\bigskip
%type = session
%variant = 57
Обробляються результати складання студентами екзаменів зимової сесії.

\underline{Записи основного файлу містять поля}: сумарна оцінка в балах за 100-бальною системою, набрані впродовж семестру бали, оцінка за державною шкалою, набрані на екзамені бали, код групи, предмет, ім’я, номер залікової книжки, прізвище.

\underline{У допоміжному файлі наявні ключі}: кількість недопусків, кількість студентів без «хвостів», довжина найкоротшого прізвища.

\underline{Знайти} всіх студентів, що по кожному предмету отримали оцінку 5 за державною шкалою. Вивести по кожному з них інформацію:\par
{\noindent}--- на першому рядку:\par
прізвище, ім’я, номер залікової книжки, кількість екзаменів, стабільність, рейтинг за балами (округлений до 1 знаку);

{\noindent}--- на наступних рядках, починаючи з табуляції, вивести для студента всі його результати (по одному на рядок):\par
бали в семестрі, бали на екзамені, сумарна оцінка, предмет

{\noindent}у такому сортуванні: бали в семестрі, предмет.

%Студентів сортувати: кількість екзаменів, рейтинг за балами (після округлення), прізвище, номер залікової книжки.

\bigskip\textit{Обмеження}

Студент не може складати предмет більше одного разу.

Бали та оцінка є значеннями цілих числових типів. Решта полів є непорожніми рядками.
Рядки розпочинатись та закінчуватись білими символами не можуть.

Оцінка за державною шкалою може приймати значення: 2-5~--- як зазвичай, 0~--- не з’явився,$-1$~--- не допущений. 
Оцінки 3-5 вважаються задовільними, решта~--- незадовільними. 
Бали є цілими та невід’ємними.
Набрані впродовж семестру бали не можуть перевищувати 60. 
Набрані на екзамені бали не можуть перевищувати 40. 
Набрані на екзамені бали повинні бути не менше 24 або дорівнювати 0(в останньому випадку оцінка не може бути задовільною). 
Сумарна оцінка в балах є сумою балів, набраних у семестрі та на екзамені.
Сумарна оцінка в балах та за державною шкалою мають бути узгоджені між собою.

Прізвище, ім’я та по-батькові (якщо наявне у варіанті) можуть мати до 24, 27, 25 відповідно символів включно кожне. 
Крім літер алфавіту можуть містити апостроф, дефіс та пробіл.

Предмет може мати від 4 до 70 символів включно. 
Крім літер алфавіту може містити подвійні лапки, пробіл та дефіс.

Код групи є непорожнім та складається не більш ніж з 4 літер. 
Крім літер алфавіту може містити десяткові цифри та дефіс.


Номер залікової книжки однозначно ідентифікує студента та складається з 7 десяткових цифр. 

\bigskip\textit{Для деяких варіантів може знадобитись}

Середній бал/оцінка за предметом визначається як середнє арифметичнесумарних балів/державних оцінок за цим предметом. 
Середній бал/оцінка студента визначається як середнє арифметичне сумарних балів/державних оцінокцього студента. 
За обчислення середньої оцінки недопущені та ті, що не з’явилися, рахуються як 2.
У решті випадків недопуски підсумовуються як $-1$, а «не з’явився» як 0.
Якщо студент не гірше ніж задовільно склав усі предмети, 
то його рейтинг за балами/оцінкою визначається як середній бал/оцінка; інакше дорівнює 0.

Стабільність студента визначається як модуль різниці між його кращим та найгіршим сумарним балом.
Стабільність предмета визначається як модуль різниці між найкращим та найгіршимсумарним балом за цим предметом.






\par
%end
\newpage
%begin
{\center Варіант  \No \textbf { 53 }\par}
\bigskip
%type = absense
%variant = 50
Обробляються дані семестрового журналу перевірок відвідування занять студентами деканатом (фіксуються відсутні). 

\underline{Записи основного файлу містять поля}: курс, код групи, по-батькові, вид занять, навчальний тиждень, ім’я, предмет, день тижня, пара, прізвище, аудиторія.

\underline{У допоміжному файлі наявні ключі}: загальна кількість пропусків третіх пар, найменша довжина назв предметів, найбільша довжина назв предметів.

\underline{Знайти} всі предмети, які найбільше пропускалися по понеділках. Вивести по кожному з них інформацію:\par
{\noindent}--- на першому рядку:\par
курс, предмет, середня кількість пропусків по днях тижня (округлена до 2 знаків), кількість пропусків практичних;

{\noindent}--- на наступних рядках, починаючи з табуляції, вивести пропуски предмета (по одному на рядок):\par
прізвище, ім’я, по-батькові, тиждень, день, пара, вид занять

{\noindent}у такому сортуванні: день, прізвище, ім’я, по-батькові, тиждень, пара, вид занять.

%Предмети сортувати: середня кількість пропусків по днях тижня (після округлення, за спаданням), кількість пропусків практичних, курс, предмет.

\bigskip\textit{Обмеження}

Навчальний тиждень може приймати цілі значення від 1 до 16.
День тижня є цілим від 1 до 5. Пара є цілим від 1 до 4.
Курс є цілим від 1 до 4. Аудиторія є додатнім цілим.
Вид занять може приймати значення:
Л --- лекція, Практ. --- практичне, Sem --- семінар, lab --- лабораторне.

Решта полів є непорожніми рядками.
Рядки починатися та закінчуватися білими символами не можуть.

Групи кодуються в межах курсу.
Код групи складається не більш ніж з 4 символів.
Може містити літери алфавіту, десяткові цифри та дефіс.

Предмет є рядком, може мати від 4 до 34 символів включно.
Може містити літери алфавіту, десяткові цифри, апостроф, дефіс, пробіл.
Предмети різних курсів вважати різними навіть за однакових назв.

Прізвище, ім’я та по-батькові (за наявності у варіанті)
можуть мати від 4 до 21, 24, 30 відповідно символів включно кожне.
Можуть містити літери алфавіту, апостроф, пробіл та дефіс.
Студенти однозначно ідентифікуються прізвищем, ім’ям та по-батькові (за наявності у варіанті).
Жоден студент не вчиться одночасно на різних курсах чи в різних групах.

Види занять сортуються так: лекція < практичне < лабораторне < семінар.

%\bigskip%\textit{Для деяких варіантів може знадобитись}




\par
%end
\newpage
%begin
{\center Варіант  \No \textbf { 54 }\par}
\bigskip
%type = absense
%variant = 46
Обробляються дані семестрового журналу перевірок відвідування занять студентами деканатом (фіксуються відсутні). 

\underline{Записи основного файлу містять поля}: вид занять, ім’я, предмет, курс, навчальний тиждень, аудиторія, прізвище, день тижня, пара.

\underline{У допоміжному файлі наявні ключі}: кількість пропусків по вівторках, загальна кількість зафіксованих пропусків.

\underline{Знайти} всі предмети, з яких пропускалися тільки лекції. Вивести по кожному з них інформацію:\par
{\noindent}--- на першому рядку:\par
курс, предмет, кількість пропусків;

{\noindent}--- на наступних рядках, починаючи з табуляції, вивести пропуски предмета (по одному на рядок):\par
прізвище, ім’я, кількість пропусків лекцій

{\noindent}у такому сортуванні: прізвище, ім’я.

%Предмети сортувати: курс, предмет.

\bigskip\textit{Обмеження}

Навчальний тиждень може приймати цілі значення від 1 до 17.
День тижня є цілим від 1 до 5. Пара є цілим від 1 до 4.
Курс є цілим від 1 до 4. Аудиторія є додатнім цілим.
Вид занять може приймати значення:
L --- лекція, Пр --- практичне, сем. --- семінар, Лаб. --- лабораторне.

Решта полів є непорожніми рядками.
Рядки починатися та закінчуватися білими символами не можуть.

Групи кодуються в межах курсу.
Код групи складається не більш ніж з 5 символів.
Може містити літери алфавіту, десяткові цифри та дефіс.

Предмет є рядком, може мати від 6 до 38 символів включно.
Може містити літери алфавіту, десяткові цифри, апостроф, дефіс, пробіл.
Предмети різних курсів вважати різними навіть за однакових назв.

Прізвище, ім’я та по-батькові (за наявності у варіанті)
можуть мати від 6 до 20, 27, 27 відповідно символів включно кожне.
Можуть містити літери алфавіту, апостроф, пробіл та дефіс.
Студенти однозначно ідентифікуються прізвищем, ім’ям та по-батькові (за наявності у варіанті).
Жоден студент не вчиться одночасно на різних курсах чи в різних групах.

Види занять сортуються так: семінар < лекція < лабораторне < практичне.

%\bigskip%\textit{Для деяких варіантів може знадобитись}




\par
%end
\newpage
%begin
{\center Варіант  \No \textbf { 55 }\par}
\bigskip
%type = tasks
%variant = 6
Обробляються результати розв’язування студентами задач на контрольних.
Одиничний факт роз\-в’я\-зу\-вання студентом задачі надалі розуміється як спроба.

\underline{Записи основного файлу містять поля}: день, прізвище, ім’я, місяць, відсоток набраних за розв’язан\-ня задачі балів, умова задачі, по-батькові, рік.

\underline{У допоміжному файлі наявні ключі}: сума всіх днів, сума всіх місяців.

\underline{Знайти} всіх студентів, що мають найбільшу кількість спроб з відсотком 90\% та більше. Вивести по кожному з них інформацію:\par
{\noindent}--- на першому рядку:\par
ім’я, по-батькові, прізвище, кількість спроб більших 90\%, кількість спроб;

{\noindent}--- на наступних рядках, починаючи з табуляції, вивести їх спроби з відсотком більше 50\% (по одному на рядок):\par
рік, відсоток набраних за задачу балів, умова

{\noindent}у такому сортуванні: умова, відсоток набраних балів, рік(за спаданням).

%Студентів сортувати: прізвище, ім’я, по-батькові.

\bigskip\textit{Обмеження}

Студент міг розв’язувати задачу кілька разів.

Відсоток набраних за розв’язання задачі балів є цілим від 0 до 100.
День, місяць, рік є числовими полями, що зображують дату.
Раніше 2001 року статистика не велася. Рік є 4-значним.
Решта полів є непорожніми рядками.
Рядки починатися та закінчуватися білими символами не можуть.

Прізвище, ім’я та по-батькові (за наявності у варіанті) можуть мати до 28, 21, 21 відповідно символів включно кожне. 
Крім літер алфавіту можуть містити апостроф, дефіс та пробіл.

Умова задачі може мати від 5 до 66 символів включно.
Крім літер алфавіту може містити подвійні лапки, десяткові цифри, пробіл, дефіс,
а також символи . , \_ \{ \} \$ $+$ $*$ {\textbackslash} ( ) \% $[$ $]$ .

Студенти однозначно ідентифікуються прізвищем, ім’ям та по-батькові (за наявності у варіанті).
Задачі однозначно ідентифікуються умовою.

\bigskip\textit{Для деяких варіантів може знадобитись}

Розв'язуваність задачі --- середній відсоток набраних за задачу балів.

Розв'язуваність студента --- середній відсоток набраних студентом балів по всіх його спробах.

Розв'язуваність --- середній відсоток набраних студентами за задачі балів (рахується по всім даним).




\par
%end
\newpage
%begin
{\center Варіант  \No \textbf { 56 }\par}
\bigskip
%type = session
%variant = 66
Обробляються результати складання студентами екзаменів зимової сесії.

\underline{Записи основного файлу містять поля}: предмет, номер залікової книжки, ім’я, набрані на екзамені бали, код групи, по-батькові, прізвище, оцінка за державною шкалою, сумарна оцінка в балах за 100-бальною системою.

\underline{У допоміжному файлі наявні ключі}: найбільша довжина назви предмета, найменша довжина назви предмета, найбільший номер залікової книжки.

\underline{Знайти} студентів, що не отримали жодної оцінки «добре». Вивести по кожному з них інформацію:\par
{\noindent}--- на першому рядку:\par
кількість «хвостів», кількість екзаменів, номер залікової книжки, прізвище, ім’я, по-батькові;

{\noindent}--- на наступних рядках, починаючи з табуляції, вивести для студента всі його результати (по одному на рядок):\par
бали в семестрі, бали на екзамені, предмет, державна оцінка

{\noindent}у такому сортуванні: бали на екзамені, предмет.

%Студентів сортувати: кількість «хвостів», ім’я, прізвище, номер залікової книжки.

\bigskip\textit{Обмеження}

Студент не може складати предмет більше одного разу.

Бали та оцінка є значеннями цілих числових типів. Решта полів є непорожніми рядками.
Рядки розпочинатись та закінчуватись білими символами не можуть.

Оцінка за державною шкалою може приймати значення: 2-5~--- як зазвичай, 0~--- не з’явився,$-1$~--- не допущений. 
Оцінки 3-5 вважаються задовільними, решта~--- незадовільними. 
Бали є цілими та невід’ємними.
Набрані впродовж семестру бали не можуть перевищувати 60. 
Набрані на екзамені бали не можуть перевищувати 40. 
Набрані на екзамені бали повинні бути не менше 24 або дорівнювати 0(в останньому випадку оцінка не може бути задовільною). 
Сумарна оцінка в балах є сумою балів, набраних у семестрі та на екзамені.
Сумарна оцінка в балах та за державною шкалою мають бути узгоджені між собою.

Прізвище, ім’я та по-батькові (якщо наявне у варіанті) можуть мати до 28, 23, 27 відповідно символів включно кожне. 
Крім літер алфавіту можуть містити апостроф, дефіс та пробіл.

Предмет може мати від 3 до 64 символів включно. 
Крім літер алфавіту може містити подвійні лапки, пробіл та дефіс.

Код групи є непорожнім та складається не більш ніж з 7 літер. 
Крім літер алфавіту може містити десяткові цифри та дефіс.


Номер залікової книжки однозначно ідентифікує студента та складається з 8 десяткових цифр. 

\bigskip\textit{Для деяких варіантів може знадобитись}

Середній бал/оцінка за предметом визначається як середнє арифметичнесумарних балів/державних оцінок за цим предметом. 
Середній бал/оцінка студента визначається як середнє арифметичне сумарних балів/державних оцінокцього студента. 
За обчислення середньої оцінки недопущені та ті, що не з’явилися, рахуються як 2.
У решті випадків недопуски підсумовуються як $-1$, а «не з’явився» як 0.
Якщо студент не гірше ніж задовільно склав усі предмети, 
то його рейтинг за балами/оцінкою визначається як середній бал/оцінка; інакше дорівнює 0.

Стабільність студента визначається як модуль різниці між його кращим та найгіршим сумарним балом.
Стабільність предмета визначається як модуль різниці між найкращим та найгіршимсумарним балом за цим предметом.






\par
%end
\newpage
%begin
{\center Варіант  \No \textbf { 57 }\par}
\bigskip
%type = meteo
%variant = 86
Обробляються дані спостережень за погодою. 

\underline{Записи основного файлу містять поля}: код метеостанції, середня денна температура, кількість опадів, рік, відносна вологість, сила вітру, місяць, день, максимальна денна температура, мінімальна денна температура.

\underline{У допоміжному файлі наявні ключі}: загальна кількість записів, сумарна сила вітру.

\underline{Знайти} 8 дат, по яких фіксувалася найбільша середня денна температура. Вивести по кожному з них інформацію:\par
{\noindent}--- на першому рядку:\par
середня середньоденна температура, найбільша середня денна температура, місяць, рік, день, кількість опадів;

{\noindent}--- на наступних рядках, починаючи з табуляції, вивести для дати спостереження метеостанції, де сила вітру була меншою за середню по даті (по одному на рядок):\par
вологість, сила вітру, середня температура, метеостанція

{\noindent}у такому сортуванні: вологість, метеостанція.

%Дати сортувати: середня середньоденна температура (за спаданням), рік, місяць, день.

\bigskip\textit{Обмеження}

Код метеостанції є рядком довжини від 5 до 11 включно.
Може містити літери алфавіту, десяткові цифри, підкреслення. Решта полів числові.

Рік є чотиризначним цілим.
Кількість опадів є дійсним значенням, подається у форматі з 1 знаком після крапки.
Сила вітру є дійсним значенням, подається у форматі з 1 знаком після крапки.
Температури є дійсними значеннями, подаються у форматі з 3 знаками після крапки.
Цілі значення подаються без десяткової крапки.
Для дійсних використовується зображення з фіксованою крапкою.

\textit{Вихідне форматування дійсних}

Усі дійсні значення виводяться у форматі з фіксованою крапкою з тією ж кількістю знаків після крапки, що і у вхідному зображенні.

Якщо у варіанті раніше не було зазначено, то \underline{обчислювані програмою} середні значення виводяться у форматі з 3 знаками після крапки.
%\bigskip%\textit{Для деяких варіантів може знадобитись}




\par
%end
\newpage
%begin
{\center Варіант  \No \textbf { 58 }\par}
\bigskip
%type = booking
%variant = 125
Обробляється інформація щодо відвантажених накладних.

\underline{Записи основного файлу містять поля}: кількість, одиниця виміру, ціна, назва товару, номер товарної позиції у межах накладної, вартість, номер накладної.

\underline{У допоміжному файлі наявні ключі}: кількість записів, кількість різних товарів.

\underline{Знайти} 7 найдешевших товарів (з найменшою середньою ціною відвантаження). Вивести по кожному з них інформацію:\par
{\noindent}--- на першому рядку:\par
назва товару, одиниця виміру, кількість накладних, де відвантажувався товар, кількість накладних, де він відвантажувався за своєю максимальною ціною, середня ціна відвантаження (з 2 знаками після крапки);

{\noindent}--- на наступних рядках, починаючи з табуляції, вивести для товару всі накладні, де він відвантажувався за своєю мінімальною ціною (по одному на рядок):\par
номер накладної, ціна, кількість, ціна, вартість

{\noindent}у такому сортуванні: номер накладної, ціна.

%Товари сортувати: кількість накладних, де відвантажувався товар, назва товару, одиниця виміру.

\bigskip\textit{Обмеження}

Кількість, ціна та вартість є значеннями дійсного типу.
Решта полів є непорожніми рядками.
Рядки починатися та закінчуватися білими символами не можуть.

Номер накладної складається з 15 символів (цифри, латинські літери).
Назва товару може мати від 3 до 25 символів включно.
Може містити тільки цифри, літери, пробіли, апостроф, дефіс, підкреслення.

Кількість, ціна мають бути додатними.
Кількість вказується з рівно 1 знаками після крапки.
Ціна вказується з рівно 2 знаками після крапки.
Вартість має бути узгоджена з кількістю та ціною та записується рівно з двома десятковими знаками.
(Застосовується округлення.)

Товари однозначно ідентифікуються назвами та одиницями виміру (за наявності).

Одиниця виміру (за наявності у варіанті) може приймати значення: шт., кг, м , кор.

Кожна накладна має свій власний унікальний номер.
У накладній можуть відвантажуватися кілька товарів (не менше 1).
Товарна позиція --- рядок накладної: товар, кількість, ціна, вартість та, за наявності, одиниця виміру.
Товарні позиції кожної накладної нумеруються послідовними натуральними числами, починаючи з 1.
В одній накладній товари можуть повторюватися, 
тобто товарні позиції однієї накладної можуть містити однакові товари 
(при цьому ціна може бути і різною, і однаковою).

\bigskip\textit{Для деяких варіантів може знадобитись}

Недоотриману вигоду розраховувати як суму виразів
{\center (найбільша ціна відвантаження товару – ціна відвантаження товару)*кількість\par}

\bigskip
{\noindent} округлених до 2 знаків після крапки.




\par
%end
\newpage
%begin
{\center Варіант  \No \textbf { 59 }\par}
\bigskip
%type = meteo
%variant = 92
Обробляються дані спостережень за погодою. 

\underline{Записи основного файлу містять поля}: відносна вологість, сила вітру, рік, кількість опадів, мінімальна денна температура, день, місяць, код метеостанції, максимальна денна температура, середня денна температура.

\underline{У допоміжному файлі наявні ключі}: кількість спостережень, в яких не спостерігалися опади; кількість спостережень за січень.

\underline{Знайти} дати, по яких фіксувалася сила вітру більша ніж в середньому по всіх спостереженнях. Вивести по кожному з них інформацію:\par
{\noindent}--- на першому рядку:\par
місяць, рік, день, середня сила вітру, найбільша сила вітру, сумарна кількість опадів за дату;

{\noindent}--- на наступних рядках, починаючи з табуляції, вивести для дати спостереження метеостанції, де сила вітру була менша ніж середня по всіх спостереженнях (всіх станцій) (по одному на рядок):\par
сила вітру, вологість, метеостанція

{\noindent}у такому сортуванні: метеостанція.

%Дати сортувати: день, рік, місяць, найбільша сила вітру.

\bigskip\textit{Обмеження}

Код метеостанції є рядком довжини від 4 до 12 включно.
Може містити літери алфавіту, десяткові цифри, підкреслення. Решта полів числові.

Рік є чотиризначним цілим.
Кількість опадів є цілим значенням.
Сила вітру є цілим значенням.
Температури є дійсними значеннями, подаються у форматі з 1 знаком після крапки.
Цілі значення подаються без десяткової крапки.
Для дійсних використовується зображення з фіксованою крапкою.

\textit{Вихідне форматування дійсних}

Усі дійсні значення виводяться у форматі з фіксованою крапкою з тією ж кількістю знаків після крапки, що і у вхідному зображенні.

Якщо у варіанті раніше не було зазначено, то \underline{обчислювані програмою} середні значення виводяться у форматі з 1 знаком після крапки.
%\bigskip%\textit{Для деяких варіантів може знадобитись}




\par
%end
\newpage
%begin
{\center Варіант  \No \textbf { 60 }\par}
\bigskip
%type = session
%variant = 76
Обробляються результати складання студентами екзаменів зимової сесії.

\underline{Записи основного файлу містять поля}: набрані на екзамені бали, оцінка за державною шкалою, ім’я, номер залікової книжки, по-батькові, сумарна оцінка в балах за 100-бальною системою, код групи, прізвище, предмет.

\underline{У допоміжному файлі наявні ключі}: кількість оцінок «добре», загальна кількість записів.

\underline{Знайти} предмети, стабільність яких вища середньої стабільності по всіх предметах. Вивести по кожному з них інформацію:\par
{\noindent}--- на першому рядку:\par
стабільність, предмет, кількість студентів, що не склали;

{\noindent}--- на наступних рядках, починаючи з табуляції, вивести для предмета всіх студентів, що його складали (по одному на рядок):\par
бали за екзамен, оцінка в балах, ім’я, прізвище, номер залікової книжки, група

{\noindent}у такому сортуванні: група, номер залікової книжки.

%Предмети сортувати: стабільність (за спаданням), кількість студентів, що не склали, предмет.

\bigskip\textit{Обмеження}

Студент не може складати предмет більше одного разу.

Бали та оцінка є значеннями цілих числових типів. Решта полів є непорожніми рядками.
Рядки розпочинатись та закінчуватись білими символами не можуть.

Оцінка за державною шкалою може приймати значення: 2-5~--- як зазвичай, 0~--- не з’явився,$-1$~--- не допущений. 
Оцінки 3-5 вважаються задовільними, решта~--- незадовільними. 
Бали є цілими та невід’ємними.
Набрані впродовж семестру бали не можуть перевищувати 60. 
Набрані на екзамені бали не можуть перевищувати 40. 
Набрані на екзамені бали повинні бути не менше 24 або дорівнювати 0(в останньому випадку оцінка не може бути задовільною). 
Сумарна оцінка в балах є сумою балів, набраних у семестрі та на екзамені.
Сумарна оцінка в балах та за державною шкалою мають бути узгоджені між собою.

Прізвище, ім’я та по-батькові (якщо наявне у варіанті) можуть мати до 24, 24, 29 відповідно символів включно кожне. 
Крім літер алфавіту можуть містити апостроф, дефіс та пробіл.

Предмет може мати від 5 до 61 символів включно. 
Крім літер алфавіту може містити подвійні лапки, пробіл та дефіс.

Код групи є непорожнім та складається не більш ніж з 5 літер. 
Крім літер алфавіту може містити десяткові цифри та дефіс.


Номер залікової книжки однозначно ідентифікує студента та складається з 6 десяткових цифр. 

\bigskip\textit{Для деяких варіантів може знадобитись}

Середній бал/оцінка за предметом визначається як середнє арифметичнесумарних балів/державних оцінок за цим предметом. 
Середній бал/оцінка студента визначається як середнє арифметичне сумарних балів/державних оцінокцього студента. 
За обчислення середньої оцінки недопущені та ті, що не з’явилися, рахуються як 2.
У решті випадків недопуски підсумовуються як $-1$, а «не з’явився» як 0.
Якщо студент не гірше ніж задовільно склав усі предмети, 
то його рейтинг за балами/оцінкою визначається як середній бал/оцінка; інакше дорівнює 0.

Стабільність студента визначається як модуль різниці між його кращим та найгіршим сумарним балом.
Стабільність предмета визначається як модуль різниці між найкращим та найгіршимсумарним балом за цим предметом.






\par
%end
\newpage
%begin
{\center Варіант  \No \textbf { 61 }\par}
\bigskip
%type = tasks
%variant = 11
Обробляються результати розв’язування студентами задач на контрольних.
Одиничний факт роз\-в’я\-зу\-вання студентом задачі надалі розуміється як спроба.

\underline{Записи основного файлу містять поля}: ім’я, місяць, день, відсоток набраних за розв’язання задачі балів, умова задачі, прізвище, рік.

\underline{У допоміжному файлі наявні ключі}: найменший рік спроби, кількість записів.

\underline{Знайти} всі задачі, за якими всі спроби лежать у межах від 60\% до 90\% включно. Вивести по кожному з них інформацію:\par
{\noindent}--- на першому рядку:\par
кількість спроб, середня розв’язуваність (округлена до 1 знаку), умова;

{\noindent}--- на наступних рядках, починаючи з табуляції, вивести спроби по цих задачах (по одному на рядок):\par
прізвище, ім’я, відсоток набраних балів

{\noindent}у такому сортуванні: прізвище, ім’я, відсоток набраних балів (за спаданням).

%Задачі сортувати: кількість спроб, середня розв’язуваність, умова.

\bigskip\textit{Обмеження}

Студент міг розв’язувати задачу кілька разів.

Відсоток набраних за розв’язання задачі балів є цілим від 0 до 100.
День, місяць, рік є числовими полями, що зображують дату.
Раніше 2001 року статистика не велася. Рік є 4-значним.
Решта полів є непорожніми рядками.
Рядки починатися та закінчуватися білими символами не можуть.

Прізвище, ім’я та по-батькові (за наявності у варіанті) можуть мати до 30, 26, 26 відповідно символів включно кожне. 
Крім літер алфавіту можуть містити апостроф, дефіс та пробіл.

Умова задачі може мати від 3 до 55 символів включно.
Крім літер алфавіту може містити подвійні лапки, десяткові цифри, пробіл, дефіс,
а також символи . , \_ \{ \} \$ $+$ $*$ {\textbackslash} ( ) \% $[$ $]$ .

Студенти однозначно ідентифікуються прізвищем, ім’ям та по-батькові (за наявності у варіанті).
Задачі однозначно ідентифікуються умовою.

\bigskip\textit{Для деяких варіантів може знадобитись}

Розв'язуваність задачі --- середній відсоток набраних за задачу балів.

Розв'язуваність студента --- середній відсоток набраних студентом балів по всіх його спробах.

Розв'язуваність --- середній відсоток набраних студентами за задачі балів (рахується по всім даним).




\par
%end
\newpage
%begin
{\center Варіант  \No \textbf { 62 }\par}
\bigskip
%type = meteo
%variant = 99
Обробляються дані спостережень за погодою. 

\underline{Записи основного файлу містять поля}: мінімальна денна температура, відносна вологість, сила вітру, середня денна температура, максимальна денна температура, код метеостанції, місяць, день, кількість опадів, рік.

\underline{У допоміжному файлі наявні ключі}: найменша мінімальна температура, найбільша сила вітру.

\underline{Знайти} дати, по яких середня кількість опадів була більша ніж в середньому по всіх спостереженнях. Вивести по кожному з них інформацію:\par
{\noindent}--- на першому рядку:\par
рік, місяць, день, кількість опадів, середня вологість;

{\noindent}--- на наступних рядках, починаючи з табуляції, вивести для дати спостереження метеостанції, де вологість була більша 50 (по одному на рядок):\par
сила вітру, вологість, метеостанція

{\noindent}у такому сортуванні: вологість, метеостанція.

%Дати сортувати: середня вологість, день, місяць, рік.

\bigskip\textit{Обмеження}

Код метеостанції є рядком довжини від 4 до 8 включно.
Може містити літери алфавіту, десяткові цифри, підкреслення. Решта полів числові.

Рік є чотиризначним цілим.
Кількість опадів є дійсним значенням, подається у форматі з 1 знаком після крапки.
Сила вітру є дійсним значенням, подається у форматі з 2 знаками після крапки.
Температури є дійсними значеннями, подаються у форматі з 3 знаками після крапки.
Цілі значення подаються без десяткової крапки.
Для дійсних використовується зображення з фіксованою крапкою.

\textit{Вихідне форматування дійсних}

Усі дійсні значення виводяться у форматі з фіксованою крапкою з тією ж кількістю знаків після крапки, що і у вхідному зображенні.

Якщо у варіанті раніше не було зазначено, то \underline{обчислювані програмою} середні значення виводяться у форматі з 3 знаками після крапки.
%\bigskip%\textit{Для деяких варіантів може знадобитись}




\par
%end
\newpage
%begin
{\center Варіант  \No \textbf { 63 }\par}
\bigskip
%type = meteo
%variant = 83
Обробляються дані спостережень за погодою. 

\underline{Записи основного файлу містять поля}: місяць, день, кількість опадів, відносна вологість, код метеостанції, середня денна температура, мінімальна денна температура, максимальна денна температура, сила вітру, рік.

\underline{У допоміжному файлі наявні ключі}: кількість різних дат, довжина найдовшого коду метеостанції.

\underline{Знайти} 6 дат, за які випала найменша кількість опадів (сумарно по всіх метеостанціях). Вивести по кожному з них інформацію:\par
{\noindent}--- на першому рядку:\par
кількість опадів, рік, місяць, день, середня вологість;

{\noindent}--- на наступних рядках, починаючи з табуляції, вивести для дати спостереження метеостанцій, де сила вітру була більша за 3.2 (по одному на рядок):\par
середня денна температура, максимальна денна температура, мінімальна денна температура, сила вітру, метеостанція

{\noindent}у такому сортуванні: середня денна температура, метеостанція.

%Дати сортувати: місяць, день, кількість опадів, рік.

\bigskip\textit{Обмеження}

Код метеостанції є рядком довжини від 5 до 6 включно.
Може містити літери алфавіту, десяткові цифри, підкреслення. Решта полів числові.

Рік є чотиризначним цілим.
Кількість опадів є цілим значенням.
Сила вітру є дійсним значенням, подається у форматі з 1 знаком після крапки.
Температури є дійсними значеннями, подаються у форматі з 1 знаком після крапки.
Цілі значення подаються без десяткової крапки.
Для дійсних використовується зображення з фіксованою крапкою.

\textit{Вихідне форматування дійсних}

Усі дійсні значення виводяться у форматі з фіксованою крапкою з тією ж кількістю знаків після крапки, що і у вхідному зображенні.

Якщо у варіанті раніше не було зазначено, то \underline{обчислювані програмою} середні значення виводяться у форматі з 2 знаками після крапки.
%\bigskip%\textit{Для деяких варіантів може знадобитись}




\par
%end
\newpage
%begin
{\center Варіант  \No \textbf { 64 }\par}
\bigskip
%type = absense
%variant = 51
Обробляються дані семестрового журналу перевірок відвідування занять студентами деканатом (фіксуються відсутні). 

\underline{Записи основного файлу містять поля}: день тижня, пара, навчальний тиждень, ім’я, вид занять, прізвище, аудиторія, предмет, курс.

\underline{У допоміжному файлі наявні ключі}: загальна кількість предметів, кількість пропусків перших пар.

\underline{Знайти} всі предмети, що по середах пропускалися більше ніж у середньому. Вивести по кожному з них інформацію:\par
{\noindent}--- на першому рядку:\par
курс, предмет, кількість пропусків, кількість пропусків по середах;

{\noindent}--- на наступних рядках, починаючи з табуляції, вивести пропуски предмета (по одному на рядок):\par
прізвище, ім’я, день, вид занять, кількість пропусків

{\noindent}у такому сортуванні: день, вид занять, прізвище, ім’я.

%Предмети сортувати: кількість пропусків по середах, курс, предмет.

\bigskip\textit{Обмеження}

Навчальний тиждень може приймати цілі значення від 1 до 18.
День тижня є цілим від 1 до 5. Пара є цілим від 1 до 4.
Курс є цілим від 1 до 4. Аудиторія є додатнім цілим.
Вид занять може приймати значення:
Л --- лекція, Practice --- практичне, Сем. --- семінар, лаб. --- лабораторне.

Решта полів є непорожніми рядками.
Рядки починатися та закінчуватися білими символами не можуть.

Групи кодуються в межах курсу.
Код групи складається не більш ніж з 6 символів.
Може містити літери алфавіту, десяткові цифри та дефіс.

Предмет є рядком, може мати від 5 до 29 символів включно.
Може містити літери алфавіту, десяткові цифри, апостроф, дефіс, пробіл.
Предмети різних курсів вважати різними навіть за однакових назв.

Прізвище, ім’я та по-батькові (за наявності у варіанті)
можуть мати від 5 до 27, 29, 21 відповідно символів включно кожне.
Можуть містити літери алфавіту, апостроф, пробіл та дефіс.
Студенти однозначно ідентифікуються прізвищем, ім’ям та по-батькові (за наявності у варіанті).
Жоден студент не вчиться одночасно на різних курсах чи в різних групах.

Види занять сортуються так: лабораторне < практичне < лекція < семінар.

%\bigskip%\textit{Для деяких варіантів може знадобитись}




\par
%end
\newpage
%begin
{\center Варіант  \No \textbf { 65 }\par}
\bigskip
%type = session
%variant = 55
Обробляються результати складання студентами екзаменів зимової сесії.

\underline{Записи основного файлу містять поля}: предмет, по-батькові, набрані на екзамені бали, оцінка за державною шкалою, прізвище, код групи, набрані впродовж семестру бали, ім’я, номер залікової книжки, сумарна оцінка в балах за 100-бальною системою.

\underline{У допоміжному файлі наявні ключі}: загальна кількість записів, кількість «не з’явився».

\underline{Знайти} 4 найменш стабільних студенти. Вивести по кожному з них інформацію:\par
{\noindent}--- на першому рядку:\par
прізвище, ім’я, по-батькові, номер залікової книжки, стабільність, рейтинг за оцінкою (округлений до 1 знаку);

{\noindent}--- на наступних рядках, починаючи з табуляції, вивести для студента всі його результати (по одному на рядок):\par
державна оцінка, предмет, сумарна оцінка, бали на екзамені, бали в семестрі

{\noindent}у такому сортуванні: державна оцінка, предмет.

%Студентів сортувати: рейтинг за балами (після округлення), прізвище, ім’я, номер залікової книжки.

\bigskip\textit{Обмеження}

Студент не може складати предмет більше одного разу.

Бали та оцінка є значеннями цілих числових типів. Решта полів є непорожніми рядками.
Рядки розпочинатись та закінчуватись білими символами не можуть.

Оцінка за державною шкалою може приймати значення: 2-5~--- як зазвичай, 0~--- не з’явився,$-1$~--- не допущений. 
Оцінки 3-5 вважаються задовільними, решта~--- незадовільними. 
Бали є цілими та невід’ємними.
Набрані впродовж семестру бали не можуть перевищувати 60. 
Набрані на екзамені бали не можуть перевищувати 40. 
Набрані на екзамені бали повинні бути не менше 24 або дорівнювати 0(в останньому випадку оцінка не може бути задовільною). 
Сумарна оцінка в балах є сумою балів, набраних у семестрі та на екзамені.
Сумарна оцінка в балах та за державною шкалою мають бути узгоджені між собою.

Прізвище, ім’я та по-батькові (якщо наявне у варіанті) можуть мати до 23, 21, 28 відповідно символів включно кожне. 
Крім літер алфавіту можуть містити апостроф, дефіс та пробіл.

Предмет може мати від 3 до 62 символів включно. 
Крім літер алфавіту може містити подвійні лапки, пробіл та дефіс.

Код групи є непорожнім та складається не більш ніж з 4 літер. 
Крім літер алфавіту може містити десяткові цифри та дефіс.


Номер залікової книжки однозначно ідентифікує студента та складається з 8 десяткових цифр. 

\bigskip\textit{Для деяких варіантів може знадобитись}

Середній бал/оцінка за предметом визначається як середнє арифметичнесумарних балів/державних оцінок за цим предметом. 
Середній бал/оцінка студента визначається як середнє арифметичне сумарних балів/державних оцінокцього студента. 
За обчислення середньої оцінки недопущені та ті, що не з’явилися, рахуються як 2.
У решті випадків недопуски підсумовуються як $-1$, а «не з’явився» як 0.
Якщо студент не гірше ніж задовільно склав усі предмети, 
то його рейтинг за балами/оцінкою визначається як середній бал/оцінка; інакше дорівнює 0.

Стабільність студента визначається як модуль різниці між його кращим та найгіршим сумарним балом.
Стабільність предмета визначається як модуль різниці між найкращим та найгіршимсумарним балом за цим предметом.






\par
%end
\newpage
%begin
{\center Варіант  \No \textbf { 66 }\par}
\bigskip
%type = meteo
%variant = 84
Обробляються дані спостережень за погодою. 

\underline{Записи основного файлу містять поля}: відносна вологість, день, місяць, код метеостанції, сила вітру, кількість опадів, рік, максимальна денна температура, середня денна температура, мінімальна денна температура.

\underline{У допоміжному файлі наявні ключі}: кількість різних дат, кількість записів.

\underline{Знайти} 4 дати, по яких фіксувалася найменша сила вітру. Вивести по кожному з них інформацію:\par
{\noindent}--- на першому рядку:\par
місяць, рік, день, сила вітру, сумарна кількість опадів;

{\noindent}--- на наступних рядках, починаючи з табуляції, вивести для дати спостереження метеостанції (по одному на рядок):\par
сила вітру, кількість опадів, середня температура, код метеостанції

{\noindent}у такому сортуванні: середня температура, метеостанція.

%Дати сортувати: день, кількість опадів, рік, місяць.

\bigskip\textit{Обмеження}

Код метеостанції є рядком довжини від 3 до 10 включно.
Може містити літери алфавіту, десяткові цифри, підкреслення. Решта полів числові.

Рік є чотиризначним цілим.
Кількість опадів є дійсним значенням, подається у форматі з 2 знаками після крапки.
Сила вітру є цілим значенням.
Температури є дійсними значеннями, подаються у форматі з 2 знаками після крапки.
Цілі значення подаються без десяткової крапки.
Для дійсних використовується зображення з фіксованою крапкою.

\textit{Вихідне форматування дійсних}

Усі дійсні значення виводяться у форматі з фіксованою крапкою з тією ж кількістю знаків після крапки, що і у вхідному зображенні.

Якщо у варіанті раніше не було зазначено, то \underline{обчислювані програмою} середні значення виводяться у форматі з 1 знаком після крапки.
%\bigskip%\textit{Для деяких варіантів може знадобитись}




\par
%end
\newpage
%begin
{\center Варіант  \No \textbf { 67 }\par}
\bigskip
%type = meteo
%variant = 115
Обробляються дані спостережень за погодою. 

\underline{Записи основного файлу містять поля}: середня денна температура, місяць, день, кількість опадів, мінімальна денна температура, рік, код метеостанції, відносна вологість, максимальна денна температура, сила вітру.

\underline{У допоміжному файлі наявні ключі}: кількість записів, максимальна зафіксована вологість, ALL IS DONE.

\underline{Знайти} метеостанції, по яких жодний з показників не відхилявся від середніх по метеостанції більше ніж на 40\%. Вивести по кожному з них інформацію:\par
{\noindent}--- на першому рядку:\par
метеостанція, кількість спостережень, найбільше, найменше та середнє значення вологості;

{\noindent}--- на наступних рядках, починаючи з табуляції, вивести для них дані по датах (по одному на рядок):\par
рік, місяць, день, середня денна температура, максимальна денна температура, мінімальна денна температура, сила вітру, вологість

{\noindent}у такому сортуванні: середня денна температура, місяць, день, рік.

%Метеостанції сортувати: найбільше значення вологості (за спаданням), кількість спостережень, метеостанція.

\bigskip\textit{Обмеження}

Код метеостанції є рядком довжини від 4 до 11 включно.
Може містити літери алфавіту, десяткові цифри, підкреслення. Решта полів числові.

Рік є чотиризначним цілим.
Кількість опадів є дійсним значенням, подається у форматі з 2 знаками після крапки.
Сила вітру є цілим значенням.
Температури є дійсними значеннями, подаються у форматі з 3 знаками після крапки.
Цілі значення подаються без десяткової крапки.
Для дійсних використовується зображення з фіксованою крапкою.

\textit{Вихідне форматування дійсних}

Усі дійсні значення виводяться у форматі з фіксованою крапкою з тією ж кількістю знаків після крапки, що і у вхідному зображенні.

Якщо у варіанті раніше не було зазначено, то \underline{обчислювані програмою} середні значення виводяться у форматі з 1 знаком після крапки.
%\bigskip%\textit{Для деяких варіантів може знадобитись}




\par
%end
\newpage
%begin
{\center Варіант  \No \textbf { 68 }\par}
\bigskip
%type = booking
%variant = 117
Обробляється інформація щодо відвантажених накладних.

\underline{Записи основного файлу містять поля}: ціна, номер накладної, назву товару, вартість, одиниця виміру, номер товарної позиції у межах накладної, кількість.

\underline{У допоміжному файлі наявні ключі}: сума цін, найменша ціна товарної позиції.

\underline{Знайти} 6 найменш популярних товарів (ті, що відвантажувалися в найменшій кількості накладних). Вивести по кожному з них інформацію:\par
{\noindent}--- на першому рядку:\par
назва товару, одиниця виміру, загальна відвантажена кількість, найменша ціна відвантаження (з 2 знаками після крапки);

{\noindent}--- на наступних рядках, починаючи з табуляції, вивести для товару всі накладні, де він відвантажувався за найбільшою ціною (по одному на рядок):\par
ціна, кількість, номер накладної

{\noindent}у такому сортуванні: ціна, номер накладної.

%Товари сортувати: загальна відвантажена кількість, найменша ціна відвантаження, назва товару, одиниця виміру.

\bigskip\textit{Обмеження}

Кількість, ціна та вартість є значеннями дійсного типу.
Решта полів є непорожніми рядками.
Рядки починатися та закінчуватися білими символами не можуть.

Номер накладної складається з 10 символів (цифри, латинські літери).
Назва товару може мати від 4 до 23 символів включно.
Може містити тільки цифри, літери, пробіли, апостроф, дефіс, підкреслення.

Кількість, ціна мають бути додатними.
Кількість вказується з рівно 4 знаками після крапки.
Ціна вказується з рівно 3 знаками після крапки.
Вартість має бути узгоджена з кількістю та ціною та записується рівно з двома десятковими знаками.
(Застосовується округлення.)

Товари однозначно ідентифікуються назвами та одиницями виміру (за наявності).

Одиниця виміру (за наявності у варіанті) може приймати значення: шт., кг, м , кор.

Кожна накладна має свій власний унікальний номер.
У накладній можуть відвантажуватися кілька товарів (не менше 1).
Товарна позиція --- рядок накладної: товар, кількість, ціна, вартість та, за наявності, одиниця виміру.
Товарні позиції кожної накладної нумеруються послідовними натуральними числами, починаючи з 1.
В одній накладній товари можуть повторюватися, 
тобто товарні позиції однієї накладної можуть містити однакові товари 
(при цьому ціна може бути і різною, і однаковою).

\bigskip\textit{Для деяких варіантів може знадобитись}

Недоотриману вигоду розраховувати як суму виразів
{\center (найбільша ціна відвантаження товару – ціна відвантаження товару)*кількість\par}

\bigskip
{\noindent} округлених до 2 знаків після крапки.




\par
%end
\newpage
%begin
{\center Варіант  \No \textbf { 69 }\par}
\bigskip
%type = meteo
%variant = 104
Обробляються дані спостережень за погодою. 

\underline{Записи основного файлу містять поля}: кількість опадів, рік, мінімальна денна температура, код метеостанції, максимальна денна температура, місяць, день, сила вітру, відносна вологість, середня денна температура.

\underline{У допоміжному файлі наявні ключі}: найбільша різниця денних температур у спостереженнях, кількість записів.

\underline{Знайти} 4 метеостанції, по яких у середньому була найвища середньоденна температура. Вивести по кожному з них інформацію:\par
{\noindent}--- на першому рядку:\par
кількість спостережень, середня середньоденна температура найменша вологість, метеостанція;

{\noindent}--- на наступних рядках, починаючи з табуляції, вивести для них дані по датах, за які не було опадів (по одному на рядок):\par
середня денна температура, максимальна денна температура, мінімальна денна температура, сила вітру, вологість, місяць, день, рік

{\noindent}у такому сортуванні: вологість, рік, місяць, день.

%Метеостанції сортувати: найменша вологість, кількість спостережень (за спаданням), метеостанція.

\bigskip\textit{Обмеження}

Код метеостанції є рядком довжини від 3 до 7 включно.
Може містити літери алфавіту, десяткові цифри, підкреслення. Решта полів числові.

Рік є чотиризначним цілим.
Кількість опадів є дійсним значенням, подається у форматі з 1 знаком після крапки.
Сила вітру є дійсним значенням, подається у форматі з 1 знаком після крапки.
Температури є дійсними значеннями, подаються у форматі з 1 знаком після крапки.
Цілі значення подаються без десяткової крапки.
Для дійсних використовується зображення з фіксованою крапкою.

\textit{Вихідне форматування дійсних}

Усі дійсні значення виводяться у форматі з фіксованою крапкою з тією ж кількістю знаків після крапки, що і у вхідному зображенні.

Якщо у варіанті раніше не було зазначено, то \underline{обчислювані програмою} середні значення виводяться у форматі з 3 знаками після крапки.
%\bigskip%\textit{Для деяких варіантів може знадобитись}




\par
%end
\newpage
%begin
{\center Варіант  \No \textbf { 70 }\par}
\bigskip
%type = session
%variant = 72
Обробляються результати складання студентами екзаменів зимової сесії.

\underline{Записи основного файлу містять поля}: код групи, предмет, прізвище, сумарна оцінка в балах за 100-бальною системою, набрані на екзамені бали, ім’я, оцінка за державною шкалою, по-батькові, номер залікової книжки.

\underline{У допоміжному файлі наявні ключі}: загальна кількість записів, кількість оцінок 100, SMILE.

\underline{Знайти} предмети з найкращою середньою оцінкою. Вивести по кожному з них інформацію:\par
{\noindent}--- на першому рядку:\par
предмет, середня оцінка (округлений до 1 знаку), кількість студентів, що не склали;

{\noindent}--- на наступних рядках, починаючи з табуляції, вивести для предмета студентів, що отримали з нього менше 95 балів (по одному на рядок):\par
прізвище, ім’я, по-батькові, номер залікової книжки, оцінка в балах, оцінка за державною шкалою

{\noindent}у такому сортуванні: номер залікової книжки.

%Предмети сортувати: кількість студентів, що не склали, середня оцінка (округлена), предмет.

\bigskip\textit{Обмеження}

Студент не може складати предмет більше одного разу.

Бали та оцінка є значеннями цілих числових типів. Решта полів є непорожніми рядками.
Рядки розпочинатись та закінчуватись білими символами не можуть.

Оцінка за державною шкалою може приймати значення: 2-5~--- як зазвичай, 0~--- не з’явився,$-1$~--- не допущений. 
Оцінки 3-5 вважаються задовільними, решта~--- незадовільними. 
Бали є цілими та невід’ємними.
Набрані впродовж семестру бали не можуть перевищувати 60. 
Набрані на екзамені бали не можуть перевищувати 40. 
Набрані на екзамені бали повинні бути не менше 24 або дорівнювати 0(в останньому випадку оцінка не може бути задовільною). 
Сумарна оцінка в балах є сумою балів, набраних у семестрі та на екзамені.
Сумарна оцінка в балах та за державною шкалою мають бути узгоджені між собою.

Прізвище, ім’я та по-батькові (якщо наявне у варіанті) можуть мати до 28, 27, 20 відповідно символів включно кожне. 
Крім літер алфавіту можуть містити апостроф, дефіс та пробіл.

Предмет може мати від 4 до 56 символів включно. 
Крім літер алфавіту може містити подвійні лапки, пробіл та дефіс.

Код групи є непорожнім та складається не більш ніж з 3 літер. 
Крім літер алфавіту може містити десяткові цифри та дефіс.


Номер залікової книжки однозначно ідентифікує студента та складається з 7 десяткових цифр. 

\bigskip\textit{Для деяких варіантів може знадобитись}

Середній бал/оцінка за предметом визначається як середнє арифметичнесумарних балів/державних оцінок за цим предметом. 
Середній бал/оцінка студента визначається як середнє арифметичне сумарних балів/державних оцінокцього студента. 
За обчислення середньої оцінки недопущені та ті, що не з’явилися, рахуються як 2.
У решті випадків недопуски підсумовуються як $-1$, а «не з’явився» як 0.
Якщо студент не гірше ніж задовільно склав усі предмети, 
то його рейтинг за балами/оцінкою визначається як середній бал/оцінка; інакше дорівнює 0.

Стабільність студента визначається як модуль різниці між його кращим та найгіршим сумарним балом.
Стабільність предмета визначається як модуль різниці між найкращим та найгіршимсумарним балом за цим предметом.






\par
%end
\newpage
%begin
{\center Варіант  \No \textbf { 71 }\par}
\bigskip
%type = absense
%variant = 32
Обробляються дані семестрового журналу перевірок відвідування занять студентами деканатом (фіксуються відсутні). 

\underline{Записи основного файлу містять поля}: ім’я, день тижня, пара, прізвище, аудиторія, вид занять, курс, предмет, навчальний тиждень.

\underline{У допоміжному файлі наявні ключі}: загальна кількість\ студентів, кількість пропусків\ по п’ятни\-цях.

\underline{Знайти} всіх студентів, що пропускали заняття на перших парах. Вивести по кожному з них інформацію:\par
{\noindent}--- на першому рядку:\par
прізвище, ім’я, кількість пропущених пар, кількість пропущених перших пар;

{\noindent}--- на наступних рядках, починаючи з табуляції, вивести пропуски студента (по одному на рядок):\par
день, пара, предмет, вид занять, кількість пропусків

{\noindent}у такому сортуванні: день, пара, предмет, вид занять.

%Студентів сортувати: кількість пропусків перших пар(за спаданням), кількість пропусків(за спаданням), прізвище, ім’я.

\bigskip\textit{Обмеження}

Навчальний тиждень може приймати цілі значення від 1 до 20.
День тижня є цілим від 1 до 5. Пара є цілим від 1 до 4.
Курс є цілим від 1 до 4. Аудиторія є додатнім цілим.
Вид занять може приймати значення:
Le --- лекція, pr --- практичне, с. --- семінар, Lab --- лабораторне.

Решта полів є непорожніми рядками.
Рядки починатися та закінчуватися білими символами не можуть.

Групи кодуються в межах курсу.
Код групи складається не більш ніж з 4 символів.
Може містити літери алфавіту, десяткові цифри та дефіс.

Предмет є рядком, може мати від 3 до 37 символів включно.
Може містити літери алфавіту, десяткові цифри, апостроф, дефіс, пробіл.
Предмети різних курсів вважати різними навіть за однакових назв.

Прізвище, ім’я та по-батькові (за наявності у варіанті)
можуть мати від 3 до 28, 23, 28 відповідно символів включно кожне.
Можуть містити літери алфавіту, апостроф, пробіл та дефіс.
Студенти однозначно ідентифікуються прізвищем, ім’ям та по-батькові (за наявності у варіанті).
Жоден студент не вчиться одночасно на різних курсах чи в різних групах.

Види занять сортуються так: семінар < лабораторне < практичне < лекція.

%\bigskip%\textit{Для деяких варіантів може знадобитись}




\par
%end
\newpage
%begin
{\center Варіант  \No \textbf { 72 }\par}
\bigskip
%type = absense
%variant = 44
Обробляються дані семестрового журналу перевірок відвідування занять студентами деканатом (фіксуються відсутні). 

\underline{Записи основного файлу містять поля}: курс, код групи, день тижня, пара, предмет, вид занять, по-батькові, прізвище, ім’я, аудиторія, навчальний тиждень.

\underline{У допоміжному файлі наявні ключі}: загальна кількість предметів, кількість пропусків, зафіксованих на 4-му курсі.

\underline{Знайти} всі предмети, з яких пропускалися четверті пари. Вивести по кожному з них інформацію:\par
{\noindent}--- на першому рядку:\par
курс, предмет, кількість пропусків, кількість пропусків четвертих пар;

{\noindent}--- на наступних рядках, починаючи з табуляції, вивести пропуски четвертих пар (по одному на рядок):\par
прізвище, ім’я, по-батькові, день, кількість пропусків

{\noindent}у такому сортуванні: день, прізвище, ім’я, по-батькові.

%Предмети сортувати: кількість пропусків, курс, предмет.

\bigskip\textit{Обмеження}

Навчальний тиждень може приймати цілі значення від 1 до 14.
День тижня є цілим від 1 до 5. Пара є цілим від 1 до 4.
Курс є цілим від 1 до 4. Аудиторія є додатнім цілим.
Вид занять може приймати значення:
Lect --- лекція, P --- практичне, Seminar --- семінар, lab --- лабораторне.

Решта полів є непорожніми рядками.
Рядки починатися та закінчуватися білими символами не можуть.

Групи кодуються в межах курсу.
Код групи складається не більш ніж з 5 символів.
Може містити літери алфавіту, десяткові цифри та дефіс.

Предмет є рядком, може мати від 3 до 23 символів включно.
Може містити літери алфавіту, десяткові цифри, апостроф, дефіс, пробіл.
Предмети різних курсів вважати різними навіть за однакових назв.

Прізвище, ім’я та по-батькові (за наявності у варіанті)
можуть мати від 3 до 25, 22, 20 відповідно символів включно кожне.
Можуть містити літери алфавіту, апостроф, пробіл та дефіс.
Студенти однозначно ідентифікуються прізвищем, ім’ям та по-батькові (за наявності у варіанті).
Жоден студент не вчиться одночасно на різних курсах чи в різних групах.

Види занять сортуються так: семінар < лекція < лабораторне < практичне.

%\bigskip%\textit{Для деяких варіантів може знадобитись}




\par
%end
\newpage
%begin
{\center Варіант  \No \textbf { 73 }\par}
\bigskip
%type = tasks
%variant = 16
Обробляються результати розв’язування студентами задач на контрольних.
Одиничний факт роз\-в’я\-зу\-вання студентом задачі надалі розуміється як спроба.

\underline{Записи основного файлу містять поля}: відсоток набраних за розв’язання задачі балів, ім’я, рік, прізвище, умова задачі, по-батькові.

\underline{У допоміжному файлі наявні ключі}: загальна кількість студентів, найбільше значення відсотка, найменше значення відсотка.

\underline{Знайти} всіх студентів, у яких всіх спроби лежать у межах від 75\% до 90\% включно. Вивести по кожному з них інформацію:\par
{\noindent}--- на першому рядку:\par
середня розв’язуваність студента (округлена до одного знаку), ім’я, по-батькові, прізвище, кількість спроб;

{\noindent}--- на наступних рядках, починаючи з табуляції, вивести їх спроби (по одному на рядок):\par
умова, рік, відсоток набраних за задачу балів

{\noindent}у такому сортуванні: рік (за спаданням), відсоток, умова.

%Студентів сортувати: середня розв’язуваність, кількість спроб, прізвище, ім’я, по-батькові.

\bigskip\textit{Обмеження}

Студент міг розв’язувати задачу кілька разів.

Відсоток набраних за розв’язання задачі балів є цілим від 0 до 100.
День, місяць, рік є числовими полями, що зображують дату.
Раніше 2001 року статистика не велася. Рік є 4-значним.
Решта полів є непорожніми рядками.
Рядки починатися та закінчуватися білими символами не можуть.

Прізвище, ім’я та по-батькові (за наявності у варіанті) можуть мати до 21, 30, 22 відповідно символів включно кожне. 
Крім літер алфавіту можуть містити апостроф, дефіс та пробіл.

Умова задачі може мати від 5 до 59 символів включно.
Крім літер алфавіту може містити подвійні лапки, десяткові цифри, пробіл, дефіс,
а також символи . , \_ \{ \} \$ $+$ $*$ {\textbackslash} ( ) \% $[$ $]$ .

Студенти однозначно ідентифікуються прізвищем, ім’ям та по-батькові (за наявності у варіанті).
Задачі однозначно ідентифікуються умовою.

\bigskip\textit{Для деяких варіантів може знадобитись}

Розв'язуваність задачі --- середній відсоток набраних за задачу балів.

Розв'язуваність студента --- середній відсоток набраних студентом балів по всіх його спробах.

Розв'язуваність --- середній відсоток набраних студентами за задачі балів (рахується по всім даним).




\par
%end
\newpage
%begin
{\center Варіант  \No \textbf { 74 }\par}
\bigskip
%type = session
%variant = 77
Обробляються результати складання студентами екзаменів зимової сесії.

\underline{Записи основного файлу містять поля}: прізвище, по-батькові, предмет, набрані впродовж семестру бали, номер залікової книжки, ім’я, набрані на екзамені бали, оцінка за державною шкалою, код групи, сумарна оцінка в балах за 100-бальною системою.

\underline{У допоміжному файлі наявні ключі}: сума державних оцінок, загальна кількість записів.

\underline{Знайти} предмети, за якими не було жодної оцінки «добре». Вивести по кожному з них інформацію:\par
{\noindent}--- на першому рядку:\par
предмет, кількість оцінок «відмінно», середній бал за предметом (округлений до цілого);

{\noindent}--- на наступних рядках, починаючи з табуляції, вивести для предмета всіх студентів, що його складали (по одному на рядок):\par
група, кількість оцінок «відмінно», середній бал по групі (округлений до цілого)

{\noindent}у такому сортуванні: група.

%Предмети сортувати: середній бал (округлений), предмет.

\bigskip\textit{Обмеження}

Студент не може складати предмет більше одного разу.

Бали та оцінка є значеннями цілих числових типів. Решта полів є непорожніми рядками.
Рядки розпочинатись та закінчуватись білими символами не можуть.

Оцінка за державною шкалою може приймати значення: 2-5~--- як зазвичай, 0~--- не з’явився,$-1$~--- не допущений. 
Оцінки 3-5 вважаються задовільними, решта~--- незадовільними. 
Бали є цілими та невід’ємними.
Набрані впродовж семестру бали не можуть перевищувати 60. 
Набрані на екзамені бали не можуть перевищувати 40. 
Набрані на екзамені бали повинні бути не менше 24 або дорівнювати 0(в останньому випадку оцінка не може бути задовільною). 
Сумарна оцінка в балах є сумою балів, набраних у семестрі та на екзамені.
Сумарна оцінка в балах та за державною шкалою мають бути узгоджені між собою.

Прізвище, ім’я та по-батькові (якщо наявне у варіанті) можуть мати до 21, 26, 26 відповідно символів включно кожне. 
Крім літер алфавіту можуть містити апостроф, дефіс та пробіл.

Предмет може мати від 5 до 69 символів включно. 
Крім літер алфавіту може містити подвійні лапки, пробіл та дефіс.

Код групи є непорожнім та складається не більш ніж з 6 літер. 
Крім літер алфавіту може містити десяткові цифри та дефіс.


Номер залікової книжки однозначно ідентифікує студента та складається з 6 десяткових цифр. 

\bigskip\textit{Для деяких варіантів може знадобитись}

Середній бал/оцінка за предметом визначається як середнє арифметичнесумарних балів/державних оцінок за цим предметом. 
Середній бал/оцінка студента визначається як середнє арифметичне сумарних балів/державних оцінокцього студента. 
За обчислення середньої оцінки недопущені та ті, що не з’явилися, рахуються як 2.
У решті випадків недопуски підсумовуються як $-1$, а «не з’явився» як 0.
Якщо студент не гірше ніж задовільно склав усі предмети, 
то його рейтинг за балами/оцінкою визначається як середній бал/оцінка; інакше дорівнює 0.

Стабільність студента визначається як модуль різниці між його кращим та найгіршим сумарним балом.
Стабільність предмета визначається як модуль різниці між найкращим та найгіршимсумарним балом за цим предметом.






\par
%end
\newpage
%begin
{\center Варіант  \No \textbf { 75 }\par}
\bigskip
%type = tasks
%variant = 13
Обробляються результати розв’язування студентами задач на контрольних.
Одиничний факт роз\-в’я\-зу\-вання студентом задачі надалі розуміється як спроба.

\underline{Записи основного файлу містять поля}: умова задачі, відсоток набраних за розв’язання задачі балів, місяць, рік, день, ім’я, прізвище, по-батькові.

\underline{У допоміжному файлі наявні ключі}: сума днів, кількість спроб на 100\%.

\underline{Знайти} всі задачі, для яких існує спроба з відсотком розв’язання рівно 77\%. Вивести по кожному з них інформацію:\par
{\noindent}--- на першому рядку:\par
умова, кількість спроб кращих за 75\%, кількість спроб, рік, місяць, день останньої найкращої спроби;

{\noindent}--- на наступних рядках, починаючи з табуляції, вивести спроби по цих задачах (по одному на рядок):\par
прізвище, ім’я, по-батькові, відсоток набраних балів, рік, місяць, день

{\noindent}у такому сортуванні: рік, місяць, день, відсоток набраних балів, прізвище, ім’я, по-батькові.

%Задачі сортувати: кількість спроб кращих за 75\%, умова.

\bigskip\textit{Обмеження}

Студент міг розв’язувати задачу кілька разів.

Відсоток набраних за розв’язання задачі балів є цілим від 0 до 100.
День, місяць, рік є числовими полями, що зображують дату.
Раніше 2001 року статистика не велася. Рік є 4-значним.
Решта полів є непорожніми рядками.
Рядки починатися та закінчуватися білими символами не можуть.

Прізвище, ім’я та по-батькові (за наявності у варіанті) можуть мати до 25, 25, 29 відповідно символів включно кожне. 
Крім літер алфавіту можуть містити апостроф, дефіс та пробіл.

Умова задачі може мати від 4 до 50 символів включно.
Крім літер алфавіту може містити подвійні лапки, десяткові цифри, пробіл, дефіс,
а також символи . , \_ \{ \} \$ $+$ $*$ {\textbackslash} ( ) \% $[$ $]$ .

Студенти однозначно ідентифікуються прізвищем, ім’ям та по-батькові (за наявності у варіанті).
Задачі однозначно ідентифікуються умовою.

\bigskip\textit{Для деяких варіантів може знадобитись}

Розв'язуваність задачі --- середній відсоток набраних за задачу балів.

Розв'язуваність студента --- середній відсоток набраних студентом балів по всіх його спробах.

Розв'язуваність --- середній відсоток набраних студентами за задачі балів (рахується по всім даним).




\par
%end
\newpage
%begin
{\center Варіант  \No \textbf { 76 }\par}
\bigskip
%type = session
%variant = 64
Обробляються результати складання студентами екзаменів зимової сесії.

\underline{Записи основного файлу містять поля}: оцінка за державною шкалою, набрані на екзамені бали, предмет, код групи, набрані впродовж семестру бали, ім’я, номер залікової книжки, сумарна оцінка в балах за 100-бальною системою, прізвище.

\underline{У допоміжному файлі наявні ключі}: найбільший сумарний бал, найбільший бал в семестрі, кількість студентів.

\underline{Знайти} всіх студентів, чий рейтинг за балами вище середнього рейтингу. Вивести по кожному з них інформацію:\par
{\noindent}--- на першому рядку:\par
ім’я, прізвище, рейтинг за балами (округлений до цілого), номер залікової книжки, кількість «відмінно»;

{\noindent}--- на наступних рядках, починаючи з табуляції, вивести для студента всі його результати (по одному на рядок):\par
державна оцінка, предмет, сумарна оцінка

{\noindent}у такому сортуванні: державна оцінка, предмет.

%Студентів сортувати: ім’я, прізвище, рейтинг за балами (після округлення), номер залікової книжки.

\bigskip\textit{Обмеження}

Студент не може складати предмет більше одного разу.

Бали та оцінка є значеннями цілих числових типів. Решта полів є непорожніми рядками.
Рядки розпочинатись та закінчуватись білими символами не можуть.

Оцінка за державною шкалою може приймати значення: 2-5~--- як зазвичай, 0~--- не з’явився,$-1$~--- не допущений. 
Оцінки 3-5 вважаються задовільними, решта~--- незадовільними. 
Бали є цілими та невід’ємними.
Набрані впродовж семестру бали не можуть перевищувати 60. 
Набрані на екзамені бали не можуть перевищувати 40. 
Набрані на екзамені бали повинні бути не менше 24 або дорівнювати 0(в останньому випадку оцінка не може бути задовільною). 
Сумарна оцінка в балах є сумою балів, набраних у семестрі та на екзамені.
Сумарна оцінка в балах та за державною шкалою мають бути узгоджені між собою.

Прізвище, ім’я та по-батькові (якщо наявне у варіанті) можуть мати до 22, 25, 24 відповідно символів включно кожне. 
Крім літер алфавіту можуть містити апостроф, дефіс та пробіл.

Предмет може мати від 3 до 59 символів включно. 
Крім літер алфавіту може містити подвійні лапки, пробіл та дефіс.

Код групи є непорожнім та складається не більш ніж з 5 літер. 
Крім літер алфавіту може містити десяткові цифри та дефіс.


Номер залікової книжки однозначно ідентифікує студента та складається з 6 десяткових цифр. 

\bigskip\textit{Для деяких варіантів може знадобитись}

Середній бал/оцінка за предметом визначається як середнє арифметичнесумарних балів/державних оцінок за цим предметом. 
Середній бал/оцінка студента визначається як середнє арифметичне сумарних балів/державних оцінокцього студента. 
За обчислення середньої оцінки недопущені та ті, що не з’явилися, рахуються як 2.
У решті випадків недопуски підсумовуються як $-1$, а «не з’явився» як 0.
Якщо студент не гірше ніж задовільно склав усі предмети, 
то його рейтинг за балами/оцінкою визначається як середній бал/оцінка; інакше дорівнює 0.

Стабільність студента визначається як модуль різниці між його кращим та найгіршим сумарним балом.
Стабільність предмета визначається як модуль різниці між найкращим та найгіршимсумарним балом за цим предметом.






\par
%end
\newpage
%begin
{\center Варіант  \No \textbf { 77 }\par}
\bigskip
%type = tasks
%variant = -4
Обробляються результати розв’язування студентами задач на контрольних.
Одиничний факт роз\-в’я\-зу\-вання студентом задачі надалі розуміється як спроба.

\underline{Записи основного файлу містять поля}: умова задачі, по-батькові, відсоток набраних за розв’язання задачі балів, день, ім’я, прізвище, місяць, рік.

\underline{У допоміжному файлі наявні ключі}: загальна сума відсотків, сума днів.

\underline{Знайти} найкращий результат по кожній розв'язуваній задачі. Вивести по кожному з них інформацію:\par
{\noindent}--- на першому рядку:\par
умова, найкраща спроба, найгірша спроба, кількість спроб;

{\noindent}--- на наступних рядках, починаючи з табуляції, вивести студентів, що їх розв'язували (по одному на рядок):\par
прізвище, ім’я, по-батькові, відсоток набраних за задачу балів (за спаданням), рік (за спаданням), місяць(за спаданням), день (за спаданням)

{\noindent}у такому сортуванні: відсоток (за спаданням), рік (за спаданням), умова.

%Задачі сортувати: найкраща спроба, найгірша спроба, кількість спроб, умова.

\bigskip\textit{Обмеження}

Студент міг розв’язувати задачу кілька разів.

Відсоток набраних за розв’язання задачі балів є цілим від 0 до 100.
День, місяць, рік є числовими полями, що зображують дату.
Раніше 2001 року статистика не велася. Рік є 4-значним.
Решта полів є непорожніми рядками.
Рядки починатися та закінчуватися білими символами не можуть.

Прізвище, ім’я та по-батькові (за наявності у варіанті) можуть мати до 24, 24, 25 відповідно символів включно кожне. 
Крім літер алфавіту можуть містити апостроф, дефіс та пробіл.

Умова задачі може мати від 3 до 69 символів включно.
Крім літер алфавіту може містити подвійні лапки, десяткові цифри, пробіл, дефіс,
а також символи . , \_ \{ \} \$ $+$ $*$ {\textbackslash} ( ) \% $[$ $]$ .

Студенти однозначно ідентифікуються прізвищем, ім’ям та по-батькові (за наявності у варіанті).
Задачі однозначно ідентифікуються умовою.

\bigskip\textit{Для деяких варіантів може знадобитись}

Розв'язуваність задачі --- середній відсоток набраних за задачу балів.

Розв'язуваність студента --- середній відсоток набраних студентом балів по всіх його спробах.

Розв'язуваність --- середній відсоток набраних студентами за задачі балів (рахується по всім даним).




\par
%end
\newpage
%begin
{\center Варіант  \No \textbf { 78 }\par}
\bigskip
%type = absense
%variant = 30
Обробляються дані семестрового журналу перевірок відвідування занять студентами деканатом (фіксуються відсутні). 

\underline{Записи основного файлу містять поля}: навчальний тиждень, ім’я, предмет, вид занять, по-батькові, прізвище, день тижня, пара, курс, код групи, аудиторія.

\underline{У допоміжному файлі наявні ключі}: найменший номер аудиторії, сума номерів пар.

\underline{Знайти} всіх студентів, що пропустили практичних більше ніж в середньому (практичних було пропущено студентами, що пропускали практичні заняття). Вивести по кожному з них інформацію:\par
{\noindent}--- на першому рядку:\par
прізвище, ім’я по-батькові, кількість пропущених пар, кількість пропущених практичних;

{\noindent}--- на наступних рядках, починаючи з табуляції, вивести пропуски студента (по одному на рядок):\par
предмет, вид занять, тиждень, день, пара

{\noindent}у такому сортуванні: предмет, вид занять, тиждень, день, пара.

%Студентів сортувати: кількість пропусків практичних (за спаданням), кількість пропусків лекцій, прізвище, ім’я, по-батькові.

\bigskip\textit{Обмеження}

Навчальний тиждень може приймати цілі значення від 1 до 15.
День тижня є цілим від 1 до 5. Пара є цілим від 1 до 4.
Курс є цілим від 1 до 4. Аудиторія є додатнім цілим.
Вид занять може приймати значення:
lect --- лекція, Практ. --- практичне, S --- семінар, ЛАБ --- лабораторне.

Решта полів є непорожніми рядками.
Рядки починатися та закінчуватися білими символами не можуть.

Групи кодуються в межах курсу.
Код групи складається не більш ніж з 6 символів.
Може містити літери алфавіту, десяткові цифри та дефіс.

Предмет є рядком, може мати від 4 до 32 символів включно.
Може містити літери алфавіту, десяткові цифри, апостроф, дефіс, пробіл.
Предмети різних курсів вважати різними навіть за однакових назв.

Прізвище, ім’я та по-батькові (за наявності у варіанті)
можуть мати від 4 до 30, 21, 27 відповідно символів включно кожне.
Можуть містити літери алфавіту, апостроф, пробіл та дефіс.
Студенти однозначно ідентифікуються прізвищем, ім’ям та по-батькові (за наявності у варіанті).
Жоден студент не вчиться одночасно на різних курсах чи в різних групах.

Види занять сортуються так: лабораторне < практичне < семінар < лекція.

%\bigskip%\textit{Для деяких варіантів може знадобитись}




\par
%end
\newpage
%begin
{\center Варіант  \No \textbf { 79 }\par}
\bigskip
%type = session
%variant = 56
Обробляються результати складання студентами екзаменів зимової сесії.

\underline{Записи основного файлу містять поля}: прізвище, предмет, код групи, по-батькові, ім’я, оцінка за державною шкалою, номер залікової книжки, сумарна оцінка в балах за 100-бальною системою, набрані на екзамені бали.

\underline{У допоміжному файлі наявні ключі}: кількість оцінок «відмінно», кількість оцінок «задовільно», загальна кількість записів.

\underline{Знайти} всіх студентів, що щось не склали. Вивести по кожному з них інформацію:\par
{\noindent}--- на першому рядку:\par
середній бал (округлений до цілого), кількість хвостів, номер залікової книжки, прізвище, ім’я, по-батькові;

{\noindent}--- на наступних рядках, починаючи з табуляції, вивести для студента всі його результати (по одному на рядок):\par
предмет, державна оцінка прописом

{\noindent}у такому сортуванні: державна оцінка прописом (впорядковуємо як рядки), предмет.

%Студентів сортувати: кількість хвостів (за спаданням), середній бал (після округлення), номер залікової книжки.

\bigskip\textit{Обмеження}

Студент не може складати предмет більше одного разу.

Бали та оцінка є значеннями цілих числових типів. Решта полів є непорожніми рядками.
Рядки розпочинатись та закінчуватись білими символами не можуть.

Оцінка за державною шкалою може приймати значення: 2-5~--- як зазвичай, 0~--- не з’явився,$-1$~--- не допущений. 
Оцінки 3-5 вважаються задовільними, решта~--- незадовільними. 
Бали є цілими та невід’ємними.
Набрані впродовж семестру бали не можуть перевищувати 60. 
Набрані на екзамені бали не можуть перевищувати 40. 
Набрані на екзамені бали повинні бути не менше 24 або дорівнювати 0(в останньому випадку оцінка не може бути задовільною). 
Сумарна оцінка в балах є сумою балів, набраних у семестрі та на екзамені.
Сумарна оцінка в балах та за державною шкалою мають бути узгоджені між собою.

Прізвище, ім’я та по-батькові (якщо наявне у варіанті) можуть мати до 20, 20, 27 відповідно символів включно кожне. 
Крім літер алфавіту можуть містити апостроф, дефіс та пробіл.

Предмет може мати від 4 до 57 символів включно. 
Крім літер алфавіту може містити подвійні лапки, пробіл та дефіс.

Код групи є непорожнім та складається не більш ніж з 7 літер. 
Крім літер алфавіту може містити десяткові цифри та дефіс.


Номер залікової книжки однозначно ідентифікує студента та складається з 8 десяткових цифр. 

\bigskip\textit{Для деяких варіантів може знадобитись}

Середній бал/оцінка за предметом визначається як середнє арифметичнесумарних балів/державних оцінок за цим предметом. 
Середній бал/оцінка студента визначається як середнє арифметичне сумарних балів/державних оцінокцього студента. 
За обчислення середньої оцінки недопущені та ті, що не з’явилися, рахуються як 2.
У решті випадків недопуски підсумовуються як $-1$, а «не з’явився» як 0.
Якщо студент не гірше ніж задовільно склав усі предмети, 
то його рейтинг за балами/оцінкою визначається як середній бал/оцінка; інакше дорівнює 0.

Стабільність студента визначається як модуль різниці між його кращим та найгіршим сумарним балом.
Стабільність предмета визначається як модуль різниці між найкращим та найгіршимсумарним балом за цим предметом.






\par
%end
\newpage
%begin
{\center Варіант  \No \textbf { 80 }\par}
\bigskip
%type = absense
%variant = 29
Обробляються дані семестрового журналу перевірок відвідування занять студентами деканатом (фіксуються відсутні). 

\underline{Записи основного файлу містять поля}: предмет, прізвище, день тижня, пара, аудиторія, вид занять, навчальний тиждень, курс, код групи, ім’я.

\underline{У допоміжному файлі наявні ключі}: найбільший номер аудиторії, кількість записів у файлі.

\underline{Знайти} всіх студентів, що хоча б з одного предмета пропустили 15 лекцій. Вивести по кожному з них інформацію:\par
{\noindent}--- на першому рядку:\par
кількість пропущених пар, кількість пропущених лекцій прізвище, ім’я;

{\noindent}--- на наступних рядках, починаючи з табуляції, вивести пропуски студента (по одному на рядок):\par
предмет, вид занять, кількість пропусків

{\noindent}у такому сортуванні: предмет, вид занять.

%Студентів сортувати: кількість пропусків, прізвище, ім’я.

\bigskip\textit{Обмеження}

Навчальний тиждень може приймати цілі значення від 1 до 18.
День тижня є цілим від 1 до 5. Пара є цілим від 1 до 4.
Курс є цілим від 1 до 4. Аудиторія є додатнім цілим.
Вид занять може приймати значення:
Lecture --- лекція, практ. --- практичне, s --- семінар, Лаб. --- лабораторне.

Решта полів є непорожніми рядками.
Рядки починатися та закінчуватися білими символами не можуть.

Групи кодуються в межах курсу.
Код групи складається не більш ніж з 4 символів.
Може містити літери алфавіту, десяткові цифри та дефіс.

Предмет є рядком, може мати від 6 до 24 символів включно.
Може містити літери алфавіту, десяткові цифри, апостроф, дефіс, пробіл.
Предмети різних курсів вважати різними навіть за однакових назв.

Прізвище, ім’я та по-батькові (за наявності у варіанті)
можуть мати від 6 до 20, 20, 24 відповідно символів включно кожне.
Можуть містити літери алфавіту, апостроф, пробіл та дефіс.
Студенти однозначно ідентифікуються прізвищем, ім’ям та по-батькові (за наявності у варіанті).
Жоден студент не вчиться одночасно на різних курсах чи в різних групах.

Види занять сортуються так: семінар < лекція < лабораторне < практичне.

%\bigskip%\textit{Для деяких варіантів може знадобитись}




\par
%end
\newpage
%begin
{\center Варіант  \No \textbf { 81 }\par}
\bigskip
%type = session
%variant = 59
Обробляються результати складання студентами екзаменів зимової сесії.

\underline{Записи основного файлу містять поля}: сумарна оцінка в балах за 100-бальною системою, набрані впродовж семестру бали, номер залікової книжки, прізвище, код групи, по-батькові, предмет, ім’я, набрані на екзамені бали, оцінка за державною шкалою.

\underline{У допоміжному файлі наявні ключі}: загальна кількість записів, кількість оцінок 77 балів, кількість оцінок 99 балів.

\underline{Знайти} всіх студентів, що хоча б з одного предмету отримали оцінку 5 за державною шкалою. Вивести по кожному з них інформацію:\par
{\noindent}--- на першому рядку:\par
прізвище, ім’я, по-батькові, номер залікової книжки, кількість «відмінно», кількість «добре», кількість «задовільно», кількість екзаменів;

{\noindent}--- на наступних рядках, починаючи з табуляції, вивести для студента всі його результати (по одному на рядок):\par
сумарна оцінка, предмет, бали на екзамені

{\noindent}у такому сортуванні: бали на екзамені, сумарна оцінка, предмет.

%Студентів сортувати: кількість екзаменів, прізвище, ім’я, по-батькові, номер залікової книжки.

\bigskip\textit{Обмеження}

Студент не може складати предмет більше одного разу.

Бали та оцінка є значеннями цілих числових типів. Решта полів є непорожніми рядками.
Рядки розпочинатись та закінчуватись білими символами не можуть.

Оцінка за державною шкалою може приймати значення: 2-5~--- як зазвичай, 0~--- не з’явився,$-1$~--- не допущений. 
Оцінки 3-5 вважаються задовільними, решта~--- незадовільними. 
Бали є цілими та невід’ємними.
Набрані впродовж семестру бали не можуть перевищувати 60. 
Набрані на екзамені бали не можуть перевищувати 40. 
Набрані на екзамені бали повинні бути не менше 24 або дорівнювати 0(в останньому випадку оцінка не може бути задовільною). 
Сумарна оцінка в балах є сумою балів, набраних у семестрі та на екзамені.
Сумарна оцінка в балах та за державною шкалою мають бути узгоджені між собою.

Прізвище, ім’я та по-батькові (якщо наявне у варіанті) можуть мати до 29, 22, 23 відповідно символів включно кожне. 
Крім літер алфавіту можуть містити апостроф, дефіс та пробіл.

Предмет може мати від 5 до 60 символів включно. 
Крім літер алфавіту може містити подвійні лапки, пробіл та дефіс.

Код групи є непорожнім та складається не більш ніж з 6 літер. 
Крім літер алфавіту може містити десяткові цифри та дефіс.


Номер залікової книжки однозначно ідентифікує студента та складається з 7 десяткових цифр. 

\bigskip\textit{Для деяких варіантів може знадобитись}

Середній бал/оцінка за предметом визначається як середнє арифметичнесумарних балів/державних оцінок за цим предметом. 
Середній бал/оцінка студента визначається як середнє арифметичне сумарних балів/державних оцінокцього студента. 
За обчислення середньої оцінки недопущені та ті, що не з’явилися, рахуються як 2.
У решті випадків недопуски підсумовуються як $-1$, а «не з’явився» як 0.
Якщо студент не гірше ніж задовільно склав усі предмети, 
то його рейтинг за балами/оцінкою визначається як середній бал/оцінка; інакше дорівнює 0.

Стабільність студента визначається як модуль різниці між його кращим та найгіршим сумарним балом.
Стабільність предмета визначається як модуль різниці між найкращим та найгіршимсумарним балом за цим предметом.






\par
%end
\newpage
%begin
{\center Варіант  \No \textbf { 82 }\par}
\bigskip
%type = absense
%variant = 33
Обробляються дані семестрового журналу перевірок відвідування занять студентами деканатом (фіксуються відсутні). 

\underline{Записи основного файлу містять поля}: по-батькові, прізвище, день тижня, пара, вид занять, курс, код групи, аудиторія, навчальний тиждень, ім’я, предмет.

\underline{У допоміжному файлі наявні ключі}: загальна кількість записів, довжина найкоротшого прізвища.

\underline{Знайти} всіх студентів, що пропускали заняття принаймні з двох предметів. Вивести по кожному з них інформацію:\par
{\noindent}--- на першому рядку:\par
прізвище, ім’я, по-батькові, кількість пропущених лекцій;

{\noindent}--- на наступних рядках, починаючи з табуляції, вивести пропуски студента (по одному на рядок):\par
предмет, вид занять, кількість пропусків

{\noindent}у такому сортуванні: предмет, вид занять.

%Студентів сортувати: кількість пропусків лекцій (за спаданням), прізвище, ім’я, по-батькові.

\bigskip\textit{Обмеження}

Навчальний тиждень може приймати цілі значення від 1 до 20.
День тижня є цілим від 1 до 5. Пара є цілим від 1 до 4.
Курс є цілим від 1 до 4. Аудиторія є додатнім цілим.
Вид занять може приймати значення:
лекц. --- лекція, ПР --- практичне, sem --- семінар, лаб. --- лабораторне.

Решта полів є непорожніми рядками.
Рядки починатися та закінчуватися білими символами не можуть.

Групи кодуються в межах курсу.
Код групи складається не більш ніж з 5 символів.
Може містити літери алфавіту, десяткові цифри та дефіс.

Предмет є рядком, може мати від 5 до 26 символів включно.
Може містити літери алфавіту, десяткові цифри, апостроф, дефіс, пробіл.
Предмети різних курсів вважати різними навіть за однакових назв.

Прізвище, ім’я та по-батькові (за наявності у варіанті)
можуть мати від 5 до 24, 27, 29 відповідно символів включно кожне.
Можуть містити літери алфавіту, апостроф, пробіл та дефіс.
Студенти однозначно ідентифікуються прізвищем, ім’ям та по-батькові (за наявності у варіанті).
Жоден студент не вчиться одночасно на різних курсах чи в різних групах.

Види занять сортуються так: семінар < лабораторне < практичне < лекція.

%\bigskip%\textit{Для деяких варіантів може знадобитись}




\par
%end
\newpage
%begin
{\center Варіант  \No \textbf { 83 }\par}
\bigskip
%type = session
%variant = 68
Обробляються результати складання студентами екзаменів зимової сесії.

\underline{Записи основного файлу містять поля}: набрані на екзамені бали, оцінка за державною шкалою, ім’я, предмет, набрані впродовж семестру бали, сумарна оцінка в балах за 100-бальною системою, код групи, прізвище, номер залікової книжки.

\underline{У допоміжному файлі наявні ключі}: сума державних оцінок, кількість записів.

\underline{Знайти} всіх студентів, що не склали жодного предмета. Вивести по кожному з них інформацію:\par
{\noindent}--- на першому рядку:\par
ім’я, прізвище, номер залікової книжки, кількість «хвостів», кількість «незадовільно»;

{\noindent}--- на наступних рядках, починаючи з табуляції, вивести для студента всі його результати (по одному на рядок):\par
предмет, бали в семестрі, бали на екзамені, сумарна оцінка, державна оцінка

{\noindent}у такому сортуванні: предмет.

%Студентів сортувати: кількість «хвостів» (за спаданням), прізвище, номер залікової книжки.

\bigskip\textit{Обмеження}

Студент не може складати предмет більше одного разу.

Бали та оцінка є значеннями цілих числових типів. Решта полів є непорожніми рядками.
Рядки розпочинатись та закінчуватись білими символами не можуть.

Оцінка за державною шкалою може приймати значення: 2-5~--- як зазвичай, 0~--- не з’явився,$-1$~--- не допущений. 
Оцінки 3-5 вважаються задовільними, решта~--- незадовільними. 
Бали є цілими та невід’ємними.
Набрані впродовж семестру бали не можуть перевищувати 60. 
Набрані на екзамені бали не можуть перевищувати 40. 
Набрані на екзамені бали повинні бути не менше 24 або дорівнювати 0(в останньому випадку оцінка не може бути задовільною). 
Сумарна оцінка в балах є сумою балів, набраних у семестрі та на екзамені.
Сумарна оцінка в балах та за державною шкалою мають бути узгоджені між собою.

Прізвище, ім’я та по-батькові (якщо наявне у варіанті) можуть мати до 27, 23, 21 відповідно символів включно кожне. 
Крім літер алфавіту можуть містити апостроф, дефіс та пробіл.

Предмет може мати від 3 до 65 символів включно. 
Крім літер алфавіту може містити подвійні лапки, пробіл та дефіс.

Код групи є непорожнім та складається не більш ніж з 3 літер. 
Крім літер алфавіту може містити десяткові цифри та дефіс.


Номер залікової книжки однозначно ідентифікує студента та складається з 8 десяткових цифр. 

\bigskip\textit{Для деяких варіантів може знадобитись}

Середній бал/оцінка за предметом визначається як середнє арифметичнесумарних балів/державних оцінок за цим предметом. 
Середній бал/оцінка студента визначається як середнє арифметичне сумарних балів/державних оцінокцього студента. 
За обчислення середньої оцінки недопущені та ті, що не з’явилися, рахуються як 2.
У решті випадків недопуски підсумовуються як $-1$, а «не з’явився» як 0.
Якщо студент не гірше ніж задовільно склав усі предмети, 
то його рейтинг за балами/оцінкою визначається як середній бал/оцінка; інакше дорівнює 0.

Стабільність студента визначається як модуль різниці між його кращим та найгіршим сумарним балом.
Стабільність предмета визначається як модуль різниці між найкращим та найгіршимсумарним балом за цим предметом.






\par
%end
\newpage
%begin
{\center Варіант  \No \textbf { 84 }\par}
\bigskip
%type = session
%variant = 75
Обробляються результати складання студентами екзаменів зимової сесії.

\underline{Записи основного файлу містять поля}: оцінка за державною шкалою, ім’я, код групи, сумарна оцінка в балах за 100-бальною системою, прізвище, номер залікової книжки, предмет, набрані на екзамені бали, по-батькові.

\underline{У допоміжному файлі наявні ключі}: загальна кількість записів, сума балів.

\underline{Знайти} предмети, середня оцінка за які, вища середньої за всю сесію. Вивести по кожному з них інформацію:\par
{\noindent}--- на першому рядку:\par
середня оцінка (округлена до одного знаку), кількість студентів, що складали, предмет;

{\noindent}--- на наступних рядках, починаючи з табуляції, вивести для предмета всіх студентів, що його складали (по одному на рядок):\par
група, середня оцінка по групі (округлена до цілого), кількість студентів

{\noindent}у такому сортуванні: група.

%Предмети сортувати: середня оцінка (округлена до цілого), предмет.

\bigskip\textit{Обмеження}

Студент не може складати предмет більше одного разу.

Бали та оцінка є значеннями цілих числових типів. Решта полів є непорожніми рядками.
Рядки розпочинатись та закінчуватись білими символами не можуть.

Оцінка за державною шкалою може приймати значення: 2-5~--- як зазвичай, 0~--- не з’явився,$-1$~--- не допущений. 
Оцінки 3-5 вважаються задовільними, решта~--- незадовільними. 
Бали є цілими та невід’ємними.
Набрані впродовж семестру бали не можуть перевищувати 60. 
Набрані на екзамені бали не можуть перевищувати 40. 
Набрані на екзамені бали повинні бути не менше 24 або дорівнювати 0(в останньому випадку оцінка не може бути задовільною). 
Сумарна оцінка в балах є сумою балів, набраних у семестрі та на екзамені.
Сумарна оцінка в балах та за державною шкалою мають бути узгоджені між собою.

Прізвище, ім’я та по-батькові (якщо наявне у варіанті) можуть мати до 26, 29, 22 відповідно символів включно кожне. 
Крім літер алфавіту можуть містити апостроф, дефіс та пробіл.

Предмет може мати від 4 до 63 символів включно. 
Крім літер алфавіту може містити подвійні лапки, пробіл та дефіс.

Код групи є непорожнім та складається не більш ніж з 4 літер. 
Крім літер алфавіту може містити десяткові цифри та дефіс.


Номер залікової книжки однозначно ідентифікує студента та складається з 7 десяткових цифр. 

\bigskip\textit{Для деяких варіантів може знадобитись}

Середній бал/оцінка за предметом визначається як середнє арифметичнесумарних балів/державних оцінок за цим предметом. 
Середній бал/оцінка студента визначається як середнє арифметичне сумарних балів/державних оцінокцього студента. 
За обчислення середньої оцінки недопущені та ті, що не з’явилися, рахуються як 2.
У решті випадків недопуски підсумовуються як $-1$, а «не з’явився» як 0.
Якщо студент не гірше ніж задовільно склав усі предмети, 
то його рейтинг за балами/оцінкою визначається як середній бал/оцінка; інакше дорівнює 0.

Стабільність студента визначається як модуль різниці між його кращим та найгіршим сумарним балом.
Стабільність предмета визначається як модуль різниці між найкращим та найгіршимсумарним балом за цим предметом.






\par
%end
\newpage
%begin
{\center Варіант  \No \textbf { 85 }\par}
\bigskip
%type = meteo
%variant = 100
Обробляються дані спостережень за погодою. 

\underline{Записи основного файлу містять поля}: місяць, день, сила вітру, середня денна температура, код метеостанції, кількість опадів, максимальна денна температура, мінімальна денна температура, рік, відносна вологість.

\underline{У допоміжному файлі наявні ключі}: кількість днів, в які велися спостереження; загальна кількість записів.

\underline{Знайти} дати, по яких найбільша середньоденна температура була більша ніж в середньому по всіх спостереженнях. Вивести по кожному з них інформацію:\par
{\noindent}--- на першому рядку:\par
день, рік, місяць, кількість опадів, середня середньоденна температура, найбільша середньоденна температура;

{\noindent}--- на наступних рядках, починаючи з табуляції, вивести для дати спостереження метеостанції, по яких денна температура відрізняється від середньої середньоденної принаймні на 10 (по одному на рядок):\par
мінімальна денна температура, максимальна денна температура, середня денна температура, вологість, метеостанція

{\noindent}у такому сортуванні: різниця максимальної та мінімальної температур, метеостанція.

%Дати сортувати: кількість опадів, місяць, рік, день.

\bigskip\textit{Обмеження}

Код метеостанції є рядком довжини від 5 до 9 включно.
Може містити літери алфавіту, десяткові цифри, підкреслення. Решта полів числові.

Рік є чотиризначним цілим.
Кількість опадів є цілим значенням.
Сила вітру є дійсним значенням, подається у форматі з 2 знаками після крапки.
Температури є дійсними значеннями, подаються у форматі з 2 знаками після крапки.
Цілі значення подаються без десяткової крапки.
Для дійсних використовується зображення з фіксованою крапкою.

\textit{Вихідне форматування дійсних}

Усі дійсні значення виводяться у форматі з фіксованою крапкою з тією ж кількістю знаків після крапки, що і у вхідному зображенні.

Якщо у варіанті раніше не було зазначено, то \underline{обчислювані програмою} середні значення виводяться у форматі з 2 знаками після крапки.
%\bigskip%\textit{Для деяких варіантів може знадобитись}




\par
%end
\newpage
%begin
{\center Варіант  \No \textbf { 86 }\par}
\bigskip
%type = absense
%variant = 25
Обробляються дані семестрового журналу перевірок відвідування занять студентами деканатом (фіксуються відсутні). 

\underline{Записи основного файлу містять поля}: навчальний тиждень, курс, день тижня, пара, аудиторія, ім’я, вид занять, прізвище, предмет.

\underline{У допоміжному файлі наявні ключі}: суму номерів пар, довжину найкоротшого предмета.

\underline{Знайти} всіх студентів, що пропускали пари тільки перші 3 навчальних тижні. Вивести по кожному з них інформацію:\par
{\noindent}--- на першому рядку:\par
прізвище, ім’я, кількість пропущених пар, кількість пропущених лабораторних;

{\noindent}--- на наступних рядках, починаючи з табуляції, вивести пропуски студента (по одному на рядок):\par
тиждень, день, пара, аудиторія, предмет, вид занять

{\noindent}у такому сортуванні: тиждень, день, пара, предмет, вид занять.

%Студентів сортувати: останній тиждень, на якому були пропуски (за спаданням), прізвище, ім’я.

\bigskip\textit{Обмеження}

Навчальний тиждень може приймати цілі значення від 1 до 17.
День тижня є цілим від 1 до 5. Пара є цілим від 1 до 4.
Курс є цілим від 1 до 4. Аудиторія є додатнім цілим.
Вид занять може приймати значення:
ЛЕ --- лекція, Pr --- практичне, Сем. --- семінар, ЛАБ --- лабораторне.

Решта полів є непорожніми рядками.
Рядки починатися та закінчуватися білими символами не можуть.

Групи кодуються в межах курсу.
Код групи складається не більш ніж з 4 символів.
Може містити літери алфавіту, десяткові цифри та дефіс.

Предмет є рядком, може мати від 5 до 40 символів включно.
Може містити літери алфавіту, десяткові цифри, апостроф, дефіс, пробіл.
Предмети різних курсів вважати різними навіть за однакових назв.

Прізвище, ім’я та по-батькові (за наявності у варіанті)
можуть мати від 5 до 23, 26, 23 відповідно символів включно кожне.
Можуть містити літери алфавіту, апостроф, пробіл та дефіс.
Студенти однозначно ідентифікуються прізвищем, ім’ям та по-батькові (за наявності у варіанті).
Жоден студент не вчиться одночасно на різних курсах чи в різних групах.

Види занять сортуються так: семінар < практичне < лекція < лабораторне.

%\bigskip%\textit{Для деяких варіантів може знадобитись}




\par
%end
\newpage
%begin
{\center Варіант  \No \textbf { 87 }\par}
\bigskip
%type = meteo
%variant = 96
Обробляються дані спостережень за погодою. 

\underline{Записи основного файлу містять поля}: рік, середня денна температура, мінімальна денна температура, відносна вологість, місяць, день, кількість опадів, код метеостанції, максимальна денна температура, сила вітру.

\underline{У допоміжному файлі наявні ключі}: найбільша зафіксована кількість опадів, загальна кількість дат, коли фіксувалися спостереження.

\underline{Знайти} дати, по яких фіксувалася максимальна денна температура менша ніж в середньому по всіх спостереженнях. Вивести по кожному з них інформацію:\par
{\noindent}--- на першому рядку:\par
день, місяць, рік, найбільша та найменша вологість за дату, середня кількість опадів;

{\noindent}--- на наступних рядках, починаючи з табуляції, вивести для дати спостереження метеостанції, де вологість була більша 60 за відсутності опадів (по одному на рядок):\par
вологість, середня температура, метеостанція

{\noindent}у такому сортуванні: вологість, метеостанція.

%Дати сортувати: місяць, середня кількість опадів, рік, день.

\bigskip\textit{Обмеження}

Код метеостанції є рядком довжини від 5 до 11 включно.
Може містити літери алфавіту, десяткові цифри, підкреслення. Решта полів числові.

Рік є чотиризначним цілим.
Кількість опадів є цілим значенням.
Сила вітру є дійсним значенням, подається у форматі з 2 знаками після крапки.
Температури є дійсними значеннями, подаються у форматі з 1 знаком після крапки.
Цілі значення подаються без десяткової крапки.
Для дійсних використовується зображення з фіксованою крапкою.

\textit{Вихідне форматування дійсних}

Усі дійсні значення виводяться у форматі з фіксованою крапкою з тією ж кількістю знаків після крапки, що і у вхідному зображенні.

Якщо у варіанті раніше не було зазначено, то \underline{обчислювані програмою} середні значення виводяться у форматі з 1 знаком після крапки.
%\bigskip%\textit{Для деяких варіантів може знадобитись}




\par
%end
\newpage
%begin
{\center Варіант  \No \textbf { 88 }\par}
\bigskip
%type = absense
%variant = 26
Обробляються дані семестрового журналу перевірок відвідування занять студентами деканатом (фіксуються відсутні). 

\underline{Записи основного файлу містять поля}: предмет, навчальний тиждень, вид занять, ім’я, прізвище, курс, код групи, день тижня, пара, по-батькові, аудиторія.

\underline{У допоміжному файлі наявні ключі}: найбільша кількість пропусків одного студента, загальна кількість записів.

\underline{Знайти} всіх студентів, що кожного навчального тижня пропускали хоча б по одній парі. Вивести по кожному з них інформацію:\par
{\noindent}--- на першому рядку:\par
кількість пропущених пар, прізвище, ім’я по-батькові, середня кількість пропущених пар на тиждень (округлена до 1 знаку);

{\noindent}--- на наступних рядках, починаючи з табуляції, вивести пропуски студента (по одному на рядок):\par
тиждень, предмет, кількість пропусків предмета на цьому тижні

{\noindent}у такому сортуванні: тиждень, предмет.

%Студентів сортувати: кількість пропусків(за спаданням), кількість пропусків лекцій, прізвище, ім’я, по-батькові.

\bigskip\textit{Обмеження}

Навчальний тиждень може приймати цілі значення від 1 до 19.
День тижня є цілим від 1 до 5. Пара є цілим від 1 до 4.
Курс є цілим від 1 до 4. Аудиторія є додатнім цілим.
Вид занять може приймати значення:
Лекц. --- лекція, Пр --- практичне, s --- семінар, Lab --- лабораторне.

Решта полів є непорожніми рядками.
Рядки починатися та закінчуватися білими символами не можуть.

Групи кодуються в межах курсу.
Код групи складається не більш ніж з 6 символів.
Може містити літери алфавіту, десяткові цифри та дефіс.

Предмет є рядком, може мати від 4 до 36 символів включно.
Може містити літери алфавіту, десяткові цифри, апостроф, дефіс, пробіл.
Предмети різних курсів вважати різними навіть за однакових назв.

Прізвище, ім’я та по-батькові (за наявності у варіанті)
можуть мати від 4 до 22, 28, 21 відповідно символів включно кожне.
Можуть містити літери алфавіту, апостроф, пробіл та дефіс.
Студенти однозначно ідентифікуються прізвищем, ім’ям та по-батькові (за наявності у варіанті).
Жоден студент не вчиться одночасно на різних курсах чи в різних групах.

Види занять сортуються так: практичне < лекція < семінар < лабораторне.

%\bigskip%\textit{Для деяких варіантів може знадобитись}




\par
%end
\newpage
%begin
{\center Варіант  \No \textbf { 89 }\par}
\bigskip
%type = absense
%variant = 43
Обробляються дані семестрового журналу перевірок відвідування занять студентами деканатом (фіксуються відсутні). 

\underline{Записи основного файлу містять поля}: день тижня, пара, прізвище, вид занять, предмет, курс, навчальний тиждень, аудиторія, ім’я.

\underline{У допоміжному файлі наявні ключі}: найбільша довжина прізвища, найменша довжина прізвища, середнє значення курсу з округленням до 1 знаку.

\underline{Знайти} всі предмети, з яких пропускалися тільки перші пари. Вивести по кожному з них інформацію:\par
{\noindent}--- на першому рядку:\par
курс, предмет, кількість пропусків;

{\noindent}--- на наступних рядках, починаючи з табуляції, вивести пропуски предмета (по одному на рядок):\par
день, пара, прізвище, ім’я, кількість пропусків

{\noindent}у такому сортуванні: прізвище, ім’я, день, пара.

%Предмети сортувати: кількість пропусків, курс, предмет.

\bigskip\textit{Обмеження}

Навчальний тиждень може приймати цілі значення від 1 до 16.
День тижня є цілим від 1 до 5. Пара є цілим від 1 до 4.
Курс є цілим від 1 до 4. Аудиторія є додатнім цілим.
Вид занять може приймати значення:
Л --- лекція, п --- практичне, с. --- семінар, lab --- лабораторне.

Решта полів є непорожніми рядками.
Рядки починатися та закінчуватися білими символами не можуть.

Групи кодуються в межах курсу.
Код групи складається не більш ніж з 4 символів.
Може містити літери алфавіту, десяткові цифри та дефіс.

Предмет є рядком, може мати від 6 до 22 символів включно.
Може містити літери алфавіту, десяткові цифри, апостроф, дефіс, пробіл.
Предмети різних курсів вважати різними навіть за однакових назв.

Прізвище, ім’я та по-батькові (за наявності у варіанті)
можуть мати від 6 до 26, 25, 26 відповідно символів включно кожне.
Можуть містити літери алфавіту, апостроф, пробіл та дефіс.
Студенти однозначно ідентифікуються прізвищем, ім’ям та по-батькові (за наявності у варіанті).
Жоден студент не вчиться одночасно на різних курсах чи в різних групах.

Види занять сортуються так: практичне < лабораторне < семінар < лекція.

%\bigskip%\textit{Для деяких варіантів може знадобитись}




\par
%end
\newpage
%begin
{\center Варіант  \No \textbf { 90 }\par}
\bigskip
%type = absense
%variant = 35
Обробляються дані семестрового журналу перевірок відвідування занять студентами деканатом (фіксуються відсутні). 

\underline{Записи основного файлу містять поля}: курс, код групи, ім’я, прізвище, навчальний тиждень, аудиторія, вид занять, предмет, день тижня, пара.

\underline{У допоміжному файлі наявні ключі}: BOOK, загальна кількість записів.

\underline{Знайти} всіх студентів, що пропускали заняття тільки по п’ятницях. Вивести по кожному з них інформацію:\par
{\noindent}--- на першому рядку:\par
кількість пропусків, прізвище, ім’я;

{\noindent}--- на наступних рядках, починаючи з табуляції, вивести пропуски студента (по одному на рядок):\par
день, пара, загальна кількість пропусків

{\noindent}у такому сортуванні: день, пара.

%Студентів сортувати: кількість пропусків, прізвище, ім’я.

\bigskip\textit{Обмеження}

Навчальний тиждень може приймати цілі значення від 1 до 15.
День тижня є цілим від 1 до 5. Пара є цілим від 1 до 4.
Курс є цілим від 1 до 4. Аудиторія є додатнім цілим.
Вид занять може приймати значення:
L --- лекція, p --- практичне, sem --- семінар, Лаб. --- лабораторне.

Решта полів є непорожніми рядками.
Рядки починатися та закінчуватися білими символами не можуть.

Групи кодуються в межах курсу.
Код групи складається не більш ніж з 5 символів.
Може містити літери алфавіту, десяткові цифри та дефіс.

Предмет є рядком, може мати від 3 до 20 символів включно.
Може містити літери алфавіту, десяткові цифри, апостроф, дефіс, пробіл.
Предмети різних курсів вважати різними навіть за однакових назв.

Прізвище, ім’я та по-батькові (за наявності у варіанті)
можуть мати від 3 до 27, 30, 30 відповідно символів включно кожне.
Можуть містити літери алфавіту, апостроф, пробіл та дефіс.
Студенти однозначно ідентифікуються прізвищем, ім’ям та по-батькові (за наявності у варіанті).
Жоден студент не вчиться одночасно на різних курсах чи в різних групах.

Види занять сортуються так: лабораторне < лекція < практичне < семінар.

%\bigskip%\textit{Для деяких варіантів може знадобитись}




\par
%end
\newpage
%begin
{\center Варіант  \No \textbf { 91 }\par}
\bigskip
%type = tasks
%variant = 10
Обробляються результати розв’язування студентами задач на контрольних.
Одиничний факт роз\-в’я\-зу\-вання студентом задачі надалі розуміється як спроба.

\underline{Записи основного файлу містять поля}: по-батькові, рік, умова задачі, відсоток набраних за розв’я\-зан\-ня задачі балів, ім’я, прізвище.

\underline{У допоміжному файлі наявні ключі}: найменший відсоток набраних балів; середній відсоток набраних балів, округлений до цілого.

\underline{Знайти} всі задачі, які хоча б раз були розв’язані не менш ніж на 90\%. Вивести по кожному з них інформацію:\par
{\noindent}--- на першому рядку:\par
середня розв’язуваність (округлена до 1 знаку), умова, кількість спроб;

{\noindent}--- на наступних рядках, починаючи з табуляції, вивести спроби по цих задачах з відсотком не більше 75\% (по одному на рядок):\par
прізвище, ім’я, по-батькові, відсоток набраних балів, рік

{\noindent}у такому сортуванні: відсоток набраних балів (за спаданням), рік, прізвище, ім’я, по-батькові.

%Задачі сортувати: середня розв’язуваність, умова.

\bigskip\textit{Обмеження}

Студент міг розв’язувати задачу кілька разів.

Відсоток набраних за розв’язання задачі балів є цілим від 0 до 100.
День, місяць, рік є числовими полями, що зображують дату.
Раніше 2001 року статистика не велася. Рік є 4-значним.
Решта полів є непорожніми рядками.
Рядки починатися та закінчуватися білими символами не можуть.

Прізвище, ім’я та по-батькові (за наявності у варіанті) можуть мати до 27, 28, 30 відповідно символів включно кожне. 
Крім літер алфавіту можуть містити апостроф, дефіс та пробіл.

Умова задачі може мати від 4 до 53 символів включно.
Крім літер алфавіту може містити подвійні лапки, десяткові цифри, пробіл, дефіс,
а також символи . , \_ \{ \} \$ $+$ $*$ {\textbackslash} ( ) \% $[$ $]$ .

Студенти однозначно ідентифікуються прізвищем, ім’ям та по-батькові (за наявності у варіанті).
Задачі однозначно ідентифікуються умовою.

\bigskip\textit{Для деяких варіантів може знадобитись}

Розв'язуваність задачі --- середній відсоток набраних за задачу балів.

Розв'язуваність студента --- середній відсоток набраних студентом балів по всіх його спробах.

Розв'язуваність --- середній відсоток набраних студентами за задачі балів (рахується по всім даним).




\par
%end
\newpage
%begin
{\center Варіант  \No \textbf { 92 }\par}
\bigskip
%type = meteo
%variant = 94
Обробляються дані спостережень за погодою. 

\underline{Записи основного файлу містять поля}: максимальна денна температура, середня денна температура, мінімальна денна температура, рік, сила вітру, кількість опадів, код метеостанції, відносна вологість, день, місяць.

\underline{У допоміжному файлі наявні ключі}: найбільша максимальна температура, найменша максимальна температура, SMILE.

\underline{Знайти} дати, по яких фіксувалася кількість опадів, більша ніж в середньому по всіх спостереженнях. Вивести по кожному з них інформацію:\par
{\noindent}--- на першому рядку:\par
рік, місяць, день, найбільша кількість опадів за цю дату, середня сила вітру;

{\noindent}--- на наступних рядках, починаючи з табуляції, вивести для дати спостереження метеостанції, коди яких починаються з 13 (по одному на рядок):\par
метеостанція, вологість, середня денна температура

{\noindent}у такому сортуванні: метеостанція.

%Дати сортувати: середня сила вітру, кількість опадів, рік, місяць, день.

\bigskip\textit{Обмеження}

Код метеостанції є рядком довжини від 3 до 6 включно.
Може містити літери алфавіту, десяткові цифри, підкреслення. Решта полів числові.

Рік є чотиризначним цілим.
Кількість опадів є дійсним значенням, подається у форматі з 1 знаком після крапки.
Сила вітру є дійсним значенням, подається у форматі з 1 знаком після крапки.
Температури є дійсними значеннями, подаються у форматі з 3 знаками після крапки.
Цілі значення подаються без десяткової крапки.
Для дійсних використовується зображення з фіксованою крапкою.

\textit{Вихідне форматування дійсних}

Усі дійсні значення виводяться у форматі з фіксованою крапкою з тією ж кількістю знаків після крапки, що і у вхідному зображенні.

Якщо у варіанті раніше не було зазначено, то \underline{обчислювані програмою} середні значення виводяться у форматі з 3 знаками після крапки.
%\bigskip%\textit{Для деяких варіантів може знадобитись}




\par
%end
\newpage
%begin
{\center Варіант  \No \textbf { 93 }\par}
\bigskip
%type = session
%variant = 62
Обробляються результати складання студентами екзаменів зимової сесії.

\underline{Записи основного файлу містять поля}: код групи, ім’я, сумарна оцінка в балах за 100-бальною системою, оцінка за державною шкалою, набрані на екзамені бали, номер залікової книжки, прізвище, предмет, набрані впродовж семестру бали.

\underline{У допоміжному файлі наявні ключі}: загальна кількість записів, сума оцінок.

\underline{Знайти} всіх студентів, що по кожному предмету отримали сумарний бал не менше 75. Вивести по кожному з них інформацію:\par
{\noindent}--- на першому рядку:\par
номер залікової книжки, прізвище, ім’я, стабільність, кількість «відмінно»;

{\noindent}--- на наступних рядках, починаючи з табуляції, вивести для студента всі його результати (по одному на рядок):\par
державна оцінка прописом, предмет

{\noindent}у такому сортуванні: сумарна оцінка, предмет.

%Студентів сортувати: рейтинг за балами (після округлення), номер залікової книжки.

\bigskip\textit{Обмеження}

Студент не може складати предмет більше одного разу.

Бали та оцінка є значеннями цілих числових типів. Решта полів є непорожніми рядками.
Рядки розпочинатись та закінчуватись білими символами не можуть.

Оцінка за державною шкалою може приймати значення: 2-5~--- як зазвичай, 0~--- не з’явився,$-1$~--- не допущений. 
Оцінки 3-5 вважаються задовільними, решта~--- незадовільними. 
Бали є цілими та невід’ємними.
Набрані впродовж семестру бали не можуть перевищувати 60. 
Набрані на екзамені бали не можуть перевищувати 40. 
Набрані на екзамені бали повинні бути не менше 24 або дорівнювати 0(в останньому випадку оцінка не може бути задовільною). 
Сумарна оцінка в балах є сумою балів, набраних у семестрі та на екзамені.
Сумарна оцінка в балах та за державною шкалою мають бути узгоджені між собою.

Прізвище, ім’я та по-батькові (якщо наявне у варіанті) можуть мати до 30, 30, 25 відповідно символів включно кожне. 
Крім літер алфавіту можуть містити апостроф, дефіс та пробіл.

Предмет може мати від 5 до 53 символів включно. 
Крім літер алфавіту може містити подвійні лапки, пробіл та дефіс.

Код групи є непорожнім та складається не більш ніж з 7 літер. 
Крім літер алфавіту може містити десяткові цифри та дефіс.


Номер залікової книжки однозначно ідентифікує студента та складається з 6 десяткових цифр. 

\bigskip\textit{Для деяких варіантів може знадобитись}

Середній бал/оцінка за предметом визначається як середнє арифметичнесумарних балів/державних оцінок за цим предметом. 
Середній бал/оцінка студента визначається як середнє арифметичне сумарних балів/державних оцінокцього студента. 
За обчислення середньої оцінки недопущені та ті, що не з’явилися, рахуються як 2.
У решті випадків недопуски підсумовуються як $-1$, а «не з’явився» як 0.
Якщо студент не гірше ніж задовільно склав усі предмети, 
то його рейтинг за балами/оцінкою визначається як середній бал/оцінка; інакше дорівнює 0.

Стабільність студента визначається як модуль різниці між його кращим та найгіршим сумарним балом.
Стабільність предмета визначається як модуль різниці між найкращим та найгіршимсумарним балом за цим предметом.






\par
%end
\newpage
%begin
{\center Варіант  \No \textbf { 94 }\par}
\bigskip
%type = absense
%variant = 38
Обробляються дані семестрового журналу перевірок відвідування занять студентами деканатом (фіксуються відсутні). 

\underline{Записи основного файлу містять поля}: аудиторія, день тижня, пара, курс, вид занять, навчальний тиждень, предмет, прізвище, ім’я.

\underline{У допоміжному файлі наявні ключі}: кількість різних предметів, довжина найкоротшого прізвища.

\underline{Знайти} всі предмети, що пропускалися найбільше. Вивести по кожному з них інформацію:\par
{\noindent}--- на першому рядку:\par
курс, предмет, кількість пропусків, кількість пропусків лекцій;

{\noindent}--- на наступних рядках, починаючи з табуляції, вивести пропуски предмета (по одному на рядок):\par
прізвище, ім’я, тиждень, день, пара, вид занять

{\noindent}у такому сортуванні: вид занять, прізвище, ім’я, тиждень, день, пара.

%Предмети сортувати: кількість пропусків(за спаданням), курс, предмет.

\bigskip\textit{Обмеження}

Навчальний тиждень може приймати цілі значення від 1 до 14.
День тижня є цілим від 1 до 5. Пара є цілим від 1 до 4.
Курс є цілим від 1 до 4. Аудиторія є додатнім цілим.
Вид занять може приймати значення:
ЛЕ --- лекція, ПР --- практичне, Sem --- семінар, lab --- лабораторне.

Решта полів є непорожніми рядками.
Рядки починатися та закінчуватися білими символами не можуть.

Групи кодуються в межах курсу.
Код групи складається не більш ніж з 6 символів.
Може містити літери алфавіту, десяткові цифри та дефіс.

Предмет є рядком, може мати від 5 до 25 символів включно.
Може містити літери алфавіту, десяткові цифри, апостроф, дефіс, пробіл.
Предмети різних курсів вважати різними навіть за однакових назв.

Прізвище, ім’я та по-батькові (за наявності у варіанті)
можуть мати від 5 до 29, 29, 22 відповідно символів включно кожне.
Можуть містити літери алфавіту, апостроф, пробіл та дефіс.
Студенти однозначно ідентифікуються прізвищем, ім’ям та по-батькові (за наявності у варіанті).
Жоден студент не вчиться одночасно на різних курсах чи в різних групах.

Види занять сортуються так: семінар < лабораторне < лекція < практичне.

%\bigskip%\textit{Для деяких варіантів може знадобитись}




\par
%end
\newpage
%begin
{\center Варіант  \No \textbf { 95 }\par}
\bigskip
%type = meteo
%variant = 107
Обробляються дані спостережень за погодою. 

\underline{Записи основного файлу містять поля}: кількість опадів, відносна вологість, рік, код метеостанції, місяць, день, середня денна температура, максимальна денна температура, мінімальна денна температура, сила вітру.

\underline{У допоміжному файлі наявні ключі}: найбільша кількість опадів у спостереженнях, найменша кількість опадів у спостереженнях, загальна кількість записів.

\underline{Знайти} 4 метеостанції, на яких була зафіксована найбільша середньоденна температура. Вивести по кожному з них інформацію:\par
{\noindent}--- на першому рядку:\par
середня вологість, середня середньоденна температура, найбільша середньоденна температура, найбільша сила вітру, метеостанція;

{\noindent}--- на наступних рядках, починаючи з табуляції, вивести для них дані по датах (по одному на рядок):\par
день, місяць, день, вологість, середня денна температура, максимальна денна температура, сила вітру, вологість, рік

{\noindent}у такому сортуванні: вологість (за спаданням), місяць, день, рік.

%Метеостанція сортувати: кількість спостережень (за спаданням), найбільша сила вітру, метеостанція.

\bigskip\textit{Обмеження}

Код метеостанції є рядком довжини від 4 до 7 включно.
Може містити літери алфавіту, десяткові цифри, підкреслення. Решта полів числові.

Рік є чотиризначним цілим.
Кількість опадів є дійсним значенням, подається у форматі з 2 знаками після крапки.
Сила вітру є цілим значенням.
Температури є дійсними значеннями, подаються у форматі з 2 знаками після крапки.
Цілі значення подаються без десяткової крапки.
Для дійсних використовується зображення з фіксованою крапкою.

\textit{Вихідне форматування дійсних}

Усі дійсні значення виводяться у форматі з фіксованою крапкою з тією ж кількістю знаків після крапки, що і у вхідному зображенні.

Якщо у варіанті раніше не було зазначено, то \underline{обчислювані програмою} середні значення виводяться у форматі з 2 знаками після крапки.
%\bigskip%\textit{Для деяких варіантів може знадобитись}




\par
%end
\newpage
%begin
{\center Варіант  \No \textbf { 96 }\par}
\bigskip
%type = booking
%variant = 127
Обробляється інформація щодо відвантажених накладних.

\underline{Записи основного файлу містять поля}: кількість, номер накладної, вартість, назва товару, ціна, номер товарної позиції у межах накладної.

\underline{У допоміжному файлі наявні ключі}: середня кількість товарних позицій у накладній, найбільша кількість товарних позицій у накладній.

\underline{Знайти} всі накладні, в яких було відвантажено рівно один товар. Вивести по кожному з них інформацію:\par
{\noindent}--- на першому рядку:\par
номер накладної, кількість товару, кількість товарних позицій, загальна сума;

{\noindent}--- на наступних рядках, починаючи з табуляції, вивести для накладної всі її товарні позиції (по одному на рядок):\par
номер товарної позиції, назва товару, ціна, кількість, вартість

{\noindent}у такому сортуванні: назва товару, ціна.

%Накладні сортувати: кількість товару (за спаданням), номер накладної.

\bigskip\textit{Обмеження}

Кількість, ціна та вартість є значеннями дійсного типу.
Решта полів є непорожніми рядками.
Рядки починатися та закінчуватися білими символами не можуть.

Номер накладної складається з 12 символів (цифри, латинські літери).
Назва товару може мати від 3 до 24 символів включно.
Може містити тільки цифри, літери, пробіли, апостроф, дефіс, підкреслення.

Кількість, ціна мають бути додатними.
Кількість вказується з рівно 4 знаками після крапки.
Ціна вказується з рівно 2 знаками після крапки.
Вартість має бути узгоджена з кількістю та ціною та записується рівно з двома десятковими знаками.
(Застосовується округлення.)

Товари однозначно ідентифікуються назвами та одиницями виміру (за наявності).

Одиниця виміру (за наявності у варіанті) може приймати значення: шт., кг, м , кор.

Кожна накладна має свій власний унікальний номер.
У накладній можуть відвантажуватися кілька товарів (не менше 1).
Товарна позиція --- рядок накладної: товар, кількість, ціна, вартість та, за наявності, одиниця виміру.
Товарні позиції кожної накладної нумеруються послідовними натуральними числами, починаючи з 1.
В одній накладній товари можуть повторюватися, 
тобто товарні позиції однієї накладної можуть містити однакові товари 
(при цьому ціна може бути і різною, і однаковою).

\bigskip\textit{Для деяких варіантів може знадобитись}

Недоотриману вигоду розраховувати як суму виразів
{\center (найбільша ціна відвантаження товару – ціна відвантаження товару)*кількість\par}

\bigskip
{\noindent} округлених до 2 знаків після крапки.




\par
%end
\newpage
%begin
{\center Варіант  \No \textbf { 97 }\par}
\bigskip
%type = booking
%variant = 120
Обробляється інформація щодо відвантажених накладних.

\underline{Записи основного файлу містять поля}: ціна, назва товару, вартість, кількість, номер товарної позиції у межах накладної, номер накладної.

\underline{У допоміжному файлі наявні ключі}: сума вартостей, кількість різних товарів.

\underline{Знайти} три найбільш популярні товари (ті, що сумарно відвантажувалися в найбільшій кількості). Вивести по кожному з них інформацію:\par
{\noindent}--- на першому рядку:\par
назва товару, сума відвантаження (з 2 знаками після крапки), найбільша відвантажена кількість в одній накладній (з 1 знаком після крапки);

{\noindent}--- на наступних рядках, починаючи з табуляції, вивести для товару всі накладні, де він відвантажувався більш ніж на 1000 (по одному на рядок):\par
вартість, ціна, кількість, номер накладної

{\noindent}у такому сортуванні: кількість, вартість, ціна, номер накладної.

%Товари сортувати: назва товару (за спаданням).

\bigskip\textit{Обмеження}

Кількість, ціна та вартість є значеннями дійсного типу.
Решта полів є непорожніми рядками.
Рядки починатися та закінчуватися білими символами не можуть.

Номер накладної складається з 13 символів (цифри, латинські літери).
Назва товару може мати від 4 до 28 символів включно.
Може містити тільки цифри, літери, пробіли, апостроф, дефіс, підкреслення.

Кількість, ціна мають бути додатними.
Кількість вказується з рівно 2 знаками після крапки.
Ціна вказується з рівно 3 знаками після крапки.
Вартість має бути узгоджена з кількістю та ціною та записується рівно з двома десятковими знаками.
(Застосовується округлення.)

Товари однозначно ідентифікуються назвами та одиницями виміру (за наявності).

Одиниця виміру (за наявності у варіанті) може приймати значення: шт., кг, м , кор.

Кожна накладна має свій власний унікальний номер.
У накладній можуть відвантажуватися кілька товарів (не менше 1).
Товарна позиція --- рядок накладної: товар, кількість, ціна, вартість та, за наявності, одиниця виміру.
Товарні позиції кожної накладної нумеруються послідовними натуральними числами, починаючи з 1.
В одній накладній товари можуть повторюватися, 
тобто товарні позиції однієї накладної можуть містити однакові товари 
(при цьому ціна може бути і різною, і однаковою).

\bigskip\textit{Для деяких варіантів може знадобитись}

Недоотриману вигоду розраховувати як суму виразів
{\center (найбільша ціна відвантаження товару – ціна відвантаження товару)*кількість\par}

\bigskip
{\noindent} округлених до 2 знаків після крапки.




\par
%end
\newpage
%begin
{\center Варіант  \No \textbf { 98 }\par}
\bigskip
%type = booking
%variant = 131
Обробляється інформація щодо відвантажених накладних.

\underline{Записи основного файлу містять поля}: кількість, номер накладної, назва товару, ціна, номер товарної позиції у межах накладної, вартість.

\underline{У допоміжному файлі наявні ключі}: сума цін, сума кількостей.

\underline{Знайти} 5 накладних, в яких було відвантажено товарів на найбільшу суму. Вивести по кожному з них інформацію:\par
{\noindent}--- на першому рядку:\par
номер накладної, кількість товару, кількість товарних позицій, загальна сума, середня вартість товарів у накладній (з 2 знаками після крапки);

{\noindent}--- на наступних рядках, починаючи з табуляції, вивести для накладної всі її товарні позиції, де відвантажувалося менше 1 одиниці товару (по одному на рядок):\par
назва товару, ціна, кількість, вартість, номер товарної позиції

{\noindent}у такому сортуванні: назва товару, кількість, ціна.

%Накладні сортувати: середня вартість товарів у накладній (за спаданням), номер накладної.

\bigskip\textit{Обмеження}

Кількість, ціна та вартість є значеннями дійсного типу.
Решта полів є непорожніми рядками.
Рядки починатися та закінчуватися білими символами не можуть.

Номер накладної складається з 11 символів (цифри, латинські літери).
Назва товару може мати від 1 до 25 символів включно.
Може містити тільки цифри, літери, пробіли, апостроф, дефіс, підкреслення.

Кількість, ціна мають бути додатними.
Кількість вказується з рівно 1 знаками після крапки.
Ціна вказується з рівно 4 знаками після крапки.
Вартість має бути узгоджена з кількістю та ціною та записується рівно з двома десятковими знаками.
(Застосовується округлення.)

Товари однозначно ідентифікуються назвами та одиницями виміру (за наявності).

Одиниця виміру (за наявності у варіанті) може приймати значення: шт., кг, м , кор.

Кожна накладна має свій власний унікальний номер.
У накладній можуть відвантажуватися кілька товарів (не менше 1).
Товарна позиція --- рядок накладної: товар, кількість, ціна, вартість та, за наявності, одиниця виміру.
Товарні позиції кожної накладної нумеруються послідовними натуральними числами, починаючи з 1.
В одній накладній товари можуть повторюватися, 
тобто товарні позиції однієї накладної можуть містити однакові товари 
(при цьому ціна може бути і різною, і однаковою).

\bigskip\textit{Для деяких варіантів може знадобитись}

Недоотриману вигоду розраховувати як суму виразів
{\center (найбільша ціна відвантаження товару – ціна відвантаження товару)*кількість\par}

\bigskip
{\noindent} округлених до 2 знаків після крапки.




\par
%end
\newpage
%begin
{\center Варіант  \No \textbf { 99 }\par}
\bigskip
%type = absense
%variant = 40
Обробляються дані семестрового журналу перевірок відвідування занять студентами деканатом (фіксуються відсутні). 

\underline{Записи основного файлу містять поля}: аудиторія, вид занять, курс, ім’я, предмет, день тижня, пара, прізвище, навчальний тиждень.

\underline{У допоміжному файлі наявні ключі}: довжина найдовшої назви предмету, загальна кількість записів.

\underline{Знайти} всі предмети, які пропускалися, але пропусків саме лабораторних занять зафіксовано не було. Вивести по кожному з них інформацію:\par
{\noindent}--- на першому рядку:\par
курс, предмет, кількість пропусків;

{\noindent}--- на наступних рядках, починаючи з табуляції, вивести пропуски предмета (по одному на рядок):\par
тиждень, день, вид занять, кількість пропусків

{\noindent}у такому сортуванні: тиждень, вид занять, день.

%Премети сортувати: курс, кількість пропусків (за спаданням), предмет.

\bigskip\textit{Обмеження}

Навчальний тиждень може приймати цілі значення від 1 до 20.
День тижня є цілим від 1 до 5. Пара є цілим від 1 до 4.
Курс є цілим від 1 до 4. Аудиторія є додатнім цілим.
Вид занять може приймати значення:
Lecture --- лекція, практ. --- практичне, S --- семінар, лаб. --- лабораторне.

Решта полів є непорожніми рядками.
Рядки починатися та закінчуватися білими символами не можуть.

Групи кодуються в межах курсу.
Код групи складається не більш ніж з 6 символів.
Може містити літери алфавіту, десяткові цифри та дефіс.

Предмет є рядком, може мати від 4 до 36 символів включно.
Може містити літери алфавіту, десяткові цифри, апостроф, дефіс, пробіл.
Предмети різних курсів вважати різними навіть за однакових назв.

Прізвище, ім’я та по-батькові (за наявності у варіанті)
можуть мати від 4 до 21, 24, 25 відповідно символів включно кожне.
Можуть містити літери алфавіту, апостроф, пробіл та дефіс.
Студенти однозначно ідентифікуються прізвищем, ім’ям та по-батькові (за наявності у варіанті).
Жоден студент не вчиться одночасно на різних курсах чи в різних групах.

Види занять сортуються так: лекція < лабораторне < практичне < семінар.

%\bigskip%\textit{Для деяких варіантів може знадобитись}




\par
%end
\newpage
%begin
{\center Варіант  \No \textbf { 100 }\par}
\bigskip
%type = booking
%variant = 129
Обробляється інформація щодо відвантажених накладних.

\underline{Записи основного файлу містять поля}: вартість, одиниця виміру, ціна, назва товару, номер товарної позиції у межах накладної, кількість, номер накладної.

\underline{У допоміжному файлі наявні ключі}: сума кількостей, середня вартість.

\underline{Знайти} всі накладні, в яких було відвантажено більше 10 різних товарів. Вивести по кожному з них інформацію:\par
{\noindent}--- на першому рядку:\par
номер накладної, кількість різних товарів, кількість товарних позицій, середня ціна товарної позиції (з 3 знаками після крапки);

{\noindent}--- на наступних рядках, починаючи з табуляції, вивести для накладної всі товарні позиції, що відвантажувалися в ній на найменшу вартість (по одному на рядок):\par
номер товарної позиції, ціна, кількість, вартість

{\noindent}у такому сортуванні: номер товарної позиції.

%Накладні сортувати: номер накладної (за спаданням).

\bigskip\textit{Обмеження}

Кількість, ціна та вартість є значеннями дійсного типу.
Решта полів є непорожніми рядками.
Рядки починатися та закінчуватися білими символами не можуть.

Номер накладної складається з 10 символів (цифри, латинські літери).
Назва товару може мати від 2 до 26 символів включно.
Може містити тільки цифри, літери, пробіли, апостроф, дефіс, підкреслення.

Кількість, ціна мають бути додатними.
Кількість вказується з рівно 3 знаками після крапки.
Ціна вказується з рівно 3 знаками після крапки.
Вартість має бути узгоджена з кількістю та ціною та записується рівно з двома десятковими знаками.
(Застосовується округлення.)

Товари однозначно ідентифікуються назвами та одиницями виміру (за наявності).

Одиниця виміру (за наявності у варіанті) може приймати значення: шт., кг, м , кор.

Кожна накладна має свій власний унікальний номер.
У накладній можуть відвантажуватися кілька товарів (не менше 1).
Товарна позиція --- рядок накладної: товар, кількість, ціна, вартість та, за наявності, одиниця виміру.
Товарні позиції кожної накладної нумеруються послідовними натуральними числами, починаючи з 1.
В одній накладній товари можуть повторюватися, 
тобто товарні позиції однієї накладної можуть містити однакові товари 
(при цьому ціна може бути і різною, і однаковою).

\bigskip\textit{Для деяких варіантів може знадобитись}

Недоотриману вигоду розраховувати як суму виразів
{\center (найбільша ціна відвантаження товару – ціна відвантаження товару)*кількість\par}

\bigskip
{\noindent} округлених до 2 знаків після крапки.




\par
%end
\newpage
%begin
{\center Варіант  \No \textbf { 101 }\par}
\bigskip
%type = session
%variant = 81
Обробляються результати складання студентами екзаменів зимової сесії.

\underline{Записи основного файлу містять поля}: номер залікової книжки, ім’я, предмет, код групи, по-батько\-ві, сумарна оцінка в балах за 100-бальною системою, оцінка за державною шкалою, прізвище, набрані на екзамені бали.

\underline{У допоміжному файлі наявні ключі}: LISTING, загальна кількість записів, кількість оцінок вище 75 балів.

\underline{Знайти} предмети, які були успішно здані найбільшою кількістю студентів. Вивести по кожному з них інформацію:\par
{\noindent}--- на першому рядку:\par
предмет, середній бал за предметом (округлений до цілого), кількість студентів, що складали, кількість студентів, що склали, кількість «відмінно»;

{\noindent}--- на наступних рядках, починаючи з табуляції, вивести для предмета студентів, що його не склали (по одному на рядок):\par
група, ім’я, прізвище, номер залікової книжки, оцінка в балах, оцінка за державною шкалою

{\noindent}у такому сортуванні: група, прізвище, ім’я, номер залікової книжки.

%Предмети сортувати: кількість «відмінно» (за спаданням), кількість студентів, що складали, предмет.

\bigskip\textit{Обмеження}

Студент не може складати предмет більше одного разу.

Бали та оцінка є значеннями цілих числових типів. Решта полів є непорожніми рядками.
Рядки розпочинатись та закінчуватись білими символами не можуть.

Оцінка за державною шкалою може приймати значення: 2-5~--- як зазвичай, 0~--- не з’явився,$-1$~--- не допущений. 
Оцінки 3-5 вважаються задовільними, решта~--- незадовільними. 
Бали є цілими та невід’ємними.
Набрані впродовж семестру бали не можуть перевищувати 60. 
Набрані на екзамені бали не можуть перевищувати 40. 
Набрані на екзамені бали повинні бути не менше 24 або дорівнювати 0(в останньому випадку оцінка не може бути задовільною). 
Сумарна оцінка в балах є сумою балів, набраних у семестрі та на екзамені.
Сумарна оцінка в балах та за державною шкалою мають бути узгоджені між собою.

Прізвище, ім’я та по-батькові (якщо наявне у варіанті) можуть мати до 25, 28, 30 відповідно символів включно кожне. 
Крім літер алфавіту можуть містити апостроф, дефіс та пробіл.

Предмет може мати від 3 до 69 символів включно. 
Крім літер алфавіту може містити подвійні лапки, пробіл та дефіс.

Код групи є непорожнім та складається не більш ніж з 3 літер. 
Крім літер алфавіту може містити десяткові цифри та дефіс.


Номер залікової книжки однозначно ідентифікує студента та складається з 7 десяткових цифр. 

\bigskip\textit{Для деяких варіантів може знадобитись}

Середній бал/оцінка за предметом визначається як середнє арифметичнесумарних балів/державних оцінок за цим предметом. 
Середній бал/оцінка студента визначається як середнє арифметичне сумарних балів/державних оцінокцього студента. 
За обчислення середньої оцінки недопущені та ті, що не з’явилися, рахуються як 2.
У решті випадків недопуски підсумовуються як $-1$, а «не з’явився» як 0.
Якщо студент не гірше ніж задовільно склав усі предмети, 
то його рейтинг за балами/оцінкою визначається як середній бал/оцінка; інакше дорівнює 0.

Стабільність студента визначається як модуль різниці між його кращим та найгіршим сумарним балом.
Стабільність предмета визначається як модуль різниці між найкращим та найгіршимсумарним балом за цим предметом.






\par
%end
\newpage
%begin
{\center Варіант  \No \textbf { 102 }\par}
\bigskip
%type = session
%variant = 78
Обробляються результати складання студентами екзаменів зимової сесії.

\underline{Записи основного файлу містять поля}: сумарна оцінка в балах за 100-бальною системою, номер залікової книжки, набрані впродовж семестру бали, код групи, прізвище, предмет, оцінка за державною шкалою, ім’я, набрані на екзамені бали.

\underline{У допоміжному файлі наявні ключі}: загальна кількість предметів, кількість «не з’явився».

\underline{Знайти} предмети, які були складені менше, ніж половиною студентів. Вивести по кожному з них інформацію:\par
{\noindent}--- на першому рядку:\par
кількість, що складали, кількість, що не склали, середня оцінка (округлена до 1 знаку), предмет;

{\noindent}--- на наступних рядках, починаючи з табуляції, вивести для предмета студентів, що його успішно склали (по одному на рядок):\par
ім’я, прізвище, номер залікової книжки, оцінка в балах, оцінка за державною шкалою, бали за екзамен, бали за семестр

{\noindent}у такому сортуванні: прізвище, ім’я, номер залікової книжки.

%Предмети сортувати: середня оцінка (округлена), предмет.

\bigskip\textit{Обмеження}

Студент не може складати предмет більше одного разу.

Бали та оцінка є значеннями цілих числових типів. Решта полів є непорожніми рядками.
Рядки розпочинатись та закінчуватись білими символами не можуть.

Оцінка за державною шкалою може приймати значення: 2-5~--- як зазвичай, 0~--- не з’явився,$-1$~--- не допущений. 
Оцінки 3-5 вважаються задовільними, решта~--- незадовільними. 
Бали є цілими та невід’ємними.
Набрані впродовж семестру бали не можуть перевищувати 60. 
Набрані на екзамені бали не можуть перевищувати 40. 
Набрані на екзамені бали повинні бути не менше 24 або дорівнювати 0(в останньому випадку оцінка не може бути задовільною). 
Сумарна оцінка в балах є сумою балів, набраних у семестрі та на екзамені.
Сумарна оцінка в балах та за державною шкалою мають бути узгоджені між собою.

Прізвище, ім’я та по-батькові (якщо наявне у варіанті) можуть мати до 21, 24, 30 відповідно символів включно кожне. 
Крім літер алфавіту можуть містити апостроф, дефіс та пробіл.

Предмет може мати від 4 до 70 символів включно. 
Крім літер алфавіту може містити подвійні лапки, пробіл та дефіс.

Код групи є непорожнім та складається не більш ніж з 4 літер. 
Крім літер алфавіту може містити десяткові цифри та дефіс.


Номер залікової книжки однозначно ідентифікує студента та складається з 6 десяткових цифр. 

\bigskip\textit{Для деяких варіантів може знадобитись}

Середній бал/оцінка за предметом визначається як середнє арифметичнесумарних балів/державних оцінок за цим предметом. 
Середній бал/оцінка студента визначається як середнє арифметичне сумарних балів/державних оцінокцього студента. 
За обчислення середньої оцінки недопущені та ті, що не з’явилися, рахуються як 2.
У решті випадків недопуски підсумовуються як $-1$, а «не з’явився» як 0.
Якщо студент не гірше ніж задовільно склав усі предмети, 
то його рейтинг за балами/оцінкою визначається як середній бал/оцінка; інакше дорівнює 0.

Стабільність студента визначається як модуль різниці між його кращим та найгіршим сумарним балом.
Стабільність предмета визначається як модуль різниці між найкращим та найгіршимсумарним балом за цим предметом.






\par
%end
\newpage
%begin
{\center Варіант  \No \textbf { 103 }\par}
\bigskip
%type = meteo
%variant = 87
Обробляються дані спостережень за погодою. 

\underline{Записи основного файлу містять поля}: кількість опадів, рік, відносна вологість, максимальна денна температура, сила вітру, код метеостанції, день, місяць, мінімальна денна температура, середня денна температура.

\underline{У допоміжному файлі наявні ключі}: сумарна кількість опадів, найменша мінімальна температура.

\underline{Знайти} 5 дат, по яких фіксувалася найменша середня денна температура. Вивести по кожному з них інформацію:\par
{\noindent}--- на першому рядку:\par
рік, місяць, день, кількість опадів (сумарна), найменша середня денна температура;

{\noindent}--- на наступних рядках, починаючи з табуляції, вивести для дати спостереження метеостанції, де середня денна температура була не більш ніж 10 більша за найменшу середню денну температуру за відповідну дату (по одному на рядок):\par
метеостанція, сила вітру, вологість

{\noindent}у такому сортуванні: метеостанція (за спаданням).

%Дати сортувати: день, місяць, кількість опадів, день.

\bigskip\textit{Обмеження}

Код метеостанції є рядком довжини від 3 до 9 включно.
Може містити літери алфавіту, десяткові цифри, підкреслення. Решта полів числові.

Рік є чотиризначним цілим.
Кількість опадів є дійсним значенням, подається у форматі з 1 знаком після крапки.
Сила вітру є дійсним значенням, подається у форматі з 1 знаком після крапки.
Температури є дійсними значеннями, подаються у форматі з 2 знаками після крапки.
Цілі значення подаються без десяткової крапки.
Для дійсних використовується зображення з фіксованою крапкою.

\textit{Вихідне форматування дійсних}

Усі дійсні значення виводяться у форматі з фіксованою крапкою з тією ж кількістю знаків після крапки, що і у вхідному зображенні.

Якщо у варіанті раніше не було зазначено, то \underline{обчислювані програмою} середні значення виводяться у форматі з 2 знаками після крапки.
%\bigskip%\textit{Для деяких варіантів може знадобитись}




\par
%end
\newpage
%begin
{\center Варіант  \No \textbf { 104 }\par}
\bigskip
%type = session
%variant = 63
Обробляються результати складання студентами екзаменів зимової сесії.

\underline{Записи основного файлу містять поля}: ім’я, номер залікової книжки, сумарна оцінка в балах за 100-бальною системою, оцінка за державною шкалою, по-батькові, набрані на екзамені бали, предмет, прізвище, код групи.

\underline{У допоміжному файлі наявні ключі}: кількість студентів, у яких були недопуски, загальна кількість записів.

\underline{Знайти} всіх студентів, що по кожному предмету отримали не менше 4 за державною шкалою. Вивести по кожному з них інформацію:\par
{\noindent}--- на першому рядку:\par
середня оцінка (округлена до 2 знаків), кількість екзаменів, номер залікової книжки;

{\noindent}--- на наступних рядках, починаючи з табуляції, вивести для студента всі його результати (по одному на рядок):\par
державна оцінка, предмет, сумарний бал

{\noindent}у такому сортуванні: бали на екзамені, державна оцінка, предмет.

%Студентів сортувати: кількість екзаменів (за спаданням), середня оцінка (після округлення), прізвище, ім’я, по-батькові, номер залікової книжки.

\bigskip\textit{Обмеження}

Студент не може складати предмет більше одного разу.

Бали та оцінка є значеннями цілих числових типів. Решта полів є непорожніми рядками.
Рядки розпочинатись та закінчуватись білими символами не можуть.

Оцінка за державною шкалою може приймати значення: 2-5~--- як зазвичай, 0~--- не з’явився,$-1$~--- не допущений. 
Оцінки 3-5 вважаються задовільними, решта~--- незадовільними. 
Бали є цілими та невід’ємними.
Набрані впродовж семестру бали не можуть перевищувати 60. 
Набрані на екзамені бали не можуть перевищувати 40. 
Набрані на екзамені бали повинні бути не менше 24 або дорівнювати 0(в останньому випадку оцінка не може бути задовільною). 
Сумарна оцінка в балах є сумою балів, набраних у семестрі та на екзамені.
Сумарна оцінка в балах та за державною шкалою мають бути узгоджені між собою.

Прізвище, ім’я та по-батькові (якщо наявне у варіанті) можуть мати до 27, 28, 27 відповідно символів включно кожне. 
Крім літер алфавіту можуть містити апостроф, дефіс та пробіл.

Предмет може мати від 5 до 63 символів включно. 
Крім літер алфавіту може містити подвійні лапки, пробіл та дефіс.

Код групи є непорожнім та складається не більш ніж з 5 літер. 
Крім літер алфавіту може містити десяткові цифри та дефіс.


Номер залікової книжки однозначно ідентифікує студента та складається з 8 десяткових цифр. 

\bigskip\textit{Для деяких варіантів може знадобитись}

Середній бал/оцінка за предметом визначається як середнє арифметичнесумарних балів/державних оцінок за цим предметом. 
Середній бал/оцінка студента визначається як середнє арифметичне сумарних балів/державних оцінокцього студента. 
За обчислення середньої оцінки недопущені та ті, що не з’явилися, рахуються як 2.
У решті випадків недопуски підсумовуються як $-1$, а «не з’явився» як 0.
Якщо студент не гірше ніж задовільно склав усі предмети, 
то його рейтинг за балами/оцінкою визначається як середній бал/оцінка; інакше дорівнює 0.

Стабільність студента визначається як модуль різниці між його кращим та найгіршим сумарним балом.
Стабільність предмета визначається як модуль різниці між найкращим та найгіршимсумарним балом за цим предметом.






\par
%end
\newpage
%begin
{\center Варіант  \No \textbf { 105 }\par}
\bigskip
%type = tasks
%variant = 8
Обробляються результати розв’язування студентами задач на контрольних.
Одиничний факт роз\-в’я\-зу\-вання студентом задачі надалі розуміється як спроба.

\underline{Записи основного файлу містять поля}: ім’я, відсоток набраних за розв’язання задачі балів, по-бать\-ко\-ві, умова задачі, рік, прізвище.

\underline{У допоміжному файлі наявні ключі}: кількість різних задач, кількість спроб з 50\% балів.

\underline{Знайти} всі задачі, для яких існує спроба з відсотком розв’язання менш ніж 60\%. Вивести по кожному з них інформацію:\par
{\noindent}--- на першому рядку:\par
умова, кількість спроб, кількість спроб гірше 60\%, середня розв’язуваність (округлена до 1 знаку);

{\noindent}--- на наступних рядках, починаючи з табуляції, вивести спроби по цих задачах з відсотком не більше 75\% (по одному на рядок):\par
ім’я, по-батькові, прізвище, відсоток набраних балів, рік

{\noindent}у такому сортуванні: рік (за спаданням), ім’я, по-батькові, прізвище, відсоток набраних балів (за спаданням).

%Задачі сортувати: кількість спроб гірше 60\%, умова.

\bigskip\textit{Обмеження}

Студент міг розв’язувати задачу кілька разів.

Відсоток набраних за розв’язання задачі балів є цілим від 0 до 100.
День, місяць, рік є числовими полями, що зображують дату.
Раніше 2001 року статистика не велася. Рік є 4-значним.
Решта полів є непорожніми рядками.
Рядки починатися та закінчуватися білими символами не можуть.

Прізвище, ім’я та по-батькові (за наявності у варіанті) можуть мати до 26, 20, 27 відповідно символів включно кожне. 
Крім літер алфавіту можуть містити апостроф, дефіс та пробіл.

Умова задачі може мати від 5 до 60 символів включно.
Крім літер алфавіту може містити подвійні лапки, десяткові цифри, пробіл, дефіс,
а також символи . , \_ \{ \} \$ $+$ $*$ {\textbackslash} ( ) \% $[$ $]$ .

Студенти однозначно ідентифікуються прізвищем, ім’ям та по-батькові (за наявності у варіанті).
Задачі однозначно ідентифікуються умовою.

\bigskip\textit{Для деяких варіантів може знадобитись}

Розв'язуваність задачі --- середній відсоток набраних за задачу балів.

Розв'язуваність студента --- середній відсоток набраних студентом балів по всіх його спробах.

Розв'язуваність --- середній відсоток набраних студентами за задачі балів (рахується по всім даним).




\par
%end
\newpage
%begin
{\center Варіант  \No \textbf { 106 }\par}
\bigskip
%type = booking
%variant = 126
Обробляється інформація щодо відвантажених накладних.

\underline{Записи основного файлу містять поля}: номер товарної позиції у межах накладної, назва товару, вартість, одиниця виміру, ціна, кількість, номер накладної.

\underline{У допоміжному файлі наявні ключі}: найменша ціна товарної позиції, кількість записів.

\underline{Знайти} 5 накладних, в яких було відвантажено найбільшу кількість різних товарів. Вивести по кожному з них інформацію:\par
{\noindent}--- на першому рядку:\par
номер накладної, кількість товару, кількість товарних позицій, кількість різних товарів, загальна сума, середня ціна товарної позиції (з 2 знаками після крапки);

{\noindent}--- на наступних рядках, починаючи з табуляції, вивести для накладної всі товарні позиції, що відвантажувалися в ній в найбільшій кількості (по одному на рядок):\par
номер товарної позиції, ціна, кількість, вартість, назва товару, одиниця виміру

{\noindent}у такому сортуванні: номер товарної позиції.

%Накладні сортувати: номер накладної (за спаданням).

\bigskip\textit{Обмеження}

Кількість, ціна та вартість є значеннями дійсного типу.
Решта полів є непорожніми рядками.
Рядки починатися та закінчуватися білими символами не можуть.

Номер накладної складається з 15 символів (цифри, латинські літери).
Назва товару може мати від 2 до 20 символів включно.
Може містити тільки цифри, літери, пробіли, апостроф, дефіс, підкреслення.

Кількість, ціна мають бути додатними.
Кількість вказується з рівно 4 знаками після крапки.
Ціна вказується з рівно 2 знаками після крапки.
Вартість має бути узгоджена з кількістю та ціною та записується рівно з двома десятковими знаками.
(Застосовується округлення.)

Товари однозначно ідентифікуються назвами та одиницями виміру (за наявності).

Одиниця виміру (за наявності у варіанті) може приймати значення: шт., кг, м , кор.

Кожна накладна має свій власний унікальний номер.
У накладній можуть відвантажуватися кілька товарів (не менше 1).
Товарна позиція --- рядок накладної: товар, кількість, ціна, вартість та, за наявності, одиниця виміру.
Товарні позиції кожної накладної нумеруються послідовними натуральними числами, починаючи з 1.
В одній накладній товари можуть повторюватися, 
тобто товарні позиції однієї накладної можуть містити однакові товари 
(при цьому ціна може бути і різною, і однаковою).

\bigskip\textit{Для деяких варіантів може знадобитись}

Недоотриману вигоду розраховувати як суму виразів
{\center (найбільша ціна відвантаження товару – ціна відвантаження товару)*кількість\par}

\bigskip
{\noindent} округлених до 2 знаків після крапки.




\par
%end
\newpage
%begin
{\center Варіант  \No \textbf { 107 }\par}
\bigskip
%type = absense
%variant = 27
Обробляються дані семестрового журналу перевірок відвідування занять студентами деканатом (фіксуються відсутні). 

\underline{Записи основного файлу містять поля}: предмет, курс, прізвище, ім’я, вид занять, навчальний тиждень, день тижня, пара, аудиторія.

\underline{У допоміжному файлі наявні ключі}: кількість пропусків лекцій, загальна кількість записів.

\underline{Знайти} всіх студентів, що пропускали лабораторні заняття. Вивести по кожному з них інформацію:\par
{\noindent}--- на першому рядку:\par
кількість пропущених пар, кількість пропущених лабораторних, відсоток пропусків лабораторних серед всіх пропусків студента (округлений до цілого), прізвище, ім’я,;

{\noindent}--- на наступних рядках, починаючи з табуляції, вивести пропуски лабораторних занять студентом (по одному на рядок):\par
предмет, тиждень, день, пара, аудиторія

{\noindent}у такому сортуванні: предмет, тиждень, день, пара.

%Студентів сортувати: прізвище, ім’я.

\bigskip\textit{Обмеження}

Навчальний тиждень може приймати цілі значення від 1 до 16.
День тижня є цілим від 1 до 5. Пара є цілим від 1 до 4.
Курс є цілим від 1 до 4. Аудиторія є додатнім цілим.
Вид занять може приймати значення:
Le --- лекція, pr --- практичне, Seminar --- семінар, Lab --- лабораторне.

Решта полів є непорожніми рядками.
Рядки починатися та закінчуватися білими символами не можуть.

Групи кодуються в межах курсу.
Код групи складається не більш ніж з 4 символів.
Може містити літери алфавіту, десяткові цифри та дефіс.

Предмет є рядком, може мати від 3 до 32 символів включно.
Може містити літери алфавіту, десяткові цифри, апостроф, дефіс, пробіл.
Предмети різних курсів вважати різними навіть за однакових назв.

Прізвище, ім’я та по-батькові (за наявності у варіанті)
можуть мати від 3 до 22, 26, 28 відповідно символів включно кожне.
Можуть містити літери алфавіту, апостроф, пробіл та дефіс.
Студенти однозначно ідентифікуються прізвищем, ім’ям та по-батькові (за наявності у варіанті).
Жоден студент не вчиться одночасно на різних курсах чи в різних групах.

Види занять сортуються так: лекція < практичне < семінар < лабораторне.

%\bigskip%\textit{Для деяких варіантів може знадобитись}




\par
%end
\newpage
%begin
{\center Варіант  \No \textbf { 108 }\par}
\bigskip
%type = absense
%variant = 23
Обробляються дані семестрового журналу перевірок відвідування занять студентами деканатом (фіксуються відсутні). 

\underline{Записи основного файлу містять поля}: вид занять, аудиторія, навчальний тиждень, день тижня, пара, ім’я, предмет, прізвище, курс.

\underline{У допоміжному файлі наявні ключі}: загальна кількість записів, загальна кількість пропусків семінарів.

\underline{Знайти} всіх студентів, що пропустили більше 10 пар. Вивести по кожному з них інформацію:\par
{\noindent}--- на першому рядку:\par
прізвище, ім’я, кількість пропущених пар, кількість пропущених практичних;

{\noindent}--- на наступних рядках, починаючи з табуляції, вивести пропуски студента (по одному на рядок):\par
предмет, вид занять, тиждень, день, пара, аудиторія

{\noindent}у такому сортуванні: предмет, вид занять, тиждень, день, пара, аудиторія.

%Студентів сортувати: кількість пропусків практичних (за спаданням), кількість пропусків, прізвище, ім’я.

\bigskip\textit{Обмеження}

Навчальний тиждень може приймати цілі значення від 1 до 19.
День тижня є цілим від 1 до 5. Пара є цілим від 1 до 4.
Курс є цілим від 1 до 4. Аудиторія є додатнім цілим.
Вид занять може приймати значення:
лекц. --- лекція, п --- практичне, сем. --- семінар, ЛАБ --- лабораторне.

Решта полів є непорожніми рядками.
Рядки починатися та закінчуватися білими символами не можуть.

Групи кодуються в межах курсу.
Код групи складається не більш ніж з 5 символів.
Може містити літери алфавіту, десяткові цифри та дефіс.

Предмет є рядком, може мати від 6 до 24 символів включно.
Може містити літери алфавіту, десяткові цифри, апостроф, дефіс, пробіл.
Предмети різних курсів вважати різними навіть за однакових назв.

Прізвище, ім’я та по-батькові (за наявності у варіанті)
можуть мати від 6 до 25, 28, 24 відповідно символів включно кожне.
Можуть містити літери алфавіту, апостроф, пробіл та дефіс.
Студенти однозначно ідентифікуються прізвищем, ім’ям та по-батькові (за наявності у варіанті).
Жоден студент не вчиться одночасно на різних курсах чи в різних групах.

Види занять сортуються так: лабораторне < лекція < практичне < семінар.

%\bigskip%\textit{Для деяких варіантів може знадобитись}




\par
%end
\newpage
%begin
{\center Варіант  \No \textbf { 109 }\par}
\bigskip
%type = absense
%variant = 41
Обробляються дані семестрового журналу перевірок відвідування занять студентами деканатом (фіксуються відсутні). 

\underline{Записи основного файлу містять поля}: по-батькові, вид занять, ім’я, предмет, курс, код групи, день тижня, пара, навчальний тиждень, прізвище, аудиторія.

\underline{У допоміжному файлі наявні ключі}: загальна кількість записів, назва предмету, що пропускався найбільше.

\underline{Знайти} всі предмети, що пропускалися тільки на перших трьох навчальних тижнях. Вивести по кожному з них інформацію:\par
{\noindent}--- на першому рядку:\par
курс, предмет, кількість пропусків на першому тижні;

{\noindent}--- на наступних рядках, починаючи з табуляції, вивести пропуски предмета (по одному на рядок):\par
тиждень, день, пара, вид занять, прізвище, ім’я, по-батькові

{\noindent}у такому сортуванні: тиждень, прізвище, ім’я, по-батькові, день, пара, вид занять.

%Предмети сортувати: курс, кількість пропусків(за спаданням), предмет.

\bigskip\textit{Обмеження}

Навчальний тиждень може приймати цілі значення від 1 до 15.
День тижня є цілим від 1 до 5. Пара є цілим від 1 до 4.
Курс є цілим від 1 до 4. Аудиторія є додатнім цілим.
Вид занять може приймати значення:
Lect --- лекція, Пр --- практичне, sem --- семінар, Лаб. --- лабораторне.

Решта полів є непорожніми рядками.
Рядки починатися та закінчуватися білими символами не можуть.

Групи кодуються в межах курсу.
Код групи складається не більш ніж з 5 символів.
Може містити літери алфавіту, десяткові цифри та дефіс.

Предмет є рядком, може мати від 3 до 22 символів включно.
Може містити літери алфавіту, десяткові цифри, апостроф, дефіс, пробіл.
Предмети різних курсів вважати різними навіть за однакових назв.

Прізвище, ім’я та по-батькові (за наявності у варіанті)
можуть мати від 3 до 20, 24, 29 відповідно символів включно кожне.
Можуть містити літери алфавіту, апостроф, пробіл та дефіс.
Студенти однозначно ідентифікуються прізвищем, ім’ям та по-батькові (за наявності у варіанті).
Жоден студент не вчиться одночасно на різних курсах чи в різних групах.

Види занять сортуються так: семінар < лабораторне < лекція < практичне.

%\bigskip%\textit{Для деяких варіантів може знадобитись}




\par
%end
\newpage
%begin
{\center Варіант  \No \textbf { 110 }\par}
\bigskip
%type = meteo
%variant = 88
Обробляються дані спостережень за погодою. 

\underline{Записи основного файлу містять поля}: рік, сила вітру, мінімальна денна температура, кількість опадів, день, місяць, відносна вологість, середня денна температура, код метеостанції, максимальна денна температура.

\underline{У допоміжному файлі наявні ключі}: середня вологість по всіх спостереженнях, загальна кількість спостережень.

\underline{Знайти} 6 дат, по яких фіксувалася найбільша максимальна денна температура. Вивести по кожному з них інформацію:\par
{\noindent}--- на першому рядку:\par
день, рік, місяць, відношення найбільшої максимальної температури до найменшої мінімальної (по всіх метеостанціях, округлене до 2 знаків);

{\noindent}--- на наступних рядках, починаючи з табуляції, вивести для дати спостереження метеостанції, де вологість не перевищувала 45 (по одному на рядок):\par
вологість, метеостанція, сила вітру, середня денна температура

{\noindent}у такому сортуванні: середня денна температура, метеостанція.

%Дати сортувати: відношення (округлене), рік, місяць, день.

\bigskip\textit{Обмеження}

Код метеостанції є рядком довжини від 4 до 10 включно.
Може містити літери алфавіту, десяткові цифри, підкреслення. Решта полів числові.

Рік є чотиризначним цілим.
Кількість опадів є дійсним значенням, подається у форматі з 2 знаками після крапки.
Сила вітру є цілим значенням.
Температури є дійсними значеннями, подаються у форматі з 3 знаками після крапки.
Цілі значення подаються без десяткової крапки.
Для дійсних використовується зображення з фіксованою крапкою.

\textit{Вихідне форматування дійсних}

Усі дійсні значення виводяться у форматі з фіксованою крапкою з тією ж кількістю знаків після крапки, що і у вхідному зображенні.

Якщо у варіанті раніше не було зазначено, то \underline{обчислювані програмою} середні значення виводяться у форматі з 1 знаком після крапки.
%\bigskip%\textit{Для деяких варіантів може знадобитись}




\par
%end
\newpage
%begin
{\center Варіант  \No \textbf { 111 }\par}
\bigskip
%type = tasks
%variant = 2
Обробляються результати розв’язування студентами задач на контрольних.
Одиничний факт роз\-в’я\-зу\-вання студентом задачі надалі розуміється як спроба.

\underline{Записи основного файлу містять поля}: рік, по-батькові, ім’я, прізвище, умова задачі, відсоток набраних за розв’язання задачі балів.

\underline{У допоміжному файлі наявні ключі}: кількість спроб з відсотком не менше 50\%, кількість записів.

\underline{Знайти} всі задачі, які були найкраще розв’язані в середньому (середнє значення округлити до цілого). Вивести по кожному з них інформацію:\par
{\noindent}--- на першому рядку:\par
умова, середня розв’язуваність, відсоток спроб краще середнього (з округленням до цілого);

{\noindent}--- на наступних рядках, починаючи з табуляції, вивести спроби, що набрали якнайменш 75\% за цю задачу (по одному на рядок):\par
відсоток набраних за задачу балів, прізвище, ім’я, по-батькові, рік

{\noindent}у такому сортуванні: відсоток набраних за задачу балів (за спаданням), прізвище, ім’я, по-батькові, рік.

%Задачі сортувати: середня розв’язуваність (за спаданням), умова.

\bigskip\textit{Обмеження}

Студент міг розв’язувати задачу кілька разів.

Відсоток набраних за розв’язання задачі балів є цілим від 0 до 100.
День, місяць, рік є числовими полями, що зображують дату.
Раніше 2001 року статистика не велася. Рік є 4-значним.
Решта полів є непорожніми рядками.
Рядки починатися та закінчуватися білими символами не можуть.

Прізвище, ім’я та по-батькові (за наявності у варіанті) можуть мати до 29, 27, 28 відповідно символів включно кожне. 
Крім літер алфавіту можуть містити апостроф, дефіс та пробіл.

Умова задачі може мати від 4 до 70 символів включно.
Крім літер алфавіту може містити подвійні лапки, десяткові цифри, пробіл, дефіс,
а також символи . , \_ \{ \} \$ $+$ $*$ {\textbackslash} ( ) \% $[$ $]$ .

Студенти однозначно ідентифікуються прізвищем, ім’ям та по-батькові (за наявності у варіанті).
Задачі однозначно ідентифікуються умовою.

\bigskip\textit{Для деяких варіантів може знадобитись}

Розв'язуваність задачі --- середній відсоток набраних за задачу балів.

Розв'язуваність студента --- середній відсоток набраних студентом балів по всіх його спробах.

Розв'язуваність --- середній відсоток набраних студентами за задачі балів (рахується по всім даним).




\par
%end
\newpage
%begin
{\center Варіант  \No \textbf { 112 }\par}
\bigskip
%type = session
%variant = 54
Обробляються результати складання студентами екзаменів зимової сесії.

\underline{Записи основного файлу містять поля}: код групи, набрані на екзамені бали, набрані впродовж семестру бали, ім’я, прізвище, предмет, номер залікової книжки, по-батькові, сумарна оцінка в балах за 100-бальною системою, оцінка за державною шкалою.

\underline{У допоміжному файлі наявні ключі}: кількість оцінок «відмінно», кількість студентів, що складали сесію.

\underline{Знайти} 3 найстабільніших студенти. Вивести по кожному з них інформацію:\par
{\noindent}--- на першому рядку:\par
прізвище, ім’я, номер залікової книжки, найменший бал, найбільший бал;

{\noindent}--- на наступних рядках, починаючи з табуляції, вивести для студента всі його результати (по одному на рядок):\par
предмет, сумарна оцінка, державна оцінка прописом, бали на екзамені, бали в семестрі

{\noindent}у такому сортуванні: сумарна оцінка (за спаданням), предмет.

%Студентів сортувати: рейтинг за балами (після округлення), прізвище, ім’я, по-батькові, номер залікової книжки.

\bigskip\textit{Обмеження}

Студент не може складати предмет більше одного разу.

Бали та оцінка є значеннями цілих числових типів. Решта полів є непорожніми рядками.
Рядки розпочинатись та закінчуватись білими символами не можуть.

Оцінка за державною шкалою може приймати значення: 2-5~--- як зазвичай, 0~--- не з’явився,$-1$~--- не допущений. 
Оцінки 3-5 вважаються задовільними, решта~--- незадовільними. 
Бали є цілими та невід’ємними.
Набрані впродовж семестру бали не можуть перевищувати 60. 
Набрані на екзамені бали не можуть перевищувати 40. 
Набрані на екзамені бали повинні бути не менше 24 або дорівнювати 0(в останньому випадку оцінка не може бути задовільною). 
Сумарна оцінка в балах є сумою балів, набраних у семестрі та на екзамені.
Сумарна оцінка в балах та за державною шкалою мають бути узгоджені між собою.

Прізвище, ім’я та по-батькові (якщо наявне у варіанті) можуть мати до 30, 29, 21 відповідно символів включно кожне. 
Крім літер алфавіту можуть містити апостроф, дефіс та пробіл.

Предмет може мати від 3 до 61 символів включно. 
Крім літер алфавіту може містити подвійні лапки, пробіл та дефіс.

Код групи є непорожнім та складається не більш ніж з 6 літер. 
Крім літер алфавіту може містити десяткові цифри та дефіс.


Номер залікової книжки однозначно ідентифікує студента та складається з 7 десяткових цифр. 

\bigskip\textit{Для деяких варіантів може знадобитись}

Середній бал/оцінка за предметом визначається як середнє арифметичнесумарних балів/державних оцінок за цим предметом. 
Середній бал/оцінка студента визначається як середнє арифметичне сумарних балів/державних оцінокцього студента. 
За обчислення середньої оцінки недопущені та ті, що не з’явилися, рахуються як 2.
У решті випадків недопуски підсумовуються як $-1$, а «не з’явився» як 0.
Якщо студент не гірше ніж задовільно склав усі предмети, 
то його рейтинг за балами/оцінкою визначається як середній бал/оцінка; інакше дорівнює 0.

Стабільність студента визначається як модуль різниці між його кращим та найгіршим сумарним балом.
Стабільність предмета визначається як модуль різниці між найкращим та найгіршимсумарним балом за цим предметом.






\par
%end
\newpage
%begin
{\center Варіант  \No \textbf { 113 }\par}
\bigskip
%type = meteo
%variant = 97
Обробляються дані спостережень за погодою. 

\underline{Записи основного файлу містять поля}: день, місяць, відносна вологість, рік, кількість опадів, максимальна денна температура, код метеостанції, середня денна температура, сила вітру, мінімальна денна температура.

\underline{У допоміжному файлі наявні ключі}: середня сила вітру, сума всіх місяців.

\underline{Знайти} дати, по яких фіксувалася мінімальна денна температура менше ніж в середньому по всіх спостереженнях. Вивести по кожному з них інформацію:\par
{\noindent}--- на першому рядку:\par
день, місяць, рік, кількість опадів у цьому місяці, максимальна сила вітру у цьому місяці;

{\noindent}--- на наступних рядках, починаючи з табуляції, вивести для дати спостереження метеостанції, по яких в цю дату не було опадів (по одному на рядок):\par
середня денна температура, максимальна денна температура, мінімальна денна температура, сила вітру, метеостанція

{\noindent}у такому сортуванні: середня денна температура, метеостанція.

%Дати сортувати: місяць, день, рік.

\bigskip\textit{Обмеження}

Код метеостанції є рядком довжини від 5 до 12 включно.
Може містити літери алфавіту, десяткові цифри, підкреслення. Решта полів числові.

Рік є чотиризначним цілим.
Кількість опадів є цілим значенням.
Сила вітру є дійсним значенням, подається у форматі з 2 знаками після крапки.
Температури є дійсними значеннями, подаються у форматі з 1 знаком після крапки.
Цілі значення подаються без десяткової крапки.
Для дійсних використовується зображення з фіксованою крапкою.

\textit{Вихідне форматування дійсних}

Усі дійсні значення виводяться у форматі з фіксованою крапкою з тією ж кількістю знаків після крапки, що і у вхідному зображенні.

Якщо у варіанті раніше не було зазначено, то \underline{обчислювані програмою} середні значення виводяться у форматі з 3 знаками після крапки.
%\bigskip%\textit{Для деяких варіантів може знадобитись}




\par
%end
\newpage
%begin
{\center Варіант  \No \textbf { 114 }\par}
\bigskip
%type = tasks
%variant = -3
Обробляються результати розв’язування студентами задач на контрольних.
Одиничний факт роз\-в’я\-зу\-вання студентом задачі надалі розуміється як спроба.

\underline{Записи основного файлу містять поля}: умова задачі, по-батькові, відсоток набраних за розв’язання задачі балів, день, ім’я, прізвище, місяць, рік.

\underline{У допоміжному файлі наявні ключі}: загальна кількість записів, довжина найкоротшої умови.

\underline{Знайти} всіх студентів, що для кожної розв'язуваної ними задачі мають спробу з відсотком не менше 60\%. Вивести по кожному з них інформацію:\par
{\noindent}--- на першому рядку:\par
прізвище, ім’я, по-батькові, найкраща спроба, найгірша спроба, кількість спроб;

{\noindent}--- на наступних рядках, починаючи з табуляції, вивести розв'язувані ними задачі (по одному на рядок):\par
кращий відсоток набраних за задачу балів, умова, рік коли це трапилось

{\noindent}у такому сортуванні: відсоток (за спаданням), рік (за спаданням), умова.

%Студентів сортувати: найкраща спроба, прізвище, ім’я, по-батькові.

\bigskip\textit{Обмеження}

Студент міг розв’язувати задачу кілька разів.

Відсоток набраних за розв’язання задачі балів є цілим від 0 до 100.
День, місяць, рік є числовими полями, що зображують дату.
Раніше 2001 року статистика не велася. Рік є 4-значним.
Решта полів є непорожніми рядками.
Рядки починатися та закінчуватися білими символами не можуть.

Прізвище, ім’я та по-батькові (за наявності у варіанті) можуть мати до 23, 22, 23 відповідно символів включно кожне. 
Крім літер алфавіту можуть містити апостроф, дефіс та пробіл.

Умова задачі може мати від 5 до 58 символів включно.
Крім літер алфавіту може містити подвійні лапки, десяткові цифри, пробіл, дефіс,
а також символи . , \_ \{ \} \$ $+$ $*$ {\textbackslash} ( ) \% $[$ $]$ .

Студенти однозначно ідентифікуються прізвищем, ім’ям та по-батькові (за наявності у варіанті).
Задачі однозначно ідентифікуються умовою.

\bigskip\textit{Для деяких варіантів може знадобитись}

Розв'язуваність задачі --- середній відсоток набраних за задачу балів.

Розв'язуваність студента --- середній відсоток набраних студентом балів по всіх його спробах.

Розв'язуваність --- середній відсоток набраних студентами за задачі балів (рахується по всім даним).




\par
%end
\newpage
%begin
{\center Варіант  \No \textbf { 115 }\par}
\bigskip
%type = meteo
%variant = 89
Обробляються дані спостережень за погодою. 

\underline{Записи основного файлу містять поля}: рік, кількість опадів, середня денна температура, місяць, день, відносна вологість, максимальна денна температура, сила вітру, мінімальна денна температура, код метеостанції.

\underline{У допоміжному файлі наявні ключі}: рік завершення спостережень, місяць завершення спостережень, середня сила вітру.

\underline{Знайти} 8 дат, по яких фіксувалася найменша максимальна  денна температура. Вивести по кожному з них інформацію:\par
{\noindent}--- на першому рядку:\par
відношення кількості опадів (середньої за дату) до вологості (середньої за дату) (округлене до 3 знаків), місяць, день, рік, середня максимальна температура;

{\noindent}--- на наступних рядках, починаючи з табуляції, вивести для дати спостереження метеостанції, де сила вітру була принаймні тричі більша за вологість (по одному на рядок):\par
середня денна температура, максимальна денна температура, мінімальна денна температура, сила вітру, вологість, метеостанція

{\noindent}у такому сортуванні: сила вітру, метеостанція.

%Дати сортувати: відношення (округлене, за спаданням), місяць, день, рік.

\bigskip\textit{Обмеження}

Код метеостанції є рядком довжини від 3 до 8 включно.
Може містити літери алфавіту, десяткові цифри, підкреслення. Решта полів числові.

Рік є чотиризначним цілим.
Кількість опадів є дійсним значенням, подається у форматі з 1 знаком після крапки.
Сила вітру є дійсним значенням, подається у форматі з 2 знаками після крапки.
Температури є дійсними значеннями, подаються у форматі з 2 знаками після крапки.
Цілі значення подаються без десяткової крапки.
Для дійсних використовується зображення з фіксованою крапкою.

\textit{Вихідне форматування дійсних}

Усі дійсні значення виводяться у форматі з фіксованою крапкою з тією ж кількістю знаків після крапки, що і у вхідному зображенні.

Якщо у варіанті раніше не було зазначено, то \underline{обчислювані програмою} середні значення виводяться у форматі з 1 знаком після крапки.
%\bigskip%\textit{Для деяких варіантів може знадобитись}




\par
%end
\newpage
%begin
{\center Варіант  \No \textbf { 116 }\par}
\bigskip
%type = session
%variant = 58
Обробляються результати складання студентами екзаменів зимової сесії.

\underline{Записи основного файлу містять поля}: оцінка за державною шкалою, ім’я, по-батькові, набрані на екзамені бали, прізвище, сумарна оцінка в балах за 100-бальною системою, предмет, код групи, номер залікової книжки.

\underline{У допоміжному файлі наявні ключі}: кількість студентів, у яких є «хвости», загальна кількість записів.

\underline{Знайти} всіх студентів, що отримали хоча б з одного предмету сумарний бал 100. Вивести по кожному з них інформацію:\par
{\noindent}--- на першому рядку:\par
кількість предметів на 100 балів, номер залікової книжки, прізвище, ім’я, середній бал (округлений до 1 знаку);

{\noindent}--- на наступних рядках, починаючи з табуляції, вивести для студента предмети, що отримали менше 100 балів (по одному на рядок):\par
сумарна оцінка, предмет

{\noindent}у такому сортуванні: сумарна оцінка (за спаданням), предмет.

%Студентів сортувати: кількість предметів на 100 балів (за спаданням), середній бал (після округлення), прізвище, ім’я, по-батькові.

\bigskip\textit{Обмеження}

Студент не може складати предмет більше одного разу.

Бали та оцінка є значеннями цілих числових типів. Решта полів є непорожніми рядками.
Рядки розпочинатись та закінчуватись білими символами не можуть.

Оцінка за державною шкалою може приймати значення: 2-5~--- як зазвичай, 0~--- не з’явився,$-1$~--- не допущений. 
Оцінки 3-5 вважаються задовільними, решта~--- незадовільними. 
Бали є цілими та невід’ємними.
Набрані впродовж семестру бали не можуть перевищувати 60. 
Набрані на екзамені бали не можуть перевищувати 40. 
Набрані на екзамені бали повинні бути не менше 24 або дорівнювати 0(в останньому випадку оцінка не може бути задовільною). 
Сумарна оцінка в балах є сумою балів, набраних у семестрі та на екзамені.
Сумарна оцінка в балах та за державною шкалою мають бути узгоджені між собою.

Прізвище, ім’я та по-батькові (якщо наявне у варіанті) можуть мати до 26, 22, 25 відповідно символів включно кожне. 
Крім літер алфавіту можуть містити апостроф, дефіс та пробіл.

Предмет може мати від 5 до 59 символів включно. 
Крім літер алфавіту може містити подвійні лапки, пробіл та дефіс.

Код групи є непорожнім та складається не більш ніж з 4 літер. 
Крім літер алфавіту може містити десяткові цифри та дефіс.


Номер залікової книжки однозначно ідентифікує студента та складається з 6 десяткових цифр. 

\bigskip\textit{Для деяких варіантів може знадобитись}

Середній бал/оцінка за предметом визначається як середнє арифметичнесумарних балів/державних оцінок за цим предметом. 
Середній бал/оцінка студента визначається як середнє арифметичне сумарних балів/державних оцінокцього студента. 
За обчислення середньої оцінки недопущені та ті, що не з’явилися, рахуються як 2.
У решті випадків недопуски підсумовуються як $-1$, а «не з’явився» як 0.
Якщо студент не гірше ніж задовільно склав усі предмети, 
то його рейтинг за балами/оцінкою визначається як середній бал/оцінка; інакше дорівнює 0.

Стабільність студента визначається як модуль різниці між його кращим та найгіршим сумарним балом.
Стабільність предмета визначається як модуль різниці між найкращим та найгіршимсумарним балом за цим предметом.






\par
%end
\newpage
%begin
{\center Варіант  \No \textbf { 117 }\par}
\bigskip
%type =     absense
%variant = 21
Обробляються дані семестрового журналу перевірок відвідування занять студентами деканатом (фіксуються відсутні). 

\underline{Записи основного файлу містять поля}: ім’я, по-батькові, день тижня, пара, вид занять, аудиторія, прізвище, курс, код групи, навчальний тиждень, предмет.

\underline{У допоміжному файлі наявні ключі}: загальна кількість зафіксованих пропусків, сума курсів.

\underline{Знайти} всіх студентів, які пропустили найбільшу кількість пар. Вивести по кожному з них інформацію:\par
{\noindent}--- на першому рядку:\par
прізвище, ім’я, по-батькові, кількість пропущених пар, кількість пропущених лекцій, кількість пропущених семінарів;

{\noindent}--- на наступних рядках, починаючи з табуляції, вивести пропуски студента (по одному на рядок):\par
тиждень, день, пара, аудиторія, предмет, вид занять

{\noindent}у такому сортуванні: предмет, тиждень, день, пара, вид занять.

%Студентів сортувати: кількість пропусків (за спаданням), кількість пропусків лекцій, прізвище, ім’я, по-батькові.

\bigskip\textit{Обмеження}

Навчальний тиждень може приймати цілі значення від 1 до 17.
День тижня є цілим від 1 до 5. Пара є цілим від 1 до 4.
Курс є цілим від 1 до 4. Аудиторія є додатнім цілим.
Вид занять може приймати значення:
lect --- лекція, P --- практичне, сем. --- семінар, Лаб. --- лабораторне.

Решта полів є непорожніми рядками.
Рядки починатися та закінчуватися білими символами не можуть.

Групи кодуються в межах курсу.
Код групи складається не більш ніж з 6 символів.
Може містити літери алфавіту, десяткові цифри та дефіс.

Предмет є рядком, може мати від 4 до 27 символів включно.
Може містити літери алфавіту, десяткові цифри, апостроф, дефіс, пробіл.
Предмети різних курсів вважати різними навіть за однакових назв.

Прізвище, ім’я та по-батькові (за наявності у варіанті)
можуть мати від 4 до 28, 27, 22 відповідно символів включно кожне.
Можуть містити літери алфавіту, апостроф, пробіл та дефіс.
Студенти однозначно ідентифікуються прізвищем, ім’ям та по-батькові (за наявності у варіанті).
Жоден студент не вчиться одночасно на різних курсах чи в різних групах.

Види занять сортуються так: практичне < семінар < лекція < лабораторне.

%\bigskip%\textit{Для деяких варіантів може знадобитись}




\par
%end
\newpage
%begin
{\center Варіант  \No \textbf { 118 }\par}
\bigskip
%type = absense
%variant = 39
Обробляються дані семестрового журналу перевірок відвідування занять студентами деканатом (фіксуються відсутні). 

\underline{Записи основного файлу містять поля}: аудиторія, навчальний тиждень, курс, код групи, предмет, вид занять, прізвище, ім’я, день тижня, пара.

\underline{У допоміжному файлі наявні ключі}: загальна кількість записів, FINE.

\underline{Знайти} всі предмети, що пропускалися більше, ніж в середньому. Вивести по кожному з них інформацію:\par
{\noindent}--- на першому рядку:\par
кількість пропусків, курс, предмет, кількість пропусків лабораторних;

{\noindent}--- на наступних рядках, починаючи з табуляції, вивести пропуски предмета (по одному на рядок):\par
прізвище, ім’я, тиждень, вид занять, кількість пропусків

{\noindent}у такому сортуванні: прізвище, ім’я, вид занять.

%Предмети сортувати: кількість пропусків (за спаданням), курс, предмет.

\bigskip\textit{Обмеження}

Навчальний тиждень може приймати цілі значення від 1 до 14.
День тижня є цілим від 1 до 5. Пара є цілим від 1 до 4.
Курс є цілим від 1 до 4. Аудиторія є додатнім цілим.
Вид занять може приймати значення:
Лекц. --- лекція, Pr --- практичне, s --- семінар, Lab --- лабораторне.

Решта полів є непорожніми рядками.
Рядки починатися та закінчуватися білими символами не можуть.

Групи кодуються в межах курсу.
Код групи складається не більш ніж з 4 символів.
Може містити літери алфавіту, десяткові цифри та дефіс.

Предмет є рядком, може мати від 6 до 30 символів включно.
Може містити літери алфавіту, десяткові цифри, апостроф, дефіс, пробіл.
Предмети різних курсів вважати різними навіть за однакових назв.

Прізвище, ім’я та по-батькові (за наявності у варіанті)
можуть мати від 6 до 29, 21, 21 відповідно символів включно кожне.
Можуть містити літери алфавіту, апостроф, пробіл та дефіс.
Студенти однозначно ідентифікуються прізвищем, ім’ям та по-батькові (за наявності у варіанті).
Жоден студент не вчиться одночасно на різних курсах чи в різних групах.

Види занять сортуються так: семінар < практичне < лекція < лабораторне.

%\bigskip%\textit{Для деяких варіантів може знадобитись}




\par
%end
\newpage
%begin
{\center Варіант  \No \textbf { 119 }\par}
\bigskip
%type = session
%variant = 69
Обробляються результати складання студентами екзаменів зимової сесії.

\underline{Записи основного файлу містять поля}: номер залікової книжки, оцінка за державною шкалою, предмет, по-батькові, набрані на екзамені бали, прізвище, ім’я, сумарна оцінка в балах за 100-бальною системою, код групи.

\underline{У допоміжному файлі наявні ключі}: кількість «відмінно», сума державних оцінок.

\underline{Знайти} всіх студентів, що щось не склали і при цьому отримували оцінку «відмінно». Вивести по кожному з них інформацію:\par
{\noindent}--- на першому рядку:\par
кількість «відмінно», кількість «хвостів», ім’я, прізвище, рейтинг за оцінкою (округлений до 1 знаку), номер залікової книжки;

{\noindent}--- на наступних рядках, починаючи з табуляції, вивести для студента предмети, з оцінкою нижче «відмінно» (по одному на рядок):\par
предмет, сумарна оцінка, державна оцінка

{\noindent}у такому сортуванні: предмет, сумарна оцінка.

%Студентів сортувати: кількість «відмінно», кількість «хвостів» (за спаданням), прізвище, номер залікової книжки.

\bigskip\textit{Обмеження}

Студент не може складати предмет більше одного разу.

Бали та оцінка є значеннями цілих числових типів. Решта полів є непорожніми рядками.
Рядки розпочинатись та закінчуватись білими символами не можуть.

Оцінка за державною шкалою може приймати значення: 2-5~--- як зазвичай, 0~--- не з’явився,$-1$~--- не допущений. 
Оцінки 3-5 вважаються задовільними, решта~--- незадовільними. 
Бали є цілими та невід’ємними.
Набрані впродовж семестру бали не можуть перевищувати 60. 
Набрані на екзамені бали не можуть перевищувати 40. 
Набрані на екзамені бали повинні бути не менше 24 або дорівнювати 0(в останньому випадку оцінка не може бути задовільною). 
Сумарна оцінка в балах є сумою балів, набраних у семестрі та на екзамені.
Сумарна оцінка в балах та за державною шкалою мають бути узгоджені між собою.

Прізвище, ім’я та по-батькові (якщо наявне у варіанті) можуть мати до 20, 20, 29 відповідно символів включно кожне. 
Крім літер алфавіту можуть містити апостроф, дефіс та пробіл.

Предмет може мати від 4 до 56 символів включно. 
Крім літер алфавіту може містити подвійні лапки, пробіл та дефіс.

Код групи є непорожнім та складається не більш ніж з 3 літер. 
Крім літер алфавіту може містити десяткові цифри та дефіс.


Номер залікової книжки однозначно ідентифікує студента та складається з 8 десяткових цифр. 

\bigskip\textit{Для деяких варіантів може знадобитись}

Середній бал/оцінка за предметом визначається як середнє арифметичнесумарних балів/державних оцінок за цим предметом. 
Середній бал/оцінка студента визначається як середнє арифметичне сумарних балів/державних оцінокцього студента. 
За обчислення середньої оцінки недопущені та ті, що не з’явилися, рахуються як 2.
У решті випадків недопуски підсумовуються як $-1$, а «не з’явився» як 0.
Якщо студент не гірше ніж задовільно склав усі предмети, 
то його рейтинг за балами/оцінкою визначається як середній бал/оцінка; інакше дорівнює 0.

Стабільність студента визначається як модуль різниці між його кращим та найгіршим сумарним балом.
Стабільність предмета визначається як модуль різниці між найкращим та найгіршимсумарним балом за цим предметом.






\par
%end
\newpage
%begin
{\center Варіант  \No \textbf { 120 }\par}
\bigskip
%type = meteo
%variant = 93
Обробляються дані спостережень за погодою. 

\underline{Записи основного файлу містять поля}: рік, сила вітру, середня денна температура, код метеостанції, місяць, день, кількість опадів, відносна вологість, максимальна денна температура, мінімальна денна температура.

\underline{У допоміжному файлі наявні ключі}: кількість спостережень за вересень, загальна кількість записів.

\underline{Знайти} дати, по яких фіксувалася відносна вологість більша ніж в середньому по всіх спостереженнях. Вивести по кожному з них інформацію:\par
{\noindent}--- на першому рядку:\par
найбільше та найменше значення вологості за дату, середня середньоденна температура, рік, місяць, день, кількість опадів;

{\noindent}--- на наступних рядках, починаючи з табуляції, вивести для дати спостереження метеостанції, де фіксувалися опади (по одному на рядок):\par
метеостанція, кількість опадів, вологість

{\noindent}у такому сортуванні: кількість опадів (за спаданням), метеостанція.

%Дати сортувати: найбільша вологість (за спаданням), рік, місяць, день.

\bigskip\textit{Обмеження}

Код метеостанції є рядком довжини від 4 до 7 включно.
Може містити літери алфавіту, десяткові цифри, підкреслення. Решта полів числові.

Рік є чотиризначним цілим.
Кількість опадів є цілим значенням.
Сила вітру є цілим значенням.
Температури є дійсними значеннями, подаються у форматі з 1 знаком після крапки.
Цілі значення подаються без десяткової крапки.
Для дійсних використовується зображення з фіксованою крапкою.

\textit{Вихідне форматування дійсних}

Усі дійсні значення виводяться у форматі з фіксованою крапкою з тією ж кількістю знаків після крапки, що і у вхідному зображенні.

Якщо у варіанті раніше не було зазначено, то \underline{обчислювані програмою} середні значення виводяться у форматі з 2 знаками після крапки.
%\bigskip%\textit{Для деяких варіантів може знадобитись}




\par
%end
\newpage
%begin
{\center Варіант  \No \textbf { 121 }\par}
\bigskip
%type = meteo
%variant = 105
Обробляються дані спостережень за погодою. 

\underline{Записи основного файлу містять поля}: рік, максимальна денна температура, сила вітру, середня денна температура, мінімальна денна температура, відносна вологість, код метеостанції, день, місяць, кількість опадів.

\underline{У допоміжному файлі наявні ключі}: кількість спостережень, найбільше значення кількості опадів, найменше значення кількості опадів.

\underline{Знайти} 6 метеостанцій, по яких у середньому була найвища максимальна денна температура. Вивести по кожному з них інформацію:\par
{\noindent}--- на першому рядку:\par
метеостанція, найбільша та найменша максимальні температури;

{\noindent}--- на наступних рядках, починаючи з табуляції, вивести для них дані по місяцях (по одному на рядок):\par
рік, місяць, кількість опадів, середня середньоденна температура

{\noindent}у такому сортуванні: місяць, рік.

%Метеостанції сортувати: кількість спостережень, найбільша максимальна температура, метеостанція.

\bigskip\textit{Обмеження}

Код метеостанції є рядком довжини від 5 до 6 включно.
Може містити літери алфавіту, десяткові цифри, підкреслення. Решта полів числові.

Рік є чотиризначним цілим.
Кількість опадів є дійсним значенням, подається у форматі з 2 знаками після крапки.
Сила вітру є дійсним значенням, подається у форматі з 1 знаком після крапки.
Температури є дійсними значеннями, подаються у форматі з 3 знаками після крапки.
Цілі значення подаються без десяткової крапки.
Для дійсних використовується зображення з фіксованою крапкою.

\textit{Вихідне форматування дійсних}

Усі дійсні значення виводяться у форматі з фіксованою крапкою з тією ж кількістю знаків після крапки, що і у вхідному зображенні.

Якщо у варіанті раніше не було зазначено, то \underline{обчислювані програмою} середні значення виводяться у форматі з 3 знаками після крапки.
%\bigskip%\textit{Для деяких варіантів може знадобитись}




\par
%end
\newpage
%begin
{\center Варіант  \No \textbf { 122 }\par}
\bigskip
%type = meteo
%variant = 109
Обробляються дані спостережень за погодою. 

\underline{Записи основного файлу містять поля}: код метеостанції, максимальна денна температура, мініма\-ль\-на денна температура, середня денна температура, кількість опадів, відносна вологість, сила вітру, рік, місяць, день.

\underline{У допоміжному файлі наявні ключі}: найбільша сила вітру, найменша сила вітру, загальна кількість спостережень.

\underline{Знайти} 5 метеостанцій, по яких у середньому випала найбільша кількість опадів. Вивести по кожному з них інформацію:\par
{\noindent}--- на першому рядку:\par
метеостанція, кількість спостережень, середня кількість опадів, сумарна кількість опадів;

{\noindent}--- на наступних рядках, починаючи з табуляції, вивести для них дані по датах, коли не було опадів (по одному на рядок):\par
середня денна температура, сила вітру, вологість, день, місяць, рік

{\noindent}у такому сортуванні: вологість (за спаданням), рік, місяць, день.

%Метеостанції сортувати: кількість спостережень (за спаданням), сумарна кількість опадів,  метеостанція.

\bigskip\textit{Обмеження}

Код метеостанції є рядком довжини від 4 до 10 включно.
Може містити літери алфавіту, десяткові цифри, підкреслення. Решта полів числові.

Рік є чотиризначним цілим.
Кількість опадів є дійсним значенням, подається у форматі з 1 знаком після крапки.
Сила вітру є дійсним значенням, подається у форматі з 1 знаком після крапки.
Температури є дійсними значеннями, подаються у форматі з 3 знаками після крапки.
Цілі значення подаються без десяткової крапки.
Для дійсних використовується зображення з фіксованою крапкою.

\textit{Вихідне форматування дійсних}

Усі дійсні значення виводяться у форматі з фіксованою крапкою з тією ж кількістю знаків після крапки, що і у вхідному зображенні.

Якщо у варіанті раніше не було зазначено, то \underline{обчислювані програмою} середні значення виводяться у форматі з 1 знаком після крапки.
%\bigskip%\textit{Для деяких варіантів може знадобитись}




\par
%end
\newpage
%begin
{\center Варіант  \No \textbf { 123 }\par}
\bigskip
%type =     session
%variant = 52
Обробляються результати складання студентами екзаменів зимової сесії.

\underline{Записи основного файлу містять поля}: ім’я, номер залікової книжки, код групи, сумарна оцінка в балах за 100-бальною системою, оцінка за державною шкалою, прізвище, предмет, набрані на екзамені бали, по-батькові.

\underline{У допоміжному файлі наявні ключі}: загальна кількість студентів, сума оцінок за державною шкалою.

\underline{Знайти} 5 найкращих студентів (за рейтингом за балами). Вивести по кожному з них інформацію:\par
{\noindent}--- на першому рядку:\par
прізвище, ім’я, по-батькові, номер залікової книжки, рейтинг за балами (округлений до цілого), рейтинг за оцінкою (округлений до 1 знаку);

{\noindent}--- на наступних рядках, починаючи з табуляції, вивести для студента всі його результати (по одному на рядок):\par
предмет, сумарна оцінка, державна оцінка, бали на екзамені, бали в семестрі

{\noindent}у такому сортуванні: сумарна оцінка, предмет.

%Студентів сортувати: рейтинг за балами (після округлення), прізвище, ім’я, по-батькові, номер залікової книжки.

\bigskip\textit{Обмеження}

Студент не може складати предмет більше одного разу.

Бали та оцінка є значеннями цілих числових типів. Решта полів є непорожніми рядками.
Рядки розпочинатись та закінчуватись білими символами не можуть.

Оцінка за державною шкалою може приймати значення: 2-5~--- як зазвичай, 0~--- не з’явився,$-1$~--- не допущений. 
Оцінки 3-5 вважаються задовільними, решта~--- незадовільними. 
Бали є цілими та невід’ємними.
Набрані впродовж семестру бали не можуть перевищувати 60. 
Набрані на екзамені бали не можуть перевищувати 40. 
Набрані на екзамені бали повинні бути не менше 24 або дорівнювати 0(в останньому випадку оцінка не може бути задовільною). 
Сумарна оцінка в балах є сумою балів, набраних у семестрі та на екзамені.
Сумарна оцінка в балах та за державною шкалою мають бути узгоджені між собою.

Прізвище, ім’я та по-батькові (якщо наявне у варіанті) можуть мати до 29, 21, 20 відповідно символів включно кожне. 
Крім літер алфавіту можуть містити апостроф, дефіс та пробіл.

Предмет може мати від 4 до 62 символів включно. 
Крім літер алфавіту може містити подвійні лапки, пробіл та дефіс.

Код групи є непорожнім та складається не більш ніж з 6 літер. 
Крім літер алфавіту може містити десяткові цифри та дефіс.


Номер залікової книжки однозначно ідентифікує студента та складається з 6 десяткових цифр. 

\bigskip\textit{Для деяких варіантів може знадобитись}

Середній бал/оцінка за предметом визначається як середнє арифметичнесумарних балів/державних оцінок за цим предметом. 
Середній бал/оцінка студента визначається як середнє арифметичне сумарних балів/державних оцінокцього студента. 
За обчислення середньої оцінки недопущені та ті, що не з’явилися, рахуються як 2.
У решті випадків недопуски підсумовуються як $-1$, а «не з’явився» як 0.
Якщо студент не гірше ніж задовільно склав усі предмети, 
то його рейтинг за балами/оцінкою визначається як середній бал/оцінка; інакше дорівнює 0.

Стабільність студента визначається як модуль різниці між його кращим та найгіршим сумарним балом.
Стабільність предмета визначається як модуль різниці між найкращим та найгіршимсумарним балом за цим предметом.






\par
%end
\newpage
%begin
{\center Варіант  \No \textbf { 124 }\par}
\bigskip
%type = absense
%variant = 36
Обробляються дані семестрового журналу перевірок відвідування занять студентами деканатом (фіксуються відсутні). 

\underline{Записи основного файлу містять поля}: \ навчальний\ тиждень,\ по-батькові,\ вид\ занять,\ прізвище, ім’я, аудиторія, предмет, день тижня, пара, курс, код групи.

\underline{У допоміжному файлі наявні ключі}: довжина найкоротшого прізвища, загальна кількість записів.

\underline{Знайти} всіх студентів, що більше інших пропускали заняття по п’ятницях. Вивести по кожному з них інформацію:\par
{\noindent}--- на першому рядку:\par
кількість пропущених пар, кількість пропущених пар по п’ятницях, прізвище, ім’я, по-батькові;

{\noindent}--- на наступних рядках, починаючи з табуляції, вивести пропуски студента саме по п’ятницях (по одному на рядок):\par
тиждень, пара, аудиторія, вид занять, предмет

{\noindent}у такому сортуванні: тиждень, пара, аудиторія, предмет, вид занять.

%Студентів сортувати: кількість пропущених пар, прізвище, ім’я, по-батькові.

\bigskip\textit{Обмеження}

Навчальний тиждень може приймати цілі значення від 1 до 18.
День тижня є цілим від 1 до 5. Пара є цілим від 1 до 4.
Курс є цілим від 1 до 4. Аудиторія є додатнім цілим.
Вид занять може приймати значення:
Le --- лекція, p --- практичне, Sem --- семінар, lab --- лабораторне.

Решта полів є непорожніми рядками.
Рядки починатися та закінчуватися білими символами не можуть.

Групи кодуються в межах курсу.
Код групи складається не більш ніж з 4 символів.
Може містити літери алфавіту, десяткові цифри та дефіс.

Предмет є рядком, може мати від 5 до 31 символів включно.
Може містити літери алфавіту, десяткові цифри, апостроф, дефіс, пробіл.
Предмети різних курсів вважати різними навіть за однакових назв.

Прізвище, ім’я та по-батькові (за наявності у варіанті)
можуть мати від 5 до 23, 29, 26 відповідно символів включно кожне.
Можуть містити літери алфавіту, апостроф, пробіл та дефіс.
Студенти однозначно ідентифікуються прізвищем, ім’ям та по-батькові (за наявності у варіанті).
Жоден студент не вчиться одночасно на різних курсах чи в різних групах.

Види занять сортуються так: лабораторне < семінар < практичне < лекція.

%\bigskip%\textit{Для деяких варіантів може знадобитись}




\par
%end
\newpage
%begin
{\center Варіант  \No \textbf { 125 }\par}
\bigskip
%type = meteo
%variant = 102
Обробляються дані спостережень за погодою. 

\underline{Записи основного файлу містять поля}: сила вітру, місяць, день, кількість опадів, максимальна денна температура, код метеостанції, середня денна температура, рік, мінімальна денна температура, відносна вологість.

\underline{У допоміжному файлі наявні ключі}: кількість спостережень, останній рік спостережень.

\underline{Знайти} дати, по яких середня мінімальна денна температура була менше ніж в середньому по всіх спостереженнях. Вивести по кожному з них інформацію:\par
{\noindent}--- на першому рядку:\par
місяць, день, середня кількість опадів, середня мінімальна температура, рік;

{\noindent}--- на наступних рядках, починаючи з табуляції, вивести для дати спостереження метеостанції, по яких в цей день опадів було менше 3 (по одному на рядок):\par
середня денна температура, метеостанція, сила вітру

{\noindent}у такому сортуванні: середня денна температура, метеостанція.

%Дати сортувати: рік, середня кількість опадів, місяць, день.

\bigskip\textit{Обмеження}

Код метеостанції є рядком довжини від 5 до 11 включно.
Може містити літери алфавіту, десяткові цифри, підкреслення. Решта полів числові.

Рік є чотиризначним цілим.
Кількість опадів є дійсним значенням, подається у форматі з 2 знаками після крапки.
Сила вітру є цілим значенням.
Температури є дійсними значеннями, подаються у форматі з 1 знаком після крапки.
Цілі значення подаються без десяткової крапки.
Для дійсних використовується зображення з фіксованою крапкою.

\textit{Вихідне форматування дійсних}

Усі дійсні значення виводяться у форматі з фіксованою крапкою з тією ж кількістю знаків після крапки, що і у вхідному зображенні.

Якщо у варіанті раніше не було зазначено, то \underline{обчислювані програмою} середні значення виводяться у форматі з 3 знаками після крапки.
%\bigskip%\textit{Для деяких варіантів може знадобитись}




\par
%end
\newpage
%begin
{\center Варіант  \No \textbf { 126 }\par}
\bigskip
%type = session
%variant = 67
Обробляються результати складання студентами екзаменів зимової сесії.

\underline{Записи основного файлу містять поля}: прізвище, предмет, код групи, набрані на екзамені бали, набрані впродовж семестру бали, оцінка за державною шкалою, ім’я, номер залікової книжки, сумарна оцінка в балах за 100-бальною системою, по-батькові.

\underline{У допоміжному файлі наявні ключі}: кількість студентів без «хвостів», загальна кількість записів.

\underline{Знайти} студентів, що в сесію отримали хоча б по одній оцінці «відмінно», «добре», «задовільно». Вивести по кожному з них інформацію:\par
{\noindent}--- на першому рядку:\par
кількість «добре», ім’я, прізвище, рейтинг за балами (округлений до цілого), номер залікової книжки;

{\noindent}--- на наступних рядках, починаючи з табуляції, вивести для студента всі його результати (по одному на рядок):\par
бали на екзамені, державна оцінка

{\noindent}у такому сортуванні: бали на екзамені, сумарна оцінка, предмет.

%Студентів сортувати: кількість «добре», номер залікової книжки.

\bigskip\textit{Обмеження}

Студент не може складати предмет більше одного разу.

Бали та оцінка є значеннями цілих числових типів. Решта полів є непорожніми рядками.
Рядки розпочинатись та закінчуватись білими символами не можуть.

Оцінка за державною шкалою може приймати значення: 2-5~--- як зазвичай, 0~--- не з’явився,$-1$~--- не допущений. 
Оцінки 3-5 вважаються задовільними, решта~--- незадовільними. 
Бали є цілими та невід’ємними.
Набрані впродовж семестру бали не можуть перевищувати 60. 
Набрані на екзамені бали не можуть перевищувати 40. 
Набрані на екзамені бали повинні бути не менше 24 або дорівнювати 0(в останньому випадку оцінка не може бути задовільною). 
Сумарна оцінка в балах є сумою балів, набраних у семестрі та на екзамені.
Сумарна оцінка в балах та за державною шкалою мають бути узгоджені між собою.

Прізвище, ім’я та по-батькові (якщо наявне у варіанті) можуть мати до 23, 25, 24 відповідно символів включно кожне. 
Крім літер алфавіту можуть містити апостроф, дефіс та пробіл.

Предмет може мати від 5 до 53 символів включно. 
Крім літер алфавіту може містити подвійні лапки, пробіл та дефіс.

Код групи є непорожнім та складається не більш ніж з 5 літер. 
Крім літер алфавіту може містити десяткові цифри та дефіс.


Номер залікової книжки однозначно ідентифікує студента та складається з 7 десяткових цифр. 

\bigskip\textit{Для деяких варіантів може знадобитись}

Середній бал/оцінка за предметом визначається як середнє арифметичнесумарних балів/державних оцінок за цим предметом. 
Середній бал/оцінка студента визначається як середнє арифметичне сумарних балів/державних оцінокцього студента. 
За обчислення середньої оцінки недопущені та ті, що не з’явилися, рахуються як 2.
У решті випадків недопуски підсумовуються як $-1$, а «не з’явився» як 0.
Якщо студент не гірше ніж задовільно склав усі предмети, 
то його рейтинг за балами/оцінкою визначається як середній бал/оцінка; інакше дорівнює 0.

Стабільність студента визначається як модуль різниці між його кращим та найгіршим сумарним балом.
Стабільність предмета визначається як модуль різниці між найкращим та найгіршимсумарним балом за цим предметом.






\par
%end
\newpage
%begin
{\center Варіант  \No \textbf { 127 }\par}
\bigskip
%type = session
%variant = 80
Обробляються результати складання студентами екзаменів зимової сесії.

\underline{Записи основного файлу містять поля}: набрані на екзамені бали, по-батькові, номер залікової книжки, прізвище, оцінка за державною шкалою, предмет, код групи, набрані впродовж семестру бали, сумарна оцінка в балах за 100-бальною системою, ім’я.

\underline{У допоміжному файлі наявні ключі}: сума державних оцінок, загальна кількість предметів.

\underline{Знайти} предмети, які складалися найбільшою кількістю студентів. Вивести по кожному з них інформацію:\par
{\noindent}--- на першому рядку:\par
предмет, середній бал за предметом (округлений до цілого), кількість студентів, що складали;

{\noindent}--- на наступних рядках, починаючи з табуляції, вивести для предмета всіх студентів, що його складали (по одному на рядок):\par
група, середній бал по групі (округлений до 1 знаку), кількість студентів

{\noindent}у такому сортуванні: група.

%Предмети сортувати: стабільність, середній бал (округлений), предмет.

\bigskip\textit{Обмеження}

Студент не може складати предмет більше одного разу.

Бали та оцінка є значеннями цілих числових типів. Решта полів є непорожніми рядками.
Рядки розпочинатись та закінчуватись білими символами не можуть.

Оцінка за державною шкалою може приймати значення: 2-5~--- як зазвичай, 0~--- не з’явився,$-1$~--- не допущений. 
Оцінки 3-5 вважаються задовільними, решта~--- незадовільними. 
Бали є цілими та невід’ємними.
Набрані впродовж семестру бали не можуть перевищувати 60. 
Набрані на екзамені бали не можуть перевищувати 40. 
Набрані на екзамені бали повинні бути не менше 24 або дорівнювати 0(в останньому випадку оцінка не може бути задовільною). 
Сумарна оцінка в балах є сумою балів, набраних у семестрі та на екзамені.
Сумарна оцінка в балах та за державною шкалою мають бути узгоджені між собою.

Прізвище, ім’я та по-батькові (якщо наявне у варіанті) можуть мати до 24, 30, 22 відповідно символів включно кожне. 
Крім літер алфавіту можуть містити апостроф, дефіс та пробіл.

Предмет може мати від 3 до 54 символів включно. 
Крім літер алфавіту може містити подвійні лапки, пробіл та дефіс.

Код групи є непорожнім та складається не більш ніж з 7 літер. 
Крім літер алфавіту може містити десяткові цифри та дефіс.


Номер залікової книжки однозначно ідентифікує студента та складається з 8 десяткових цифр. 

\bigskip\textit{Для деяких варіантів може знадобитись}

Середній бал/оцінка за предметом визначається як середнє арифметичнесумарних балів/державних оцінок за цим предметом. 
Середній бал/оцінка студента визначається як середнє арифметичне сумарних балів/державних оцінокцього студента. 
За обчислення середньої оцінки недопущені та ті, що не з’явилися, рахуються як 2.
У решті випадків недопуски підсумовуються як $-1$, а «не з’явився» як 0.
Якщо студент не гірше ніж задовільно склав усі предмети, 
то його рейтинг за балами/оцінкою визначається як середній бал/оцінка; інакше дорівнює 0.

Стабільність студента визначається як модуль різниці між його кращим та найгіршим сумарним балом.
Стабільність предмета визначається як модуль різниці між найкращим та найгіршимсумарним балом за цим предметом.






\par
%end
\newpage
%begin
{\center Варіант  \No \textbf { 128 }\par}
\bigskip
%type = absense
%variant = 49
Обробляються дані семестрового журналу перевірок відвідування занять студентами деканатом (фіксуються відсутні). 

\underline{Записи основного файлу містять поля}: предмет, день тижня, пара, курс, прізвище, ім’я, аудиторія, вид занять, навчальний тиждень.

\underline{У допоміжному файлі наявні ключі}: кількість пропусків третіх пар, кількість різних предметів.

\underline{Знайти} всі предмети, які пропускалися тільки по понеділках. Вивести по кожному з них інформацію:\par
{\noindent}--- на першому рядку:\par
курс, предмет, кількість пропусків;

{\noindent}--- на наступних рядках, починаючи з табуляції, вивести пропуски предмета (по одному на рядок):\par
прізвище, ім’я, тиждень, кількість пропусків

{\noindent}у такому сортуванні: прізвище, ім’я, тиждень.

%Предмети сортувати: предмет, курс.

\bigskip\textit{Обмеження}

Навчальний тиждень може приймати цілі значення від 1 до 18.
День тижня є цілим від 1 до 5. Пара є цілим від 1 до 4.
Курс є цілим від 1 до 4. Аудиторія є додатнім цілим.
Вид занять може приймати значення:
Lecture --- лекція, Практ. --- практичне, Seminar --- семінар, ЛАБ --- лабораторне.

Решта полів є непорожніми рядками.
Рядки починатися та закінчуватися білими символами не можуть.

Групи кодуються в межах курсу.
Код групи складається не більш ніж з 6 символів.
Може містити літери алфавіту, десяткові цифри та дефіс.

Предмет є рядком, може мати від 6 до 40 символів включно.
Може містити літери алфавіту, десяткові цифри, апостроф, дефіс, пробіл.
Предмети різних курсів вважати різними навіть за однакових назв.

Прізвище, ім’я та по-батькові (за наявності у варіанті)
можуть мати від 6 до 24, 30, 30 відповідно символів включно кожне.
Можуть містити літери алфавіту, апостроф, пробіл та дефіс.
Студенти однозначно ідентифікуються прізвищем, ім’ям та по-батькові (за наявності у варіанті).
Жоден студент не вчиться одночасно на різних курсах чи в різних групах.

Види занять сортуються так: лекція < семінар < практичне < лабораторне.

%\bigskip%\textit{Для деяких варіантів може знадобитись}




\par
%end
\newpage
%begin
{\center Варіант  \No \textbf { 129 }\par}
\bigskip
%type = absense
%variant = 45
Обробляються дані семестрового журналу перевірок відвідування занять студентами деканатом (фіксуються відсутні). 

\underline{Записи основного файлу містять поля}: прізвище, курс, код групи, ім’я, аудиторія, вид занять, предмет, день тижня, пара, навчальний тиждень.

\underline{У допоміжному файлі наявні ключі}: загальна кількість предметів, кількість пропусків лабораторних занять.

\underline{Знайти} всі предмети, що пропускалися кожного навчального тижня. Вивести по кожному з них інформацію:\par
{\noindent}--- на першому рядку:\par
кількість пропусків, предмет, курс;

{\noindent}--- на наступних рядках, починаючи з табуляції, вивести пропуски предмета (по одному на рядок):\par
тиждень, вид занять, кількість пропусків

{\noindent}у такому сортуванні: тиждень, вид занять.

%Предмети сортувати: кількість пропусків, предмет, курс.

\bigskip\textit{Обмеження}

Навчальний тиждень може приймати цілі значення від 1 до 16.
День тижня є цілим від 1 до 5. Пара є цілим від 1 до 4.
Курс є цілим від 1 до 4. Аудиторія є додатнім цілим.
Вид занять може приймати значення:
ЛЕ --- лекція, Practice --- практичне, с. --- семінар, лаб. --- лабораторне.

Решта полів є непорожніми рядками.
Рядки починатися та закінчуватися білими символами не можуть.

Групи кодуються в межах курсу.
Код групи складається не більш ніж з 5 символів.
Може містити літери алфавіту, десяткові цифри та дефіс.

Предмет є рядком, може мати від 5 до 38 символів включно.
Може містити літери алфавіту, десяткові цифри, апостроф, дефіс, пробіл.
Предмети різних курсів вважати різними навіть за однакових назв.

Прізвище, ім’я та по-батькові (за наявності у варіанті)
можуть мати від 5 до 21, 23, 20 відповідно символів включно кожне.
Можуть містити літери алфавіту, апостроф, пробіл та дефіс.
Студенти однозначно ідентифікуються прізвищем, ім’ям та по-батькові (за наявності у варіанті).
Жоден студент не вчиться одночасно на різних курсах чи в різних групах.

Види занять сортуються так: лекція < семінар < лабораторне < практичне.

%\bigskip%\textit{Для деяких варіантів може знадобитись}




\par
%end
\newpage
%begin
{\center Варіант  \No \textbf { 130 }\par}
\bigskip
%type = tasks
%variant = 3
Обробляються результати розв’язування студентами задач на контрольних.
Одиничний факт роз\-в’я\-зу\-вання студентом задачі надалі розуміється як спроба.

\underline{Записи основного файлу містять поля}: відсоток набраних за розв’язання задачі балів, день, прізвище, місяць, рік, ім’я, умова задачі.

\underline{У допоміжному файлі наявні ключі}: загальна кількість записів, сумарний відсоток по всім записам файлу.

\underline{Знайти} всі задачі, які не розв’язав ніхто (з відсотком 0). Вивести по кожному з них інформацію:\par
{\noindent}--- на першому рядку:\par
умова, кількість спроб розв’язати;

{\noindent}--- на наступних рядках, починаючи з табуляції, вивести спроби по цих задачах (по одному на рядок):\par
прізвище, ім’я, рік, місяць, день

{\noindent}у такому сортуванні: рік (за спаданням), місяць, день, прізвище, ім’я.

%Задачі сортувати: кількість спроб (за спаданням), умова.

\bigskip\textit{Обмеження}

Студент міг розв’язувати задачу кілька разів.

Відсоток набраних за розв’язання задачі балів є цілим від 0 до 100.
День, місяць, рік є числовими полями, що зображують дату.
Раніше 2001 року статистика не велася. Рік є 4-значним.
Решта полів є непорожніми рядками.
Рядки починатися та закінчуватися білими символами не можуть.

Прізвище, ім’я та по-батькові (за наявності у варіанті) можуть мати до 22, 23, 24 відповідно символів включно кожне. 
Крім літер алфавіту можуть містити апостроф, дефіс та пробіл.

Умова задачі може мати від 3 до 67 символів включно.
Крім літер алфавіту може містити подвійні лапки, десяткові цифри, пробіл, дефіс,
а також символи . , \_ \{ \} \$ $+$ $*$ {\textbackslash} ( ) \% $[$ $]$ .

Студенти однозначно ідентифікуються прізвищем, ім’ям та по-батькові (за наявності у варіанті).
Задачі однозначно ідентифікуються умовою.

\bigskip\textit{Для деяких варіантів може знадобитись}

Розв'язуваність задачі --- середній відсоток набраних за задачу балів.

Розв'язуваність студента --- середній відсоток набраних студентом балів по всіх його спробах.

Розв'язуваність --- середній відсоток набраних студентами за задачі балів (рахується по всім даним).




\par
%end
\newpage
%begin
{\center Варіант  \No \textbf { 131 }\par}
\bigskip
%type = absense
%variant = 31
Обробляються дані семестрового журналу перевірок відвідування занять студентами деканатом (фіксуються відсутні). 

\underline{Записи основного файлу містять поля}: предмет, вид занять, ім’я, курс, код групи, день тижня, пара, прізвище, навчальний тиждень, аудиторія.

\underline{У допоміжному файлі наявні ключі}: кількість пропусків на четвертих парах, кількість пропусків лекцій.

\underline{Знайти} всіх студентів, що пропускали заняття виключно на четвертих парах. Вивести по кожному з них інформацію:\par
{\noindent}--- на першому рядку:\par
прізвище, ім’я, кількість пропущених пар;

{\noindent}--- на наступних рядках, починаючи з табуляції, вивести пропуски студента (по одному на рядок):\par
тиждень, день, пара, кількість пропусків

{\noindent}у такому сортуванні: тиждень, день, пара.

%Студентів сортувати: кількість пропусків, прізвище, ім’я.

\bigskip\textit{Обмеження}

Навчальний тиждень може приймати цілі значення від 1 до 14.
День тижня є цілим від 1 до 5. Пара є цілим від 1 до 4.
Курс є цілим від 1 до 4. Аудиторія є додатнім цілим.
Вид занять може приймати значення:
L --- лекція, п --- практичне, Сем. --- семінар, лаб. --- лабораторне.

Решта полів є непорожніми рядками.
Рядки починатися та закінчуватися білими символами не можуть.

Групи кодуються в межах курсу.
Код групи складається не більш ніж з 6 символів.
Може містити літери алфавіту, десяткові цифри та дефіс.

Предмет є рядком, може мати від 3 до 29 символів включно.
Може містити літери алфавіту, десяткові цифри, апостроф, дефіс, пробіл.
Предмети різних курсів вважати різними навіть за однакових назв.

Прізвище, ім’я та по-батькові (за наявності у варіанті)
можуть мати від 3 до 30, 20, 27 відповідно символів включно кожне.
Можуть містити літери алфавіту, апостроф, пробіл та дефіс.
Студенти однозначно ідентифікуються прізвищем, ім’ям та по-батькові (за наявності у варіанті).
Жоден студент не вчиться одночасно на різних курсах чи в різних групах.

Види занять сортуються так: практичне < лабораторне < лекція < семінар.

%\bigskip%\textit{Для деяких варіантів може знадобитись}




\par
%end
\newpage
%begin
{\center Варіант  \No \textbf { 132 }\par}
\bigskip
%type = absense
%variant = 37
Обробляються дані семестрового журналу перевірок відвідування занять студентами деканатом (фіксуються відсутні). 

\underline{Записи основного файлу містять поля}: курс, код групи, вид занять, навчальний тиждень, предмет, день тижня, пара, аудиторія, прізвище, ім’я.

\underline{У допоміжному файлі наявні ключі}: кількість зафіксованих пропусків, кількість пропусків лекцій.

\underline{Знайти} всіх студентів, що по вівторках пропускали заняття більше ніж у середньому. Вивести по кожному з них інформацію:\par
{\noindent}--- на першому рядку:\par
прізвище, ім’я, середня кількість пропусків (по днях тижня, округлена до 1 знаку), кількість пропусків по вівторках;

{\noindent}--- на наступних рядках, починаючи з табуляції, вивести пропуски студента (по одному на рядок):\par
предмет, вид занять, тиждень, день, пара

{\noindent}у такому сортуванні: день, тиждень, пара, предмет, вид занять.

%Студентів сортувати: середня кількість пропусків (округлена), кількість пропусків по вівторках, прізвище, ім’я.

\bigskip\textit{Обмеження}

Навчальний тиждень може приймати цілі значення від 1 до 17.
День тижня є цілим від 1 до 5. Пара є цілим від 1 до 4.
Курс є цілим від 1 до 4. Аудиторія є додатнім цілим.
Вид занять може приймати значення:
лекц. --- лекція, Практ. --- практичне, S --- семінар, lab --- лабораторне.

Решта полів є непорожніми рядками.
Рядки починатися та закінчуватися білими символами не можуть.

Групи кодуються в межах курсу.
Код групи складається не більш ніж з 5 символів.
Може містити літери алфавіту, десяткові цифри та дефіс.

Предмет є рядком, може мати від 4 до 35 символів включно.
Може містити літери алфавіту, десяткові цифри, апостроф, дефіс, пробіл.
Предмети різних курсів вважати різними навіть за однакових назв.

Прізвище, ім’я та по-батькові (за наявності у варіанті)
можуть мати від 4 до 27, 22, 23 відповідно символів включно кожне.
Можуть містити літери алфавіту, апостроф, пробіл та дефіс.
Студенти однозначно ідентифікуються прізвищем, ім’ям та по-батькові (за наявності у варіанті).
Жоден студент не вчиться одночасно на різних курсах чи в різних групах.

Види занять сортуються так: семінар < лабораторне < практичне < лекція.

%\bigskip%\textit{Для деяких варіантів може знадобитись}




\par
%end
\newpage
%begin
{\center Варіант  \No \textbf { 133 }\par}
\bigskip
%type = absense
%variant = 47
Обробляються дані семестрового журналу перевірок відвідування занять студентами деканатом (фіксуються відсутні). 

\underline{Записи основного файлу містять поля}: по-батькові, аудиторія, навчальний тиждень, день тижня, пара, предмет, ім’я, вид занять, курс, код групи, прізвище.

\underline{У допоміжному файлі наявні ключі}: середня кількість пропусків предметів, округлена до 2 знаків, SMILE.

\underline{Знайти} всі предмети, більша половина пропусків з яких – пропуски лекцій. Вивести по кожному з них інформацію:\par
{\noindent}--- на першому рядку:\par
курс, предмет, кількість пропусків лекцій, кількість пропусків;

{\noindent}--- на наступних рядках, починаючи з табуляції, вивести пропуски предмета (по одному на рядок):\par
прізвище, ім’я, по-батькові, тиждень, вид занять, кількість пропусків

{\noindent}у такому сортуванні: вид занять, прізвище, ім’я, по-батькові, тиждень.

%Предмети сортувати: кількість пропусків лекцій, курс, предмет.

\bigskip\textit{Обмеження}

Навчальний тиждень може приймати цілі значення від 1 до 19.
День тижня є цілим від 1 до 5. Пара є цілим від 1 до 4.
Курс є цілим від 1 до 4. Аудиторія є додатнім цілим.
Вид занять може приймати значення:
Lect --- лекція, практ. --- практичне, s --- семінар, Lab --- лабораторне.

Решта полів є непорожніми рядками.
Рядки починатися та закінчуватися білими символами не можуть.

Групи кодуються в межах курсу.
Код групи складається не більш ніж з 4 символів.
Може містити літери алфавіту, десяткові цифри та дефіс.

Предмет є рядком, може мати від 4 до 34 символів включно.
Може містити літери алфавіту, десяткові цифри, апостроф, дефіс, пробіл.
Предмети різних курсів вважати різними навіть за однакових назв.

Прізвище, ім’я та по-батькові (за наявності у варіанті)
можуть мати від 4 до 26, 25, 25 відповідно символів включно кожне.
Можуть містити літери алфавіту, апостроф, пробіл та дефіс.
Студенти однозначно ідентифікуються прізвищем, ім’ям та по-батькові (за наявності у варіанті).
Жоден студент не вчиться одночасно на різних курсах чи в різних групах.

Види занять сортуються так: лабораторне < семінар < лекція < практичне.

%\bigskip%\textit{Для деяких варіантів може знадобитись}




\par
%end
\newpage
%begin
{\center Варіант  \No \textbf { 134 }\par}
\bigskip
%type = meteo
%variant = 113
Обробляються дані спостережень за погодою. 

\underline{Записи основного файлу містять поля}: рік, середня денна температура, код метеостанції, місяць, день, відносна вологість, кількість опадів, максимальна денна температура, мінімальна денна температура, сила вітру.

\underline{У допоміжному файлі наявні ключі}: останній рік, зазначений у спостереженнях; кількість метеостанцій, по яких велися спостереження в останньому році.

\underline{Знайти} метеостанції, по яких фіксувалася сила вітру у два рази більша за середню по метеостанції. Вивести по кожному з них інформацію:\par
{\noindent}--- на першому рядку:\par
метеостанція, найбільша сила вітру, середня різниця максимальної та мінімальної денних температур;

{\noindent}--- на наступних рядках, починаючи з табуляції, вивести для них дані по датах (по одному на рядок):\par
середня денна температура, сила вітру, вологість, місяць, день, рік

{\noindent}у такому сортуванні: місяць, день, рік.

%Метеостанції сортувати: найбільша сила вітру (за спаданням), метеостанція.

\bigskip\textit{Обмеження}

Код метеостанції є рядком довжини від 3 до 12 включно.
Може містити літери алфавіту, десяткові цифри, підкреслення. Решта полів числові.

Рік є чотиризначним цілим.
Кількість опадів є цілим значенням.
Сила вітру є дійсним значенням, подається у форматі з 2 знаками після крапки.
Температури є дійсними значеннями, подаються у форматі з 2 знаками після крапки.
Цілі значення подаються без десяткової крапки.
Для дійсних використовується зображення з фіксованою крапкою.

\textit{Вихідне форматування дійсних}

Усі дійсні значення виводяться у форматі з фіксованою крапкою з тією ж кількістю знаків після крапки, що і у вхідному зображенні.

Якщо у варіанті раніше не було зазначено, то \underline{обчислювані програмою} середні значення виводяться у форматі з 2 знаками після крапки.
%\bigskip%\textit{Для деяких варіантів може знадобитись}




\par
%end
\newpage
%begin
{\center Варіант  \No \textbf { 135 }\par}
\bigskip
%type = tasks
%variant = 15
Обробляються результати розв’язування студентами задач на контрольних.
Одиничний факт роз\-в’я\-зу\-вання студентом задачі надалі розуміється як спроба.

\underline{Записи основного файлу містять поля}: ім’я, по-батькові, рік, відсоток набраних за розв’язання задачі балів, прізвище, умова задачі.

\underline{У допоміжному файлі наявні ключі}: найбільша довжина умови, загальна кількість записів.

\underline{Знайти} всіх студентів, що не мають спроб з відсотком не менше 90\%. Вивести по кожному з них інформацію:\par
{\noindent}--- на першому рядку:\par
ім’я, по-батькові, прізвище, середня розв’язуваність студента (округлена до цілого), кількість спроб;

{\noindent}--- на наступних рядках, починаючи з табуляції, вивести їх спроби (по одному на рядок):\par
умова, рік, відсоток набраних за задачу балів

{\noindent}у такому сортуванні: умова, відсоток, рік.

%Студентів сортувати: середня розв’язуваність, кількість спроб (за спаданням), прізвище, ім’я, по-батькові.

\bigskip\textit{Обмеження}

Студент міг розв’язувати задачу кілька разів.

Відсоток набраних за розв’язання задачі балів є цілим від 0 до 100.
День, місяць, рік є числовими полями, що зображують дату.
Раніше 2001 року статистика не велася. Рік є 4-значним.
Решта полів є непорожніми рядками.
Рядки починатися та закінчуватися білими символами не можуть.

Прізвище, ім’я та по-батькові (за наявності у варіанті) можуть мати до 20, 29, 20 відповідно символів включно кожне. 
Крім літер алфавіту можуть містити апостроф, дефіс та пробіл.

Умова задачі може мати від 5 до 58 символів включно.
Крім літер алфавіту може містити подвійні лапки, десяткові цифри, пробіл, дефіс,
а також символи . , \_ \{ \} \$ $+$ $*$ {\textbackslash} ( ) \% $[$ $]$ .

Студенти однозначно ідентифікуються прізвищем, ім’ям та по-батькові (за наявності у варіанті).
Задачі однозначно ідентифікуються умовою.

\bigskip\textit{Для деяких варіантів може знадобитись}

Розв'язуваність задачі --- середній відсоток набраних за задачу балів.

Розв'язуваність студента --- середній відсоток набраних студентом балів по всіх його спробах.

Розв'язуваність --- середній відсоток набраних студентами за задачі балів (рахується по всім даним).




\par
%end
\newpage
%begin
{\center Варіант  \No \textbf { 136 }\par}
\bigskip
%type = meteo
%variant = 98
Обробляються дані спостережень за погодою. 

\underline{Записи основного файлу містять поля}: код метеостанції, сила вітру, рік, максимальна денна температура, відносна вологість, кількість опадів, середня денна температура, мінімальна денна температура, день, місяць.

\underline{У допоміжному файлі наявні ключі}: загальна кількість записів, перший рік спостережень.

\underline{Знайти} дати, по яких фіксувалась відносна вологість більша ніж в середньому по всіх спостереженнях. Вивести по кожному з них інформацію:\par
{\noindent}--- на першому рядку:\par
кількість опадів, середня середньоденна температура, місяць, день, рік;

{\noindent}--- на наступних рядках, починаючи з табуляції, вивести для дати спостереження метеостанції, коди яких закінчуються на 17 (по одному на рядок):\par
метеостанція, вологість, середня денна температура, максимальна денна температура, мінімальна денна температура, сила вітру

{\noindent}у такому сортуванні: метеостанція.

%Дати сортувати: середня середньоденна температура, рік (за спаданням), місяць, день.

\bigskip\textit{Обмеження}

Код метеостанції є рядком довжини від 3 до 8 включно.
Може містити літери алфавіту, десяткові цифри, підкреслення. Решта полів числові.

Рік є чотиризначним цілим.
Кількість опадів є дійсним значенням, подається у форматі з 2 знаками після крапки.
Сила вітру є дійсним значенням, подається у форматі з 1 знаком після крапки.
Температури є дійсними значеннями, подаються у форматі з 3 знаками після крапки.
Цілі значення подаються без десяткової крапки.
Для дійсних використовується зображення з фіксованою крапкою.

\textit{Вихідне форматування дійсних}

Усі дійсні значення виводяться у форматі з фіксованою крапкою з тією ж кількістю знаків після крапки, що і у вхідному зображенні.

Якщо у варіанті раніше не було зазначено, то \underline{обчислювані програмою} середні значення виводяться у форматі з 3 знаками після крапки.
%\bigskip%\textit{Для деяких варіантів може знадобитись}




\par
%end
\newpage

%begin
{\center Варіант  \No \textbf { 137 }\par}
\bigskip
%type = tasks
%variant = 17
Обробляються результати розв’язування студентами задач на контрольних.
Одиничний факт роз\-в’я\-зу\-вання студентом задачі надалі розуміється як спроба.

\underline{Записи основного файлу містять поля}: день, ім’я, відсоток набраних за розв’язання задачі балів, рік, умова задачі, прізвище, місяць.

\underline{У допоміжному файлі наявні ключі}: загальна кількість студентів; найменше значення року, що зустрічається у файлі.

\underline{Знайти} всіх студентів, що не робили спроб ані з відсотком 0\%, ані з відсотком 100\%. Вивести по кожному з них інформацію:\par
{\noindent}--- на першому рядку:\par
прізвище, ім’я, по-батькові, кількість спроб, кількість спроб з відсотком 60\%;

{\noindent}--- на наступних рядках, починаючи з табуляції, вивести їх спроби (по одному на рядок):\par
відсоток набраних за задачу балів, рік, місяць, день умова

{\noindent}у такому сортуванні: умова, відсоток (за спаданням), рік, місяць, день.

%Студентів сортувати: кількість спроб (за спаданням), прізвище, ім’я.

\bigskip\textit{Обмеження}

Студент міг розв’язувати задачу кілька разів.

Відсоток набраних за розв’язання задачі балів є цілим від 0 до 100.
День, місяць, рік є числовими полями, що зображують дату.
Раніше 2001 року статистика не велася. Рік є 4-значним.
Решта полів є непорожніми рядками.
Рядки починатися та закінчуватися білими символами не можуть.

Прізвище, ім’я та по-батькові (за наявності у варіанті) можуть мати до 29, 20, 23 відповідно символів включно кожне. 
Крім літер алфавіту можуть містити апостроф, дефіс та пробіл.

Умова задачі може мати від 3 до 65 символів включно.
Крім літер алфавіту може містити подвійні лапки, десяткові цифри, пробіл, дефіс,
а також символи . , \_ \{ \} \$ $+$ $*$ {\textbackslash} ( ) \% $[$ $]$ .

Студенти однозначно ідентифікуються прізвищем, ім’ям та по-батькові (за наявності у варіанті).
Задачі однозначно ідентифікуються умовою.

\bigskip\textit{Для деяких варіантів може знадобитись}

Розв'язуваність задачі --- середній відсоток набраних за задачу балів.

Розв'язуваність студента --- середній відсоток набраних студентом балів по всіх його спробах.

Розв'язуваність --- середній відсоток набраних студентами за задачі балів (рахується по всім даним).




\par
%end
\newpage

%begin
{\center Варіант  \No \textbf { 138 }\par}
\bigskip
%type = absense
%variant = 22
Обробляються дані семестрового журналу перевірок відвідування занять студентами деканатом (фіксуються відсутні). 

\underline{Записи основного файлу містять поля}: аудиторія, вид занять, прізвище, ім’я, предмет, навчальний тиждень, день тижня, пара, курс, код групи.

\underline{У допоміжному файлі наявні ключі}: кількість лекцій, кількість записів.

\underline{Знайти} всіх студентів, що пропустили пар не більше ніж в середньому. Вивести по кожному з них інформацію:\par
{\noindent}--- на першому рядку:\par
прізвище, ім’я, кількість пропущених пар, середню кількість пропусків по всіх студентах (округлена до 1 знаку);

{\noindent}--- на наступних рядках, починаючи з табуляції, вивести пропуски студента (по одному на рядок):\par
тиждень, день, пара, аудиторія, предмет, вид занять

{\noindent}у такому сортуванні: вид занять, предмет, тиждень, день, пара.

%Студентів сортувати: кількість пропусків, прізвище, ім’я.

\bigskip\textit{Обмеження}

Навчальний тиждень може приймати цілі значення від 1 до 19.
День тижня є цілим від 1 до 5. Пара є цілим від 1 до 4.
Курс є цілим від 1 до 4. Аудиторія є додатнім цілим.
Вид занять може приймати значення:
лекц. --- лекція, п --- практичне, Seminar --- семінар, лаб. --- лабораторне.

Решта полів є непорожніми рядками.
Рядки починатися та закінчуватися білими символами не можуть.

Групи кодуються в межах курсу.
Код групи складається не більш ніж з 6 символів.
Може містити літери алфавіту, десяткові цифри та дефіс.

Предмет є рядком, може мати від 4 до 28 символів включно.
Може містити літери алфавіту, десяткові цифри, апостроф, дефіс, пробіл.
Предмети різних курсів вважати різними навіть за однакових назв.

Прізвище, ім’я та по-батькові (за наявності у варіанті)
можуть мати від 4 до 29, 28, 29 відповідно символів включно кожне.
Можуть містити літери алфавіту, апостроф, пробіл та дефіс.
Студенти однозначно ідентифікуються прізвищем, ім’ям та по-батькові (за наявності у варіанті).
Жоден студент не вчиться одночасно на різних курсах чи в різних групах.

Види занять сортуються так: лекція < семінар < практичне < лабораторне.

%\bigskip%\textit{Для деяких варіантів може знадобитись}




\par
%end
\newpage

%begin
{\center Варіант  \No \textbf { 139 }\par}
\bigskip
%type = session
%variant = 61
Обробляються результати складання студентами екзаменів зимової сесії.

\underline{Записи основного файлу містять поля}: по-батькові, сумарна оцінка в балах за 100-бальною системою, набрані впродовж семестру бали, предмет, прізвище, ім’я, код групи, номер залікової книжки, набрані на екзамені бали, оцінка за державною шкалою.

\underline{У допоміжному файлі наявні ключі}: кількість «незадовільно», кількість сумарних оцінок від 60 до 70 включно.

\underline{Знайти} всіх студентів, що по кожному предмету отримали оцінку 3 за державною шкалою. Вивести по кожному з них інформацію:\par
{\noindent}--- на першому рядку:\par
середня оцінка в балах (округлена до цілого), прізвище, ім’я, номер залікової книжки, кількість екзаменів;

{\noindent}--- на наступних рядках, починаючи з табуляції, вивести для студента всі його результати (по одному на рядок):\par
предмет, сумарна оцінка, бали на екзамені, бали в семестрі

{\noindent}у такому сортуванні: бали на екзамені, предмет.

%Студентів сортувати: середня оцінка в балах (після округлення), прізвище, ім’я, номер залікової книжки.

\bigskip\textit{Обмеження}

Студент не може складати предмет більше одного разу.

Бали та оцінка є значеннями цілих числових типів. Решта полів є непорожніми рядками.
Рядки розпочинатись та закінчуватись білими символами не можуть.

Оцінка за державною шкалою може приймати значення: 2-5~--- як зазвичай, 0~--- не з’явився,$-1$~--- не допущений. 
Оцінки 3-5 вважаються задовільними, решта~--- незадовільними. 
Бали є цілими та невід’ємними.
Набрані впродовж семестру бали не можуть перевищувати 60. 
Набрані на екзамені бали не можуть перевищувати 40. 
Набрані на екзамені бали повинні бути не менше 24 або дорівнювати 0(в останньому випадку оцінка не може бути задовільною). 
Сумарна оцінка в балах є сумою балів, набраних у семестрі та на екзамені.
Сумарна оцінка в балах та за державною шкалою мають бути узгоджені між собою.

Прізвище, ім’я та по-батькові (якщо наявне у варіанті) можуть мати до 30, 28, 26 відповідно символів включно кожне. 
Крім літер алфавіту можуть містити апостроф, дефіс та пробіл.

Предмет може мати від 4 до 51 символів включно. 
Крім літер алфавіту може містити подвійні лапки, пробіл та дефіс.

Код групи є непорожнім та складається не більш ніж з 7 літер. 
Крім літер алфавіту може містити десяткові цифри та дефіс.


Номер залікової книжки однозначно ідентифікує студента та складається з 7 десяткових цифр. 

\bigskip\textit{Для деяких варіантів може знадобитись}

Середній бал/оцінка за предметом визначається як середнє арифметичнесумарних балів/державних оцінок за цим предметом. 
Середній бал/оцінка студента визначається як середнє арифметичне сумарних балів/державних оцінокцього студента. 
За обчислення середньої оцінки недопущені та ті, що не з’явилися, рахуються як 2.
У решті випадків недопуски підсумовуються як $-1$, а «не з’явився» як 0.
Якщо студент не гірше ніж задовільно склав усі предмети, 
то його рейтинг за балами/оцінкою визначається як середній бал/оцінка; інакше дорівнює 0.

Стабільність студента визначається як модуль різниці між його кращим та найгіршим сумарним балом.
Стабільність предмета визначається як модуль різниці між найкращим та найгіршимсумарним балом за цим предметом.






\par
%end
\newpage

%begin
{\center Варіант  \No \textbf { 140 }\par}
\bigskip
%type = tasks
%variant = 7
Обробляються результати розв’язування студентами задач на контрольних.
Одиничний факт роз\-в’я\-зу\-вання студентом задачі надалі розуміється як спроба.

\underline{Записи основного файлу містять поля}: ім’я, прізвище, відсоток набраних за розв’язання задачі балів, по-батькові, умова задачі, рік, місяць, день.

\underline{У допоміжному файлі наявні ключі}: сума довжин всіх прізвищ, кількість записів.

\underline{Знайти} всі задачі, які в кожній спробі були розв’язані не менш ніж на 60\%. Вивести по кожному з них інформацію:\par
{\noindent}--- на першому рядку:\par
умова, кількість спроб, кількість спроб не гірше 60\%, відсоток спроб не гірше 60\% до кількості спроб по задачі (з округленням до цілого);

{\noindent}--- на наступних рядках, починаючи з табуляції, вивести спроби по цих задачах (по одному на рядок):\par
рік, прізвище, ім’я, по-батькові, відсоток набраних балів

{\noindent}у такому сортуванні: прізвище, ім’я, по-батькові, відсоток набраних балів (за спаданням), рік.

%Задачі сортувати: відсоток спроб не гірше 60\% (з виконаним округленням, за спаданням), умова.

\bigskip\textit{Обмеження}

Студент міг розв’язувати задачу кілька разів.

Відсоток набраних за розв’язання задачі балів є цілим від 0 до 100.
День, місяць, рік є числовими полями, що зображують дату.
Раніше 2001 року статистика не велася. Рік є 4-значним.
Решта полів є непорожніми рядками.
Рядки починатися та закінчуватися білими символами не можуть.

Прізвище, ім’я та по-батькові (за наявності у варіанті) можуть мати до 25, 26, 30 відповідно символів включно кожне. 
Крім літер алфавіту можуть містити апостроф, дефіс та пробіл.

Умова задачі може мати від 4 до 54 символів включно.
Крім літер алфавіту може містити подвійні лапки, десяткові цифри, пробіл, дефіс,
а також символи . , \_ \{ \} \$ $+$ $*$ {\textbackslash} ( ) \% $[$ $]$ .

Студенти однозначно ідентифікуються прізвищем, ім’ям та по-батькові (за наявності у варіанті).
Задачі однозначно ідентифікуються умовою.

\bigskip\textit{Для деяких варіантів може знадобитись}

Розв'язуваність задачі --- середній відсоток набраних за задачу балів.

Розв'язуваність студента --- середній відсоток набраних студентом балів по всіх його спробах.

Розв'язуваність --- середній відсоток набраних студентами за задачі балів (рахується по всім даним).




\par
%end
\newpage

\end{document}
